% 本文件由 LaTeX 排版 Agent 根据有效 translation.json 自动生成。
% author=Augustine of Hippo; work_id=1501; title=论功绩与罪的赦免，及婴孩的圣洗

\chapter{论功绩与罪的赦免，及婴孩的圣洗}

% structure: p0001 source=h1 target=omitted reason=page-title-already-rendered-by-parent

\Emph{第一卷}\newline{}\Emph{第二卷}\newline{}\Emph{第三卷}\par

\SourceRecord{Translated by Peter Holmes and Robert Ernest Wallis, and revised by Benjamin B. Warfield.；Nicene and Post-Nicene Fathers, First Series；Vol. 5.；Edited by Philip Schaff.；Buffalo, NY: Christian Literature Publishing Co.,；1887.；<http://www.newadvent.org/fathers/1501.htm>.}

% structure: p0001 source=h1 target=section reason=page-title

\section{论功绩与赦罪及婴儿圣洗（第一卷）}

\Emph{在此卷中，他驳斥那些主张即使\Person{亚当}从未犯罪也必定会死的人，以及那些认为\Person{亚当}的罪没有通过自然遗传传给后代的人。他还表明，死亡并非由于人性的必然性而临于人，而是作为罪的惩罚；他继而证明\Person{亚当}的罪牵连了他的所有后代，指出婴儿受洗正是为了获得原罪的赦免。}\par

% structure: p0003 source=h2 target=subsection reason=relative-heading-rank; base=h2

\subsection{第一章 [I.]——以题献形式写给友人\Person{玛策利诺}的引论}

尽管我们在与那些离弃天主法律的罪人周旋的纷扰和热情中，深陷焦虑与烦恼——即使我们完全可以将这些恶事归咎于自己的罪过——但我却不愿意，而且，说实话，也无法再继续欠债，我至亲的\Person{玛策利诺}，欠你那份热诚的爱，这份爱只加深了我对你感恩而愉悦的评价。我受［双重情感］的推动：一方面，正是那爱使我们不可分割地合而为一，共同盼望那改善的转变；另一方面，是怕在你身上得罪天主，因祂赐予你如此恳切的渴望；满足这渴望，我必是在侍奉那赐予你的主。这推动如此强烈地引导并吸引我，尽我微薄之力解决你以书面提出的问题，以至于我的心智逐渐视此探究凌驾于一切之上；［它如今不让我安息］，直到我完成某事，足以表明我已顺从——即便不是充分的，也是服从的——你的善意和那些为此问题焦虑之人的愿望。\par

% structure: p0005 source=h2 target=subsection reason=relative-heading-rank; base=h2

\subsection{第二章 [II.]——如果\Person{亚当}没有犯罪，他将永远不会死}

那些主张\Person{亚当}受造如此，即使没有任何罪的过失也必定会死，不是作为罪的惩罚，而是由于他存在的必然性的人，确实试图将法律中那段话：\BibleQuote{你们吃那日子，必定要死。}\BibleRef{创 2:17}不是指肉身的死，而是指那发生在罪中的灵魂的死。是那些不信者经历了这种死，主指着他们说：\Quote{让死人埋葬他们的死人。}当我们读到天主在\Person{亚当}犯罪后责备并判决他时，对他说：\BibleQuote{你既是灰土，你还要归于灰土。}\BibleRef{创 3:19}他们又将如何回答呢？因为他被称为\Quote{灰土}不是就他的灵魂而言，而显然是由于他的身体，并且正是通过这同一个身体的死，他注定要\Quote{归于灰土}。然而，虽然就他的身体而言他是灰土，虽然他带着受造时所拥有的自然身体，但如果他没有犯罪，他本应变成属灵的身体，进入那不腐朽的状态，这是为信徒和圣人所应许的，而没有死亡的危险。\BibleRef{格前 15:52}对于这事，我们不仅在自己内意识到有热切的渴望，也从宗徒的暗示中得知，他说：\BibleQuote{为此我们叹息，渴望穿上我们那来自天上的住处；只要穿上了，就不会被看见是赤裸的。我们在这帐幕里叹息，深觉负担，并非想脱去，而是想穿上，好使这必死的被生命吞灭。}\BibleRef{格后 5:2}因此，如果\Person{亚当}没有犯罪，他不会被剥夺身体，而是穿上不朽和不腐，使\Quote{这必死的被生命吞灭}；也就是说，他可能从自然身体变成属灵身体。\par

% structure: p0007 source=h2 target=subsection reason=relative-heading-rank; base=h2

\subsection{第三章 [III.]——会死是一回事，注定要死是另一回事}

那完全不必担忧：如果他碰巧在自然身体中活得更久，他会不会被衰老压倒，并逐渐因年岁增长而走向死亡。因为如果天主赐予以色列人的衣服和鞋子，使它们在这么多年里\BibleQuote{没有变旧}\BibleRef{申 29:5}，那么，如果因着服从，同一天主的能力允许人，即使他有自然可死的身体，仍能在其中拥有某种状态，使他可以寿终正寝而不衰朽，并且在天主乐意的时候，不经死亡就从可死变为不死，这又有什么可惊奇的呢？因为正如我们现在拥有的这血肉之躯，并非因为不必受伤就是不可伤害的；同样，他的身体也并非因为不必死亡就是不朽的。我相信，这种状态，即仍在自然可死的身体中，甚至赐给了那些不经死亡而从此处被提升的人。因为\Person{哈诺客}和\Person{厄里亚}并没有因长寿而衰朽。但我并不认为他们当时已变成那种属灵的身体，如同复活时所应许的，并且主是第一个领受的；只是他们可能不需要那些通过使用而给身体提供滋养的食物；但从他们被提升以来，他们这样活着，享受着如同\Person{厄里亚}在四十天内靠水壶和饼维生、无实质性食物时所得到的那种充足\BibleRef{列上 19:8}；或者，如果有需要这样的食物，他们可能像\Person{亚当}在乐园中被供养，在他因罪而招致被驱逐之前。而我猜想，\Person{亚当}是从各种树木的果实中得到对抗腐朽的滋养，并从生命树得到对抗衰老的保障。\par

% structure: p0009 source=h2 target=subsection reason=relative-heading-rank; base=h2

\subsection{第四章 [IV.]——甚至肉身的死也来自罪}

但除了天主在惩罚中所说的那段话：\BibleQuote{你原是灰土，将来还要归于灰土}\BibleRef{创 3:19}——我无法理解任何人怎能将这段话应用在肉体死亡之外——还有其他同样明确的证据，足以充分证明人类因罪而招致的不仅是灵魂的死亡，也包括肉体的死亡。宗徒在致罗马人书中说：\BibleQuote{如果基督在你们内，身体固然因罪而死亡，但神魂却因正义而生活。再者，如果那使耶稣从死者中复活者的圣神住在你们内，那么，那使基督从死者中复活的，也必要藉那住在你们内的圣神，使你们有死的身体复活。}\BibleRef{罗 8:10}我认为这样明白清楚的句子只需阅读，无需解释。\Emph{身体}，他说，\Emph{是死的}，并非因为尘世的脆弱，好像是由地上的尘土所造，而是\Emph{由于罪}；我们还需要什么？他措辞非常谨慎：他没有说\Quote{是\Emph{能死的}}，而是说\Quote{\Emph{死的}}。\par

% structure: p0011 source=h2 target=subsection reason=relative-heading-rank; base=h2

\subsection{第5章 [V.]——这些词：\OriginalTerm{能死的（Capable of Dying）}{Mortale}、\OriginalTerm{死的（Dead）}{Mortuum}和\OriginalTerm{注定要死的（Destined to Die）}{Moriturus}}

在进入圣徒复活中所应许的不朽状态之前，身体可以是\OriginalTerm{能死的（capable of dying）}{mortale}，尽管不一定\OriginalTerm{注定要死的（destined to die）}{moriturus}；正如我们现在的身体，可以说可以生病，尽管不一定生病。因为谁的肉体是不生病的，即使偶然在生病之前就死亡呢？同样，当时人的身体是能死的；如果人持守了正义，即服从，这种能死性原本会被永恒的不朽所取代；但即使是\OriginalTerm{能死的（capable of dying）}{mortale}也没有成为\OriginalTerm{死的（dead）}{mortuum}，除非是因为罪。因为将来在复活时发生的改变，事实上，不仅不会遭遇死亡（这因罪而发生），而且也不会有能死性（或死亡的可能性），即犯罪前自然身体所具有的。他没有说：\Quote{那使基督耶稣从死者中复活的，也必使你们\Emph{死的}身体复生}（尽管他先前说过：\Quote{身体是死的}\BibleRef{罗 8:10}）；而是说：\Quote{他也必使你们\Emph{能死的}身体复生}\BibleRef{罗 8:11}；这样，它们不仅不再死亡，而且不再能死，因为属自然的被复活为属神的，这能死的身体将穿上不朽，死亡将被生命吞灭。\par

% structure: p0013 source=h2 target=subsection reason=relative-heading-rank; base=h2

\subsection{第6章 [VI.]——身体如何因罪而死亡}

有人惊讶，还有什么比我们已提供的证据更清楚呢。但或许我们必须满足于听到这清晰的说明被争论所反驳，即我们应理解此处的\Quote{死身体}\BibleRef{罗 8:10} 正如那节经文所说：\Quote{你们要致死属于地上的肢体。}\BibleRef{哥 3:5} 但身体之所以如此被致死，是因 \Emph{义德}而非因罪恶；因为我们致死那属于地上的肢体，是为了行义德的工作。或者，如果他们以为\Quote{因罪恶}一语是附加的，并非要我们理解为\Quote{因为罪恶已犯}，而是\Quote{为使罪恶不犯}——就好像是说：\Quote{身体固然是因罪恶而死，为使罪恶不犯}——那么，他在下一句中又加上\Quote{因义德}一词，说\Quote{圣神却是生命}\BibleRef{罗 8:10} 是什么意思呢？因为只需简单地加上\Quote{圣神是生命}，便可同样理解为\Quote{为使罪恶不犯}；这样，两个命题便指向同一件事——即\Quote{身体是死的}和\Quote{圣神是生命}都是为了共同的目的\Quote{为使罪恶不犯}。同样，如果他仅仅想说\Quote{因义德}，即\Quote{为行义德}的意思，那么两个分句也可能指向此一目的——即\Quote{身体是死的}和\Quote{圣神是生命}都是\Quote{为行义德}。但按经文实际所说，\Quote{身体因罪恶而死，圣神因义德而活}，将不同的功过归于不同的事物——将罪的恶报归于身体的死亡，将义德的功绩归于圣神的生命。因此，既然无人能怀疑\Quote{圣神因义德而活}，即因义德的功绩；那么，我们如何或能够理解\Quote{身体因罪恶而死}这句话，除了身体是因罪恶的功绩而死，除非我们试图按自己武断的意愿曲解或强解圣经最明显的含义？此外，后面的经文也提供了进一步的启发。因为他是限于现在的时候说：一方面，\Quote{身体因罪恶而死}，因为身体在未因复活而更新时，仍保留着罪恶的功绩，即死的必然性；另一方面，\Quote{圣神因义德而活}，因为尽管我们仍背负着\Quote{这死身体}\BibleRef{罗 7:24}，但藉着在我们内已开始的更新，我们对因信德而来的义德有了新的渴望。然而，恐怕人在无知中失去对身体复活的希望，他说：那在此世他刚说成\Quote{因罪恶而死}的身体，在来世将\Quote{因义德而活}——不仅是死中复生，而且是从可朽变为不朽。\par

% structure: p0015 source=h2 target=subsection reason=relative-heading-rank; base=h2

\subsection{第七章 [VII]——身体的生命是希望的对象，圣神的生命是其前奏。}

虽然我非常担心如此清楚的事情可能会因解释而变得模糊，但我仍恳请你们注意宗徒那光明的陈述。他说：\Quote{如果基督在你们内，身体固然因罪恶而死，圣神却因义德而活。}\BibleRef{罗 8:10} 这话是为使人不以为他们从基督的恩宠中毫无或只得到微少的益处，因为他们的身体免不了死。因为他们必须记得，虽然他们的身体仍带着那必然归于死亡状况的罪恶的功绩，但他们的圣神已因信德的义德开始活了，尽管它曾因不信的死亡而实际熄灭。因此，他说，你们必须认为，因基督在你们内而赐予你们的恩赐不小；即在因罪恶而死的身体中，你们的圣神现今已因义德而活；所以，你们不应连自己身体的性命也绝望。\Quote{如果那使基督从死者中复活的圣神住在你们内，那么那使基督从死者中复活的，也必藉那住在你们内的圣神，使你们有死的身体复活。}\BibleRef{罗 8:11} 争论的烟雾怎能遮蔽如此清晰的光明呢？宗徒明确告诉你们：虽然身体在你们内因罪恶而死，但你们有死的身体却要因义德而复活——正是那义德使你们的圣神现今已活着——整个过程都要藉基督的恩宠，即住在你们内的圣神来完成：而人仍然在反驳！他接着告诉我们，生命如何藉着致死死亡而将死亡转化为自身。他说：\Quote{所以，弟兄们！我们并不欠肉债，去顺从肉性生活；如果你们顺从肉性生活，必要死亡；如果你们藉着圣神致死肉体的行为，必要生活。}\BibleRef{罗 8:12} 这不就是说：如果你们顺从死亡而生活，你们将完全死亡；但如果你们藉着顺从生命而致死死亡，你们将完全生活？\par

% structure: p0017 source=h2 target=subsection reason=relative-heading-rank; base=h2

\subsection{第八章 [VIII]——身体之死来自亚当的罪。}

当他说类似的话时：\Quote{藉着人来了死亡，也藉着人来了死者的复活，}\BibleRef{格前15:21}除了指身体的死亡，这段经文还能作何理解？因为他提及此事，进而讲到身体的复活，并以最恳切隆重的论述肯定了它。他在致格林多人书中说：\Quote{藉着人来了死亡，也藉着人来了死者的复活；因为就如众人都在亚当内死了，照样众人也都要在基督内复活，}\BibleRef{格前15:21}——这岂不是与他对罗马人所说的那句话表达了同样的含义：\Quote{藉着一人，罪恶进入了世界，罪恶又带来了死亡？}\BibleRef{罗5:12}现在他们却坚持认为，这里的死亡不是身体的死亡，而是灵魂的死亡，借口说对格林多人讲的是另一回事，那里他们完全无法理解为灵魂的死亡，因为所讨论的是身体的复活，这是与身体的死亡相对立的。此外，这里为何只说死亡是由人带来的，而不提罪，是因为论述的重点不是关于\OriginalTerm{义德}{righteousness}（义德是与罪相对的），而是关于身体的复活（复活是与身体的死亡相对的）。\par

% structure: p0019 source=h2 target=subsection reason=relative-heading-rank; base=h2

\subsection{第九章——罪藉自然遗传而非仅藉模仿传到众人}

你在来信中告诉我，他们试图将宗徒的话——\Quote{藉着一人，罪恶进入了世界，死亡藉着罪恶也进入了世界；}\BibleRef{罗5:12}——曲解成新的含义，但你并未告诉我他们以为这些字句是什么意思。但我从其他人那里了解到，他们认为这里所说的死亡不是身体的死亡（他们不允许亚当因自己的罪而应受身体的死亡），而是灵魂的死亡，即在实际犯罪中发生的死亡；并且这实际罪并非由第一人藉自然遗传传给他人，而是藉模仿。因此，他们也拒绝相信在婴儿身上，原罪藉圣洗得以赦免，因为他们争辩说人出生时根本不存在这样的原罪。但如果宗徒想要断言罪进入世界不是藉自然遗传，而是藉模仿，他就会指出第一个犯罪者不是亚当，而是魔鬼，关于他写道：\Quote{他从起初就犯罪。}\BibleRef{若一3:8}我们也在\BibleBook{}{智慧篇}中读到：\Quote{但因魔鬼的嫉妒，死亡进入了世界。}\BibleRef{智2:24}因为这死亡是从魔鬼临到人，不是因为他们由他繁衍，而是因为他们效法他的榜样，所以紧接着就加了：\Quote{凡属于他的，都效法他。}\BibleRef{智2:25}因此，宗徒在同时提及罪和死亡时——它们是由一人藉自然遗传临到众人的——就把那一位立为这些的引入者，人类繁衍从他开始。\par

% structure: p0021 source=h2 target=subsection reason=relative-heading-rank; base=h2

\subsection{第十章——恩宠的类比}

毫无疑问，所有因不顺从而违反天主诫命的人，都效法亚当；但他在那些故意犯罪的人面前是一个榜样，同时又是所有带着罪出生的人的始祖。一切圣人也效法基督追求义德；因此我们刚才引用的那位宗徒说：\Quote{你们当效法我，如同我效法基督一样。}\BibleRef{格前11:1}但除了这种效法之外，祂的恩宠也藉着某种运行在我们内里光照并成义我们，关于这运行，那同一位宣讲祂圣名的传道者说：\Quote{可见栽种的算不得什么，浇灌的也算不得什么，只在那使生长的天主。}\BibleRef{格前3:7}因为这恩宠，祂甚至将受洗的婴儿接枝到祂的身体上，这些婴儿当然还不能效法任何人。因此，正如那位使众人得以复活的基督，除了将自己作为义德的榜样给效法祂的人以外，也赐给信祂的人祂圣神隐藏的恩宠，这恩宠祂秘密地倾注在婴儿身上；同样，那位使众人死亡的亚当，除了给故意违反主诫命的人提供一个效法的榜样以外，也藉着自己肉性的私欲偏情的隐藏败坏，败坏了所有从他血统而出的人。宗徒说这话完全出于这个原因，没有其它理由：\Quote{藉着一人，罪恶进入了世界，死亡藉着罪恶也进入了世界，这样死亡就临到了众人，因为众人都犯了罪。}\BibleRef{罗5:12}如果现在我\Emph{我}这样说，他们就会反对，并大声坚持说我用词和意思都不正确，因为这些字句出自一个普通人时，他们看不到任何意义，除非是他们在宗徒身上拒绝看到的那种意义。然而，既然这是那位他们所服从权威和教义的人的话，他们指责我们理解迟钝，同时却努力把那些意思清楚明白的话曲解成某种不可理解的意思。\Quote{藉着一人，}他说，\Quote{罪恶进入了世界，死亡藉着罪恶也进入了世界。}这表明是遗传，而非模仿；因为如果意思是模仿，他就会说：\Quote{藉着魔鬼。}但无人怀疑，他指的是那被称为亚当的第一个人：\Quote{这样，}他说，\Quote{它临到了众人。}\par

% structure: p0023 source=h2 target=subsection reason=relative-heading-rank; base=h2

\subsection{第十一章——实际罪与原罪的区别}

再次，在接下来的句子中：\BibleQuote{因为在它内众人犯了罪}，这陈述是多么谨慎、正确且不含糊！因为如果你理解那由一人进入世界的罪，\BibleQuote{在它内（罪）众人犯了罪}，那么显然，每个人特有的罪（即他们自己所犯、仅仅属于他们的罪）是一回事，而那使众人犯罪的一罪又是另一回事；因为众人都曾是那一人。然而，如果理解的不是罪，而是那一人，\BibleQuote{在他内（一人）众人犯了罪}，那么还有什么比这明确的陈述更清楚呢？我们确实读到，那些信靠基督的人在祂内因着那属灵恩宠的秘密共融与灵感而成义，这恩宠使每个与主结合的人与祂成为\BibleQuote{一神}，\BibleRef{格前 6:17}虽然祂的圣人也效法祂的榜样；但我能否找到任何类似的陈述，关于那些效法祂的圣人的人？能说任何人在\Person{保禄}或\Person{伯多禄}内，或在任何一位在天主子民中享有崇高权威的卓越人物内成义吗？我们无疑被称为在亚巴郎内受祝福，根据对他说的那句话：\BibleQuote{在你内万民都要蒙福}——是为了基督，祂按肉体是亚巴郎的后裔；这更清楚地在平行经文中表达出来：\BibleQuote{在你的后裔内万民都要蒙福。}我不相信任何人能在圣经中找到任何一处说一个人曾经或仍然\Quote{在魔鬼内}，尽管所有邪恶不敬的人都\Quote{模仿}他。然而，宗徒论及第一人时宣告说：\BibleQuote{在他内众人犯了罪}；\BibleRef{罗 5:12}而关于罪的传播仍有争论，人们用我不知道什么模糊的\Quote{模仿}论来反对它。\par

% structure: p0025 source=h2 target=subsection reason=relative-heading-rank; base=h2

\subsection{第十二章——法律不能除去罪}

也要注意接下来的内容。他说了\BibleQuote{因为在它内众人犯了罪}之后，立即加上\BibleQuote{因为直到法律，罪已在世界上。} \BibleRef{罗 5:13}这意味着罪甚至不能被法律除去，法律进入是为了使罪更显多，\BibleRef{罗 5:20}无论是自然律，在自然律下，每个人到达明辨是非的年龄时，只会将自己本罪加于原罪；或者是梅瑟赐予百姓的那法律。\BibleQuote{因为如果曾赐下能赐予生命的法律，那么义德确实应由法律而来。但圣经把所有的人都归在罪权之下，为使那藉着对耶稣基督的信德而来的恩许，赐予相信的人。} \BibleRef{迦 3:21} \BibleQuote{然而，没有法律的地方，罪就不被归算。} \BibleRef{罗 5:13}现在，\BibleQuote{即不被归算}，除了\BibleQuote{被忽略}，或\BibleQuote{不被视为罪？}这句话是什么意思？然而上主天主自己不视它为从未发生，因为经上记载：\BibleQuote{凡没有法律而犯罪的人，也要在没有法律而丧亡。} \BibleRef{罗 2:12}\par

% structure: p0027 source=h2 target=subsection reason=relative-heading-rank; base=h2

\subsection{第十三章[XI.]——宗徒短语\Quote{死亡的统治}的意义}

\BibleQuote{然而，}他说，\BibleQuote{死亡从亚当直到梅瑟也统治了，} \BibleRef{罗 5:14}——这就是说，从第一个人直到那由神权颁布的法律，因为连它也不能废除死亡的统治。现在，必须理解死亡\Quote{统治}，每当罪的罪责在人内如此主宰，以致阻止他们获得那唯一真实生命的永生，并将他们拉下甚至到那作为惩罚永恒的第二死亡时。这死亡的统治只有藉着救主的恩宠才能在任何人内被摧毁，这恩宠也曾在古时的圣人内运作，他们所有人，虽然在基督降生成人之前，却生活在对祂助佑之恩的关系中，而非对法律的文字，法律只知道命令，却不知道帮助他们。的确，在旧约中，那隐藏的（与时期完美正义的安排相符）如今在新约中启示出来。因此\BibleQuote{死亡从亚当直到梅瑟统治了}，在所有未被基督恩宠帮助的人内，为使他们内死亡的国度被摧毁，\BibleQuote{甚至在那些没有像亚当违命那样犯罪的人身上，} \BibleRef{罗 5:14}这就是说，他们尚未像亚当那样出于自己个人意志犯罪，但从他那里承受了原罪，\BibleQuote{他是那将来之人的预像，} \BibleRef{罗 5:14}因为在他内确立了对他未来后裔的定罪形式，他们应由他藉自然血统而生；如此，从一人众人出生归于定罪，这定罪除了在救主的恩宠内别无解救。我确实知道，有一些拉丁文圣经抄本这样读：\BibleQuote{死亡从亚当到梅瑟统治了那些按亚当违命之相似而犯罪的人；}但即使这个版本，那些如此读的人也引向完全同样的意义，因为他们理解那些在他内犯罪的人，是按亚当违命之相似而犯罪；所以他们被造成他的相似，不仅是作为由人而生的人，而且是作为由罪人而生的罪人，由必死者而生的必死者，由被定罪者而被定罪者。然而，拉丁译本所依据的希腊文抄本，全部或几乎全部，都有我最初引用的读法。\par

% structure: p0029 source=h2 target=subsection reason=relative-heading-rank; base=h2

\subsection{第十四章——恩宠的丰盛}

\Quote{但是，}他说，\BibleQuote{不像过犯，恩赐也是如此。因为如果因一人的过犯，众人都死了，那么天主的恩宠和那因一人耶稣基督所赐的恩赐，更加丰富地临到众人。}\BibleRef{罗 5:15}不是\Emph{更多人}，即更多人数，因为被成义的人并不比被定罪的人多；而是说\Emph{更加丰富地临到}，因为亚当从一次过犯产生了罪人，而基督却因祂的恩宠获得了无偿的宽恕，甚至宽恕了人们因实际过犯而主动加给原罪的那些罪。他在后文中更清楚地说明了这一点。\par

% structure: p0031 source=h2 target=subsection reason=relative-heading-rank; base=h2

\subsection{第十五章 [XII.]——所有人共有的那一罪}

但请更仔细地观察他所说的话：\BibleQuote{因一人的过犯，众人都死了。}为什么是因一人的罪，而不是因他们自己的罪呢？如果这段话应理解为\OriginalTerm{仿效}{imitation}，而非\OriginalTerm{遗传}{propagation}？但请看下文：\BibleQuote{并且，如同由一人犯罪而来的，恩赐也是如此；因为审判是由一人到定罪，恩宠却是由许多过犯到成义。}\BibleRef{罗 5:16}现在请他们告诉我们，在这些话中，\OriginalTerm{仿效}{imitation}在哪里？他说：\BibleQuote{由一人到定罪。}由一人，除了由一次过犯，还能是什么呢？这一点他在随后的话语中清楚地暗示了：\BibleQuote{但恩宠是由许多过犯到成义。}为什么审判是从一次过犯到定罪，而恩宠却是从许多过犯到成义？如果原罪是虚妄的，那么不仅恩宠将人从许多过犯中引到成义，审判也岂不同样从许多过犯中引到定罪？因为恩宠固然不宽恕许多过犯，而审判也同样不把许多过犯定罪。否则，如果人因一次过犯被定罪，是基于所有被定罪的过犯都是仿效那一次过犯而犯下的；那么，同样有理由认为，人也因一次过犯被解脱到成义，因为所有被宽恕给成义者的过犯都是仿效那一次过犯而犯下的。但这绝不是宗徒的意思，他说：\BibleQuote{审判确实是从一次过犯到定罪，但恩宠却是从许多过犯到成义。}我们这方面确实能理解宗徒，并看到审判之所以仅因一次过犯就被宣判为定罪，完全是因为即使在人身上只有原罪，也足以使他们定罪。因为那些在原始过犯上加添自己罪过的人，其定罪会更为严重（且各人越严重，按其罪过）；然而，即使那原本遗传给人的罪本身，不仅将人排除于天主的国度之外——这是婴儿无法进入的（正如他们自己所承认的），除非他们在死前领受基督的恩宠——而且使人远离救恩和永生，这救恩和永生只能是天主的国度，唯有与基督的共融才能引入。\par

% structure: p0033 source=h2 target=subsection reason=relative-heading-rank; base=h2

\subsection{第十六章 [XIII.]——死如何由一人来，生命如何由一人来}

由此我们得出，从亚当（我们在他内都犯了罪）那里，我们继承的不是我们所有的实际罪过，而仅仅是原罪；而从基督（我们在祂内都得以成义）那里，我们不仅获得原罪的赦免，也获得其他我们添加的罪过的赦免。因此说：\BibleQuote{不像由那犯罪的一人，恩赐也是如此。}因为审判确实从一次过犯而来——若不被赦免，就是原罪——足以将我们引向定罪；而恩宠却通过赦免许多罪过——也就是说，不仅原罪，而且所有其他的罪——将我们引至成义。\par

% structure: p0035 source=h2 target=subsection reason=relative-heading-rank; base=h2

\subsection{第十七章——罪人仿效谁}

\BibleQuote{因为如果因一人的过犯，死亡因那一人而掌权；那么那些领受丰盛恩宠和义德的人，更要因耶稣基督一人在生命中掌权。}\BibleRef{罗 5:17} 为什么死亡只因一人的罪而掌权？除非是因为所有人在那人内被死亡的锁链束缚，众人在他内都犯了罪，即使他们没有加添自己的罪？否则，死亡并不是通过一人因一人的罪而掌权，而是通过每个个别的罪人因众多过犯而掌权。因为如果人们因他人的过犯而死，是因为他们效法他，跟随他作为过犯的前驱，那么必然的结果是，而且\Emph{\InnerQuote{更甚}}，那人因另一人的过犯而死，而魔鬼在他之前过犯，以致说服他去犯罪。然而，亚当并没有使用任何影响力去说服他的跟随者；而那些据说效法他的许多人，实际上要么根本没有听说过他的存在，要么听说过他犯过那样的罪却不信。因此，正如我已经指出的，如果使徒在此处想说的是传染而非效法，他本会更加正确地将魔鬼设定为那\BibleQuote{一人}，从那\BibleQuote{一人}他说罪和死亡传给了所有人。因为说亚当是魔鬼的效法者更有理由，因为他有实际的煽动者；如果一个人可以效法一个从未使用过任何劝说的人，或者他完全不知道的人。那么经句\BibleQuote{他们领受丰盛恩宠和义德}意味着什么？岂不就是说赦罪的恩宠不仅赐予了众人所犯的罪，也赐予了人们实际犯下的其他过犯；并且在这些[人]身上白白地赐予了如此伟大的义德，以至于尽管亚当向劝说他犯罪的人让步，他们却连同样的诱惑者的胁迫也不屈服？再者，\BibleQuote{他们更要在生命中作王}这句话是什么意思？事实上，死亡的权势将更多人拖向永罚，除非我们理解在两种经文中真正提到的都是那些从亚当过渡到基督的人，换句话说，就是从死亡到生命的人；因为在永生中他们将无休止地作王，从而超越死亡的权势，这权势只在暂时且有终结地统治他们。\par

% structure: p0037 source=h2 target=subsection reason=relative-heading-rank; base=h2

\subsection{第十八章——唯有基督使人成义。}

\BibleQuote{因此，正如因一人的过犯，众人被定罪；同样，因一人的义行，众人被称义而得生命。}\BibleRef{罗 5:18} 这个\Quote{一人的过犯}，如果我们坚持\Quote{效法}，只能是魔鬼的过犯。然而，既然它明确地是针对亚当而非魔鬼说的，我们别无选择，只能理解这里暗示的是本性遗传的原则，而非效法。[XIV.] 现在当他说到基督时，\BibleQuote{因一人的\Emph{义行}}，他比他说\BibleQuote{因一人的\Emph{义德}}更明确地陈述了我们的教义；因为他提到的那义行是基督使不义之人成义的义行，祂并没有将其视为效法的对象，因为只有祂能够实现这一点。宗徒完全可以这样说并且说得对：\Quote{你们当效法我，如同我效法基督一样；}但他绝不能这样说：你们被我称义，如同我被基督称义一样；——因为可能并且实际上有许多义人值得效法；但唯有基督是义人又是使人成义者。因此经上说：\BibleQuote{那信使不义之人成义的天主的人，他的信德算为义德。}\BibleRef{罗 4:5} 如果有人有能力自信地宣布：\Quote{我使你成义，}那么他必然也能说：\Quote{信从我吧。}但天主的任何圣人都没有能力这样说，除了圣中之圣，祂说：\BibleQuote{你们信从天主，也当信从我；}\BibleRef{若 14:1} 因此，既然是祂使不义之人成义，那么对于信祂的人，祂的信德就被算为义德。\par

% structure: p0039 source=h2 target=subsection reason=relative-heading-rank; base=h2

\subsection{第十九章[XV.]——罪恶源于本性的传承，正如义德源于重生；\Quote{所有}人藉着亚当成为罪人，\Quote{所有}人藉着基督成为义人。}

若是只有模仿才使人通过亚当成为罪人，何以模仿不能同样使人通过基督成为义人？\Quote{因为，}他说，\BibleQuote{正如因一人的过犯，众人都被定罪；照样，因一人的成义，众人也得到生命的成义。}\BibleRef{罗 5:18}（依模仿论而言），则这里的\Quote{\Emph{一}}与\Quote{\Emph{一}}不应视为亚当与基督，而应是亚当与亚伯尔。因为虽有众多罪人在今世生活中先于我们，并且后来犯罪的人模仿了他们的罪，但他们却坚持：只提及亚当作为众人因模仿而犯罪的那一位，因为他是第一个犯罪的人。同理，亚伯尔当然也应当被提及，作为众人\Quote{在他内}同样因模仿而成义的那一位，因为他是第一个按照正义生活的人。然而，若认为有必要考虑与新约开端相关的某个关键时期，以基督为义人的首领和效法对象，那么出卖祂的犹达斯便应被列为罪人一类的首领。此外，如果基督是众人在其内成义的唯一一位，其根据并非仅仅因模仿祂的榜样使人成为义人，而是祂的恩宠借圣神重生人，那么亚当也是众人在其内犯罪的唯一一位，其根据并非仅仅因追随他的恶表使人成为罪人，而是那通过肉身产生人的处罚。因此出现了\Quote{\Emph{众人}}与\Quote{\Emph{众人}}的说法。因为通过亚当而生的人并非实际上与通过基督而重生的人完全相同；但宗徒的措辞仍然严格正确，因为正如无人不经亚当而参与肉身出生，同样无人不经基督而分享属灵重生。因为若有人能不经亚当而生在肉身中，同样若有人能不经基督而生在圣神内，那么显然\Quote{\Emph{众人}}就不能用于任何一类。但宗徒后来将这些\Quote{\Emph{众人}}描述为\Quote{\Emph{许多人}}；因为显然，在某些情况下，\Quote{众人}可能只是少数。然而，肉身的出生涵盖\Quote{\Emph{许多人}}，属灵的重生也包括\Quote{\Emph{许多人}}；虽然属灵的\Quote{许多人}在数量上少于肉身的\Quote{许多人}。但正如一类涵盖\Emph{一切}人，另一类也包括\Emph{一切}义人；因为在前一种情况下，无人能不经肉身出生而成为人，在后一种情况下，无人能不经属灵重生而成为义人；因此两者中都有\Quote{许多人}：\BibleQuote{因为正如因一人的悖逆，\Emph{许多人}成了罪人，照样，因一人的服从，\Emph{许多人}将成为义人。}\BibleRef{罗 5:19}\par

% structure: p0041 source=h2 target=subsection reason=relative-heading-rank; base=h2

\subsection{第二十章——唯独原罪由自然出生而来}

\BibleQuote{再者，法律介入，好叫过犯增多。}\BibleRef{罗 5:20}人们现在凭自己的任性在原罪上增加了过犯，并非通过亚当；但连这也被基督消除和补救，因为\BibleQuote{那里罪过增多了，恩宠却越格外丰富；好叫罪过怎样以死亡为王}\BibleRef{罗 5:21}——即使那罪并非人从亚当继承的，而是他们自己的意志增加的——\BibleQuote{照样，恩宠也借着正义为王，进入永生}\BibleRef{罗 5:21}。然而，除了基督之外，还有别的义，如同除了亚当之外还有别的罪。因此，在他说\BibleQuote{正如罪过怎样以死亡为王}之后，他没有在同一子句中加上\Quote{\Emph{因一人}}或\Quote{\Emph{因亚当}}，因为他已经谈到那当法律进入时增多的罪，那当然不是原罪，而是人自己任意犯的罪。但在他说完\BibleQuote{照样，恩宠也借着正义为王，进入永生}之后，他立即加上\BibleQuote{借着我们的主耶稣基督}\BibleRef{罗 5:21}；因为借着肉身的出生，只染上那属于原罪之罪；但借着圣神的重生，却带来赦免，不仅是原罪，也包括人自己意愿和实际所犯的罪。\par

% structure: p0043 source=h2 target=subsection reason=relative-heading-rank; base=h2

\subsection{第二十一章[第十六章]——未受洗的婴儿被定罪，但处罚最轻；亚当之罪的刑罚：失去其身体的恩宠}

因此，可以正确地断言，那些未受圣洗便离世的婴孩将承受最轻的定罪。所以，那教导说他们不会受定罪的人，大大欺骗了自己和他人；因为宗徒说：\BibleQuote{审判出于一次的过犯而至定罪，} \BibleRef{罗 5:16} 稍后又说：\BibleQuote{因一人的过犯，众人被定罪。} \BibleRef{罗 5:18} 当亚当因不服从天主而犯罪时，他的身体——虽然是自然而有死的身体——便失去了那曾使身体各部分顺从灵魂的恩宠。于是，在人身上产生了与野兽相同的情欲，这些情欲使人感到羞耻，并使人对自己赤身露体感到羞愧。\BibleRef{创 3:10} 然后，由于一种突然注入的致死的腐败在人身上产生的疾病，人失去了被造时所拥有的生命的稳定，并且因生命各阶段的变迁，最终走向死亡。无论他们此后活了多少年，他们从接受死亡法律的那一天起就开始死亡，因为他们不断趋向衰老。因为那东西没有片刻的稳定，而是不停地流逝，通过持续的变化，可察觉地走向一个终点，这终点不是带来完美，而是彻底衰竭。如此，天主所说的就应验了：\BibleQuote{在你吃的那一天，你必死。} \BibleRef{创 2:17} 因此，由于肉体的悖逆和罪与死的法律，凡从肉体生的，都需要属灵的重生——不仅为进入天国，也为脱离罪的罚。所以，人一方面因第一个亚当而从肉体出生，并处于罪与死之下；另一方面在圣洗中重生，与第二个亚当的义德和永生相连；正如\WorkTitle{德训篇}所记载的：\BibleQuote{罪恶始于女人，因着她，我们都死。} \BibleRef{德 25:24} 现在，无论是指女人还是指亚当，两者都属于第一个人；因为（我们知道）女人出自男人，两人成为一体。因此也记载着：\BibleQuote{二人要成为一体；所以} 主说，\BibleQuote{他们不再是两个，而是一体。} \BibleRef{玛 19:5}\par

% structure: p0045 source=h2 target=subsection reason=relative-heading-rank; base=h2

\subsection{第二十二章 [XVII.]——不应将本罪归给婴孩。}

因此，那些说婴孩受圣洗的理由是为了得罪他们自己一生所犯的罪，而非从亚当所遗传的罪的人，不难被驳倒。因为每当这些人稍加内省，不受任何争论精神的影响，反思自己说法的荒谬性——它事实上是多么不值得严肃讨论——他们就会立刻改变看法。但如果他们不这样做，我们也不至于对人的常识如此绝望，以至于担心他们会诱使他人接受他们的观点。如果我没有弄错，他们自己是被对其他某种理论的偏见所驱使而接受这种意见；并且，因为他们觉得自己必须承认罪被赦免给受圣洗者，而不愿意承认那罪是从亚当遗传的（他们承认那罪被赦免给婴孩），所以他们不得不将本罪归咎于婴儿本身；仿佛通过将这项控诉加于婴儿，他本人在被指控且无法回应攻击者时能变得更加安全一样！然而，按照我的建议，让我们绕过这样的反对者；事实上，我们既不需要言语也不需要圣经引文来证明婴孩在其生活中的行为是无罪的；这段生命，由于他们的出生如此之近，是在他们自身之内度过的，因为它逃脱了人类感知的认知，而人类感知没有任何数据或支持来就此主题进行任何争论。\par

% structure: p0047 source=h2 target=subsection reason=relative-heading-rank; base=h2

\subsection{第二十三章 [XVIII.]——他驳斥那些声称婴孩受圣洗不是为了得罪之赦，而是为了获得天国的人。}

但那些人提出一个问题，并似乎提出一个值得考虑和讨论的论点，他们说新生婴儿领受圣洗不是为了赦免罪过，而是因为他们的生育不是属灵的，为使他们得以在基督内被创造，成为天国的分享者，并借此成为天主的儿女和继承人，以及基督的同继承人。然而，当你问他们，那些未受洗、未成为基督的同继承人和天国的分享者的人，是否至少在死者复活中享有永生的福分呢，他们极其困惑，找不到出路。因为哪个基督徒会允许有人说，任何人若不重生在基督内——祂正是通过圣洗来实现这一点，当时这圣事特意为在永生希望中重生而设立——就能获得永生的救恩呢？因此宗徒说：\Quote{祂救了我们，并不是由于我们本着义德所立的功劳，而是按照祂的怜悯，借着重生的水洗。}\BibleRef{铎 3:5} 然而，他说这救恩在于希望，当我们生活在此下土时，他说：\Quote{我们得救是在于希望，但看得见的希望不是希望，因为人看见的，为什么还要希望呢？但我们若希望那看不见的，就必须坚忍等待。}\BibleRef{罗 8:24} 那么，谁敢如此大胆地断言，没有宗徒所说的重生，婴儿就能获得永生的救恩，好像基督并未为他们而死？因为\Quote{基督为\Emph{不虔敬者}死了。}\BibleRef{罗 5:6} 然而，对于那些（显然）一生从未行过任何不虔敬行为的人，如果他们在原初本性中也不受任何罪债的束缚，那么基督如何为他们死了呢？祂是为那不虔敬者死的。如果他们未受原罪的任何疾病伤害，为何那些虔诚焦虑的人要跑到祂面前，特意将他们带到基督这位医生那里，以领受永生圣事呢？为何不在教会中对他们说：把这些无罪者带走，\Quote{健康的人不需要医生，而是有病的人；}——基督 \Quote{来不是召义人，而是召罪人。}\BibleRef{路 5:31} 关于基督的这样一种虚构，在教会中从未听说过，从未听说过，也永远不会听到。\par

% structure: p0049 source=h2 target=subsection reason=relative-heading-rank; base=h2

\subsection{第二十四章 [XIX.]——婴儿作为罪人得救。}

并且没有人应当认为，婴儿应该被带到圣洗前，是因为他们既非罪人，也非义人；那么，怎么有些人提醒我们，主称赞这幼小的年龄为有功的，说：\Quote{让小孩子到我这里来，不要阻止他们，因为天国正是属于这样的人。}\BibleRef{玛 19:14} 因为如果这话（\Quote{这样的人}）不是由于谦卑的相似（因为谦卑使人成为孩子），而是由于儿童值得称赞的生活，那么婴儿当然必须是义人；否则，就不能正确地说\Quote{天国正是属于这样的人}，因为天国只属于义人。但也许，对主的话的理解并不正确，以为祂在说\Quote{天国正是属于这样的人}时称赞了婴儿的生命；其实\Emph{那}可能是它们的真实意义，使基督引用幼小的年龄作为谦卑的比喻。即便如此，也许我们必须回到我刚才提到的论点：婴儿应当受洗，因为虽然他们不是罪人，但他们也不义。但当祂说：\Quote{我来不是召义人}时，好像回应了这一点，那么祢来召谁呢？祂紧接着说：\Quote{——而是召罪人悔改。}因此，无论他们多么义，如果他们也不是罪人，祂就没有来召他们，因为祂自称：\Quote{我来不是召义人，而是召罪人。}他们因此似乎\Emph{不仅是徒劳地，而且是邪恶地}冲向祂的圣洗，而祂并不邀请他们——这种想法，愿天主禁止我们怀有。那么，祂召唤他们，如同一位医生，不需要健康的人，而是有病的人；祂来不是召义人，而是召罪人悔改。现在，既然婴儿并不受自己实际生活中任何罪的约束，那么原罪的罪责就在他们内被那借着重生水洗拯救他们的恩宠所治愈。\par

% structure: p0051 source=h2 target=subsection reason=relative-heading-rank; base=h2

\subsection{第二十五章——婴儿被描述为信徒和忏悔者。只有罪将天主与人分开。}

有人会说：那么婴儿如何被召悔改呢？像他们这样的人怎能悔改任何事呢？对此的回答是：如果他们不能被称为忏悔者，因为他们没有悔改的意识，那么他们也不能被称为信徒，因为他们同样没有相信的意识。但如果他们被称为信徒是正确的，因为他们某种程度上借父母的话语宣认信仰，那么为什么他们不被视为在之前就是忏悔者呢，当他们借同样父母的宣认表明弃绝魔鬼和这个世界时？这一切都是在希望中、在圣事的力量和主赐予教会的天主的恩宠中完成的。但谁不知道，受洗的婴儿如果到了理性年龄后不信并不断绝不法的欲望，他所领受的便无法获益？然而，如果婴儿在领受圣洗后离开现世生命，原罪所牵涉的罪责被消除，他将在那永恒不变的真理之光中得以成全，这光照耀着在造物主面前成义的人。因为只有罪将人与天主分开；这些罪借着基督的恩宠被消除，借着祂，作为中保，我们得以和好，当祂使不虔敬者成义时。\par

% structure: p0053 source=h2 target=subsection reason=relative-heading-rank; base=h2

\subsection{第二十六章 [XX.]——除非受洗，没有人能正当地来到主的祭台前。}

现在他们从主的话中感到惊恐，因为主说：\BibleQuote{人若不重生，就不能见天主的国；}\BibleRef{若 3:3} 因为在主亲自解释这段经文时，祂断言：\BibleQuote{人若不由水和圣神而生，不能进天主的国。}\BibleRef{若 3:5} 于是他们试图将救恩和永生的恩赐归于未受洗的婴孩，凭借他们的无辜的功绩，但同时由于他们未受洗，将他们排除在天国之外。但这样的假设是多么新奇而惊人，好像可能有不与基督同嗣、没有天国的救恩和永生！当然他们有避难所，可以逃避和隐藏，因为主没有说：\Quote{人若不由水和圣神而生，不能有生命}，而是说：\Quote{不能进天主的国。} 如果主确实说了前者，那么就不会有一刻的怀疑。那么，让我们消除怀疑；现在让我们听主的话，而不是人的观念和猜测；让我们，我说，听主所说的——不是关于圣洗圣事，而是关于祂自己圣体圣事，只有受洗的人才有权走近：\BibleQuote{你们若不吃我的肉，不喝我的血，在你们内便没有生命。}\BibleRef{若 6:53} 我们还需要什么呢？对此能提出什么回答呢，除非是那种总是抵制明显真理的固执？\par

% structure: p0055 source=h2 target=subsection reason=relative-heading-rank; base=h2

\subsection{第27章——婴孩必须领受基督。}

然而，有谁胆敢说这话与婴孩无关，他们不领受祂的体血也能有生命——因为祂不是说：\Quote{除非\Emph{一人}吃}，而是说：\Quote{除非\Emph{你们}吃}；好像祂是对那些能听到和理解的人说话，婴孩当然不能这样？但说这话的人是不经心的；因为除非\Emph{所有人}都包含在这句话中，即没有圣子的体血，人就不能有生命，那么即使成年人为此焦虑也是徒劳的。因为如果你只注意主说话时的字面，而不注意意思，这段话似乎只是对当时在场的人说的；因为祂不是说\Quote{除非一人吃}，而是\Quote{除非你们吃}。关于这一点，祂在同样的上下文中说的话又如何呢：\BibleQuote{我要赐给的食粮，就是我的肉，是为世界的生命而赐的。}\BibleRef{若 6:51} 因为根据这句话，我们发现这圣体圣事也属于我们，我们当时并不存在，但我们不能说我们不属于\Quote{世界}，基督为世界的生命赐下了祂的肉。谁又能怀疑在\Quote{世界}这个词中包含了所有因出生而进入世界的人呢？正如祂在另一处所说：\Quote{今世之子也娶也嫁。} 由此可知，甚至为了婴孩的生命，祂的肉也被赐下了，祂为世界的生命赐下了肉；并且即使他们若不食圣子的肉，他们也不会有生命。\par

% structure: p0057 source=h2 target=subsection reason=relative-heading-rank; base=h2

\subsection{第28章——受洗的婴孩属于信者；未受洗的属于丧亡者。}

因此还有另一句话：\BibleQuote{父爱子，且把一切交在祂手中。信从子的人，有永生；不信从子的人，不得见生命，天主的震怒却常在他身上。}\BibleRef{若 3:35} 现在，婴孩应放在哪一类中——是信从子的人中，还是不信从子的人中？有些人说，两者都不属，因为他们还不能信，所以也不应被视为不信者。然而，教会的规矩并非如此，因为它将已受洗的婴孩算在信者之中。既然他们因如此伟大圣事的卓越与管理，仍被算在信者之列，尽管他们自己的心和口并没有实际行出属于信德行为和宣认的事；那么，无疑那些缺乏圣事的人必须被归入不信从子的人中，因此，如果他们离开此生而没有这恩宠，他们将要面对关于这样的人所写的：他们不得见生命，天主的震怒常在他们身上。对那些显然没有自身之罪的人，如果他们不因原罪而有罪，这又如何可能呢？\par

% structure: p0059 source=h2 target=subsection reason=relative-heading-rank; base=h2

\subsection{第29章[XXI.]——为何有些人得救，其他人却不得救，是一个不可探究的奥秘。}

现在，祂没有说\Quote{天主的义怒\Emph{将来临}到他身上}，而说\Quote{\Emph{常在他身上}}，这有很大的深意。因为，除了藉着我们的主耶稣基督，天主的恩宠之外，没有什么能将我们从这义怒中（我们所有的人都在罪中陷入这义怒，关于这义怒，宗徒说：\BibleQuote{因为我们从前也生来就是义怒之子，和别人一样}\BibleRef{弗 2:3}）拯救出来。这恩宠为什么临到这个人而不临到那个人，原因可能是隐藏的，但绝不可能是非义的。因为\BibleQuote{难道天主有不义吗？绝对不可。}\BibleRef{罗 9:14}。但我们必须先顺服圣经的权威，好使我们各人藉着信德获得知识和理解。因为经上并非徒然说：\Quote{你的判断如同深渊。}这个\Quote{深渊}的深，宗徒似乎怀着战栗之情，在那惊叹中指出了：\Quote{啊，天主的智慧和知识的财富多么深远！}他之前已经指出了这奇妙深渊的意义，当他说：\BibleQuote{因为天主把众人都禁锢在不信服之中，是要怜悯众人。}\BibleRef{罗 11:32}。然后，仿佛被这深渊的可怕所震慑，他说：\BibleQuote{啊，天主的智慧和知识的财富多么深远！他的判断是多么不可测量，他的道路是多么不可探察！谁曾知道上主的心意？或者谁曾当过他的顾问？谁又先给了祂什么，以致祂应当还报他呢？因为万物都出于祂、依靠祂、归于祂。愿光荣归于祂，至于永世。阿们。}\BibleRef{罗 11:33}。那么，我们讨论天主判断的公义以及祂白白赐予的恩宠（这恩宠，因为人没有先行的功德配得它，所以不能是偏袒的或不义的）的能力是多么微不足道啊！这恩宠当赐予不配的人时，并不令我们困扰，就像当它被拒绝给同样不配的人时一样。\par

% structure: p0061 source=h2 target=subsection reason=relative-heading-rank; base=h2

\subsection{第三十章——为何一人受洗而另一人不，何尝不是同样深不可测。}

现在，那些认为婴孩在未获得基督恩宠而离世，不仅被剥夺了天主的国（他们自己也承认，只有藉着圣洗重生的人才能进入），而且也被剥夺了永生和救恩，这是不义的人，当他们问：两个人的状况相同，为何一人得以从原罪中获得释放而另一人却不能，这怎么可能是正义的？他们可以根据自己的意见来回答自己的问题：按照他们的看法，为何同样是两个人，一人得以受洗而进入天主的国，另一人却未得到如此恩待，尽管两人的情况相同，这常常是公义和正当的？因为，如果这个问题困扰他：为何两个同样生来有罪的人，一个藉着圣洗从捆绑中得到释放，而另一个没有获得这恩宠，所以不得释放？他为何不被另一事实同样困扰：两个生来无罪的人，一个接受圣洗，得以进入天主的国，而另一个没有接受，因此不能接近天主的国？在这两种情况下，人们都会想起宗徒的惊叹：\Quote{啊，天主的智慧和知识的财富多么深远！}再者，请告诉我，为什么在已受洗的婴孩中，一个被接走，使他的心智不致因恶行而改变\BibleRef{智 4:11}，而另一个则存活下来，注定成为不虔敬的人？假设两者都被接走，难道他们不都进入天国吗？然而，在天主并没有不义。\BibleRef{罗 9:14}为什么没有人被触动，没有人因这些深奥的事而发出惊叹，鉴于有些儿童被邪魔折磨，而另一些没有遭受这种污染，还有的，如耶肋米亚，甚至在母胎中就被圣化\BibleRef{耶 1:5}？相反，所有的人，如果有原罪，就都有同样的罪债；或者，如果有原罪，就同样是无辜的？这种巨大的差异从何而来，除非是因为天主的判断不可测量，祂的道路不可探察？\par

% structure: p0063 source=h2 target=subsection reason=relative-heading-rank; base=h2

\subsection{第三十一章[第二十二篇]——他驳斥那些认为灵魂因在另一状态中所犯的罪，而被投入与其功过相称的身体，在其中或多或少受折磨的人。}

但也许必须重新提出那如今已被驳斥和摒弃的见解，即：灵魂曾经在天上的居所中犯过罪，然后按阶段和程度降下，进入与它们功过相称的身体，并作为对前世生活的惩罚，或多或少地受到肉体惩戒的折磨。对于这种见解，圣经确实提出了最明显的矛盾；因为当它推荐天主恩宠时，它说：\BibleQuote{孩子还没有出生，也没有行善或作恶，为使天主按拣选所定的旨意得以坚立，不是由于行为，而是由于召叫的那一位，就有话说：年长的要服事年小的。}\BibleRef{罗 9:11}然而，那些持这种见解的人实际上无法逃脱这个问题的困惑，反而被它们困扰和束缚，被迫像其他人一样呼喊：\Quote{深哉！}因为怎么会发生这样的事：一个人从最早的童年起就表现出更大的节制、智力优秀和节德，并在很大程度上克服情欲，憎恨贪财，厌恶奢侈，在其他美德上达到更高的卓越和资质，却生活在这样一个地方，以至于无法听到宣讲基督的恩宠？——因为\BibleQuote{他们怎能呼求他们尚未相信的那一位？他们怎能相信他们尚未听见的那一位？没有宣讲者，他们怎能听见？}\BibleRef{罗 10:14}而另一个人，虽然心智迟钝，沉溺于情欲，被耻辱和罪恶覆盖，却被如此引导，以至于听见，相信，受洗并被接走——或者，如果允许他在这里停留更久，则以一种会给他带来赞誉的方式度过余生。那么，这两个人从哪里获得如此不同的功德呢？——我不是说一个人相信而另一个人不相信，因为那是人自己的意志之事；而是一个人得以听见以便相信，另一个人却没有听见，因为这不是人力所能及的？我说，他们从哪里获得不同的功德呢？如果他们确实曾在天上度过一部分生命，以至于被推下或沉入这个世界，并居住在与其前世生活相宜的身体容器中，那么当然应该认为，那个没有太应受此负担的人，在现在的有死身体之前，曾度过更好的生活，以至于既有良好的性情，又受到他能轻易克服的温和欲望的困扰；然而，他却不应得那唯一能使他从第二次死亡的毁灭中得救的恩宠被宣讲给他。而另一个人，他被一个粗重的身体所阻碍，作为一种惩罚——如他们所推测的——因为更坏的功德，因此拥有更迟钝的情感，当他处于情欲的猛烈炽热中，屈服于肉体的陷阱，并通过他邪恶的生活加重了他先前导致他陷入如此境地的罪过，通过更为放纵的世俗欢乐——要么在十字架上听见：\BibleQuote{今天你就要同我一起在乐园里。}\BibleRef{路 23:43}，要么依附于某位宗徒，通过他的宣讲而改变，并藉着重生的洗涤而得救——好叫罪恶在哪里增多，恩宠在那里就更加增多。我不知道那些想凭借人的猜测来维护天主的公义、且对恩宠的深奥一无所知、却编织了不可能寓言的网罗的人，对此能给出什么回答。\par

% structure: p0065 source=h2 target=subsection reason=relative-heading-rank; base=h2

\subsection{第三十二章 某些愚昧人和头脑简单之人的情形}

现今关于人奇特的召叫，可以谈论不少——无论是我们读到的，还是亲身经历的——这些事足以推翻那些人的意见，他们认为灵魂在拥有身体之前，曾经过某种独特的生活，在这种生活中，他们必须来到今生，并根据各自功过的差异，在现世经历善恶。然而，我急于结束这部作品，不容我在此更多停留。不过有一件事，我在许多事例中发现它非常奇特，我绝不保持缄默。若我们追随那些人，他们认为灵魂根据前一生活于属天身体中的功过（在取得现世身体之前），而受到或多或少的粗重属地上体的压制，那么谁不会断言，那些人在今世之前犯有特别重大的罪，以致他们配得完全失去心智之光，生来就具有类似野兽的本能？——他们（我不说心智极为迟钝，因为这点也常被用于他人）是如此愚笨，以致用白痴的怪相取悦聪明人，俗人用一个源自希腊文的名称称他们为\OriginalTerm{愚人}{Moriones}？然而曾有一个这样的人，他是如此基督徒，以至于他对自己所受的任何伤害都耐心到一种奇怪的愚蠢地步，但对基督之名，或他所浸润的宗教，若受到丝毫侮辱，他就极不耐烦；每当他那群轻浮聪明的听众为了激怒他的耐心而有时亵渎圣名时，他绝忍不住用石头砸他们；在这些场合，他甚至对贵族也不留情。那么，依我看来，这样的人被预定并被带到世上，为的是让那些有能力的人明白：天主的恩宠和那\Quote{风随意向哪里吹}的圣神\BibleRef{若 3:8}并没有越过仁慈之子的任何能力，同样也没有越过革赫纳之子的任何能力，好使\Quote{凡夸耀的，应因主而夸耀}\BibleRef{格前 1:31}。然而，那些主张灵魂根据前一生活的功过而领受不同身体（或粗或细），并且人的心智能力也根据同样的功过而各异（因此有的更敏锐，有的更迟钝），而且天主的恩宠也根据前生的功过而分施以解救人们脱离罪恶的人——对于这个人，他们会说什么呢？他们如何能归给他一个如此可耻的前生，以致他应当生为白痴，同时又如此功高，以致在领受基督恩宠的分配上比许多最敏锐的人更受青睐呢？\par

% structure: p0067 source=h2 target=subsection reason=relative-heading-rank; base=h2

\subsection{第33章——基督甚至也是婴孩的救主和救赎者。}

因此，让我们屈服并赞同圣经的权威，它是不会受骗也不会骗人的；我们既不相信尚未出生的人曾行过任何善恶以造成他们道德功过的差异，就绝不要怀疑所有人都在罪恶之下，这罪由一人进入了世界，并传给了众人；除赖我们的主耶稣基督的天主的恩宠外，没有什么能解救我们。[XXIII.] 祂医治的来临是为病人，不是为健康的人：因为祂来不是召义人，而是召罪人；凡不从水和圣神重生的人，不能进入祂的国；除了在祂的国内，没有人能达到救恩和永生——因为不信从子、不吃祂肉的人，不得有生命，天主的义怒常在他身上。现在，从这罪、这病、这天主的义怒（凡有原罪的人，即使因年幼没有本罪，也按本性是义怒之子）中，没有谁能解救他们，除了一天除免世界罪孽的天主羔羊\BibleRef{若 1:29}，除了那来不是为健康的人、而是为病人的医生，除了那救主，关于祂曾对人族说：\Quote{今天为你们诞生了一位救主}\BibleRef{路 2:11}，除了那以自己的血抹去我们债务的救赎者。因为谁敢说基督不是婴孩的救主和救赎者呢？但如果他们里面没有原罪的疾病，祂救他们脱离什么？如果他们因从第一人起源而未曾被卖于罪下，祂赎他们脱离什么？那么，不要凭我们自己的意见，说未受基督圣洗的婴孩有永生；因为在圣经中并没有应许，而圣经应凌驾于一切人为的权威和意见之上。\par

% structure: p0069 source=h2 target=subsection reason=relative-heading-rank; base=h2

\subsection{第34章 [XXIV.]——在迦太基的基督徒称圣洗为救恩，称圣体为生命。}

迦太基的基督徒对圣事有一个极好的称呼，他们说圣洗不是别的，就是\Quote{救恩}，而基督身体的圣事不是别的，就是\Quote{生命}。然而，这源自何处？我认为是源自那原始的、宗徒的传统，基督的教会以此持守一个固有原则：没有圣洗和分享主的晚餐，任何人都不可能获得天主的国、救恩和永生。圣经也根据我们已经引用的经文证实了这一点。因为他们称圣洗为\Emph{救恩}的看法，与所写的\BibleQuote{祂藉着重生洗涤\Emph{拯救了我们}}\BibleRef{铎 3:5}，或伯多禄的话\BibleQuote{与此预表相应的圣洗现在也\Emph{拯救你们}}\BibleRef{伯前 3:21}，有什么不同呢？那些称主的晚餐圣事为\Emph{生命}的人，又说了什么，不就是所写的\BibleQuote{我是从天上降下来的\Emph{生活的}食粮}\BibleRef{若 6:51}，以及\BibleQuote{我所要赐给的饼，就是我的肉，是为\Emph{世界的生命}}\BibleRef{若 6:51}，和\BibleQuote{你们若不吃人子的肉，不喝祂的血，在你们内便\Emph{没有生命}}\BibleRef{若 6:53}吗？因此，既然如此众多且神圣的见证一致认为，没有圣洗和主的身体与血，任何人都不能指望救恩和永生，那么向婴儿许诺这些恩惠而不给他们这些圣事，便是徒劳的。此外，如果只有罪使人远离救恩和永生，那么这些圣事在婴儿身上所能除去的，就只能是罪的罚——关于这有罪的本性，经上写着：\BibleQuote{没有人是洁净的，即使他的生命只有一天}\BibleRef{约 14:4}。圣咏作者也如此呼喊：\BibleQuote{看，我是在罪恶中生成的，我母亲在罪恶中怀了我！}这话要么是以我们共同人性的口吻说的，要么如果只是达味说自己，那并不意味着他是从奸淫而生，而是在合法的婚姻中。因此，我们不应怀疑，就连为尚未受洗的婴儿，那宝血也曾流下，这血在它实际倾流之前，就这样在圣事中赐予并应用，以致经上说：\BibleQuote{这是我的血，将为多人倾流，以赦免罪过}\BibleRef{玛 26:28}。如今，那些不承认他们处于罪下的人，就否认有任何解放。因为如果人们被认为不受任何罪恶的束缚，那么他们被解放出来是为什么呢？\par

% structure: p0071 source=h2 target=subsection reason=relative-heading-rank; base=h2

\subsection{第35章——除非婴儿受洗，否则他们仍留在黑暗中}

基督说：\BibleQuote{我来到世界上，为光，凡信我的，不留在黑暗中}\BibleRef{若 12:46}。这段话向我们表明了什么？岂不是说，凡不信祂的人都在黑暗中，而藉着信祂，就能脱离这种永久的黑暗状态？我们所说的\Emph{黑暗}是指什么？就是罪。无论它还可能包含什么意思，至少不信基督的人将\Quote{留在黑暗中}——这当然是一种受罚的状态，不像黑夜的黑暗那样是为了活物的休息。[XXV.] 所以，婴儿若不通过为此目的而神圣设立的圣事进入信徒的行列，无疑将留在这黑暗中。\par

% structure: p0073 source=h2 target=subsection reason=relative-heading-rank; base=h2

\subsection{第36章——婴儿并非一出生就被光照}

有些人却认为，孩子一出生就被光照了；他们从这段经文得出这个意见：\BibleQuote{那普照每个人的真光，正在进入这世界}\BibleRef{若 1:9}。如果真是这样，那就很奇怪了：那些被在起初就与天主同在，且是天主的独生子这样光照的人，怎么不能进入天主的国，不能成为天主的承继者和与基督同承继者呢？因为这种继承不通过圣洗就不会赐予他们，即使持这种见解的人也承认这一点。那么，如果他们（虽然已被光照）仍不适合进入天主的国，他们至少应该欣然接受那使他们适合的圣洗；但奇怪的是，我们看到婴儿是多么不愿意接受圣洗，甚至大声哭闹反抗。我们轻视他们这个年纪的无知，所以尽管他们抗拒，我们仍充分施行我们知道对他们有益的圣事。而且，如果他们的心智已经被那真光——即天主的圣言——光照了，那么宗徒为什么说：\BibleQuote{在理解上不要作孩子}\BibleRef{格前 14:20}呢？\par

% structure: p0075 source=h2 target=subsection reason=relative-heading-rank; base=h2

\subsection{第37章——天主如何光照每一个人}

因此，福音中的那句话：\BibleQuote{那普照每个人的真光，正在进入这世界}\BibleRef{若 1:9}，其意思是：没有人被光照，除非被那真理之光——就是天主——光照；所以没有人应该认为，他是由他所听信的老师光照的，即使那位导师——我不说是什么伟人——甚至是天使本身。因为真理之言是通过身体声音的职务从外部应用于人，但\BibleQuote{栽种的不是什么，浇灌的也不是什么，而是那使之生长的天主}\BibleRef{格前 3:7}。人确实听到讲话者，无论是人还是天使，但为了他能够领悟并知道所说的是真实的，他的内心被那永远存在、甚至照耀黑暗的光所浸透。然而，就像太阳不被盲人看见，尽管他们仿佛被它的光线包裹，同样，真理之光不会被愚昧的黑暗所理解。\par

% structure: p0077 source=h2 target=subsection reason=relative-heading-rank; base=h2

\subsection{第38章——\Quote{光}一词的含义}

但是，为什么在说了\BibleQuote{照耀每一个人}之后，他又加上\BibleQuote{来到这世界上}\BibleRef{若 1:9}——正是这个分句引发了这样的意见，认为祂照明了新生婴儿的心灵，而此时他们刚从母腹中出生不久？这些词语在希腊文中无疑是如此排列的，以至于它们可以被理解为表达光本身\Quote{来到了世界上}。然而，如果这个分句必须被理解为表达那进入这世界的人，我认为它要么是一个简单的短语，就像在圣经中常见的许多其他短语一样，删除它并不损害整体意思；要么，如果它被视为一个独特的附加说明，那么它或许是为了区分属灵的光照与那身体的光照——后者或藉着天上的光体，或藉着普通的火光，照亮肉身的眼睛。因此，他提到内在的人来到世界上，因为外在的人具有属物质的本性，就如这世界本身一样；就好像他说：\Quote{照耀每一个进入身体的人}，这与所记载的是一致的：\BibleQuote{我获得了良好的精神，并以无玷的身体来到世上。}\BibleRef{智 8:19}或者，这段经文——\Quote{照耀每一个来到世界上的}——如果它是为了表达某种区别而添加的，那么或许可以意味着：照耀每一个内在的人，因为内在的人，当他成为真正有智慧时，只有那真光才能光照他。或者，再一次，如果意图是用\Emph{光照}一词来指称理智本身，它使人的灵魂被称为理性的（而这理智，虽然尚安静且仿佛沉睡，但仍在婴儿内隐藏着，是与生俱来的，可说是内在的），仿佛它是内在眼睛的创造，那么，当灵魂被造时，这光照就发生了，这不能否认；而设想它在人来到世界时发生，也并非荒谬。然而，虽然他的眼睛已被创造，他本人却必须仍在黑暗中，如果他不相信那一位所说的：\BibleQuote{我来到世界上，作为光，使凡信我的，不住在黑暗中。}\BibleRef{若 12:46}至于在婴儿身上，通过圣洗圣事，这事发生，慈母教会毫不怀疑；教会为他们使用母亲的心和口，使他们能浸润于神圣奥迹中，因为他们尚不能以自己的心\BibleQuote{相信而成义}，也不能以自己的口\BibleQuote{宣认而得救}\BibleRef{罗 10:10}。的确，在信徒中，没有人会犹豫称这样的婴儿为\Emph{信徒}，仅仅因为这个称呼来源于相信的行动；因为虽然他们自己不能做出这样的行动，但有其他人在圣事中为他们作保。\par

% structure: p0079 source=h2 target=subsection reason=relative-heading-rank; base=h2

\subsection{第三十九章 [第二十六]——得出的结论：所有人都涉及原罪}

这将是冗长的，如果我们以同样的篇幅讨论所有与这问题相关的证言。我想更便捷的方式是，简单地将可能出现的或那些足以显明真理的经文收集在一起，即主耶稣基督成了肉身，取了奴仆的形体，\BibleQuote{服从至死，且死在十字架上}\BibleRef{斐 2:8}，除了藉此极其仁慈恩宠的救恩计划，将生命赐给所有那些作为祂身体的肢体、被嫁接在祂身上、成为祂的头以得享天国的人之外，没有其他原因；祂来拯救、释放、救赎、光照他们——他们先前已被陷于死亡、软弱、奴役、囚禁和罪的黑暗中，在罪的作者魔鬼的权势之下；如此，祂成为天主与人之间的中保，藉着祂（在我们不敬神明的敌意因祂恩助而终止之后），我们得以与天主和好，获得永生，从那威胁我们的永死中被救出。当这一点被足够证据显明后，那些不需要基督藉降卑所执行的这一救恩计划的人，就不能与那计划相关，因为他们不需要生命、救恩、解放、救赎和光照。既然圣洗属于这计划，我们在洗礼中与基督同葬，为要成为祂的肢体（即那些信祂的人），那么，对那些不需要藉中保所获得的宽恕与和好益处的人，圣洗当然是不必要的。现在，既然他们承认给婴儿施洗的必要——无法反对普世教会的权威，这权威无疑是由主和祂的宗徒传承下来的——他们不能避免进一步的承认：婴儿需要中保的同样益处，为能藉着圣事和信徒的爱德被洗净，从而被接入基督的身体——即教会——与天主和好，并在祂内生活，得救、被释放、被救赎、被光照。但若说不是从死亡、罪恶的恶习、罪债、奴役和黑暗中，那又是从哪里呢？既然他们在幼小的年龄没有通过实际的过犯犯任何罪，那么只留下原罪。\par

% structure: p0081 source=h2 target=subsection reason=relative-heading-rank; base=h2

\subsection{第四十章 [第二十七]——圣经证言汇集。来自福音书。}

这段推理，在收集了我所应引用的大量圣经证词之后，会更有分量。我们已经引用了\BibleQuote{我来不是召义人，而是召罪人。}\BibleRef{路 5:32}同样，[主]进入\Person{匝凯}的家时说：\BibleQuote{今天救恩来到了这家，因为他也是\Person{亚巴郎}之子；因为人子来，是为寻找并拯救那失落的。}\BibleRef{路 19:9}同样的真理也表达在那迷失的羊和那九十九只羊直到失落的被寻回的比喻中；\BibleRef{路 15:4}以及那十个银币中失落一个的比喻中。\BibleRef{路 15:8}因此，正如祂所说：\BibleQuote{必须从\Place{耶路撒冷}开始，向万邦宣讲悔改和赦罪。}\BibleRef{路 24:46}同样，\Person{马尔谷}在他的\WorkTitle{马尔谷福音}结束时告诉我们，主说：\BibleQuote{你们往普天下去，向一切受造物宣讲福音。信而受洗的，必得救；不信的，必被定罪。}\BibleRef{谷 16:15}现在，谁能不知道，在婴孩的案例中，受洗就是相信，不受洗就是不相信？从\WorkTitle{若望福音}我们已经引用了一些段落。不过，我也要请你们注意以下内容：\Person{若翰洗者}论及基督说：\BibleQuote{请看天主的羔羊，请看除免世罪者！}\BibleRef{若 1:29}祂也论及自己说：\BibleQuote{我的羊听我的声音，我认识他们，他们跟随我；我赐给他们永生，他们永远不丧亡。}\BibleRef{若 10:27}既然婴孩只能藉着圣洗成为祂的羊，那么如果他们不受洗，就必须丧亡，因为他们将得不到祂赐予祂的羊的永生。所以祂在另一处说：\BibleQuote{我是道路、真理、生命；除非经过我，谁也不能到父那里去。}\BibleRef{若 14:6}\par

% structure: p0083 source=h2 target=subsection reason=relative-heading-rank; base=h2

\subsection{第四十一章——取自伯多禄前书}

请看宗徒们接受这道理后，是多么恳切地宣讲它！\Person{伯多禄}在他的\WorkTitle{伯多禄前书}中说：\BibleQuote{愿我们的主耶稣基督的天主和父受赞美！祂因自己的大仁慈，借耶稣基督由死者中复活，重生了我们，为获得那充满生命的希望，为你们保留着不朽、无玷、不衰的嗣业，这嗣业已为他们保存在天上。你们因信德，赖天主的德能保全，为获得那在末世要显示的救恩。}\BibleRef{伯前 1:3}稍后他又说：\BibleQuote{为使你们的信德得到称誉、光荣和尊敬，你们虽然从前不认识祂，却相信祂；虽然现在看不见祂，却仍相信祂；当你们看见祂时，将有无可言喻、充满光荣的喜乐，获得信德的成效，就是灵魂的救恩。}\BibleRef{伯前 1:7}在另一处他又说：\BibleQuote{你们却是特选的种族、王家的司祭、圣洁的国民、属于天主的民族，为叫你们宣扬那召你们出离黑暗进入奇妙光明者的德能。}\BibleRef{伯前 2:9}他又说：\BibleQuote{基督一次为罪而死，义者为不义者而死，为把我们领到天主面前。}\BibleRef{伯前 3:18}在提到\Person{诺厄}方舟中八人得救后，他接着说：\BibleQuote{与此预像相似的圣洗，现在拯救了你们。}\BibleRef{伯前 3:21}婴孩与这救恩和光明无缘，除非他们因着嗣养成为天主的子民，并且持守为不义者而死的义者基督，被祂领到天主面前，否则他们将留在丧亡和黑暗中。\par

% structure: p0085 source=h2 target=subsection reason=relative-heading-rank; base=h2

\subsection{第四十二章——取自若望一书}

此外，从\WorkTitle{若望一书}中我遇到下面的话，似乎对解决这个问题必不可少。他说：\BibleQuote{但我们若在光明中行走，如同祂在光明中，我们就彼此相通，祂圣子耶稣基督的血就会洗净我们的一切罪。}\BibleRef{若一 1:7}同样，在另一处他说：\BibleQuote{如果我们接受人的见证，天主的见证更大：因为天主的见证是基于祂为祂圣子所作的见证。信天主圣子的，在自己内就有这见证；不信天主的，就是以天主为撒谎者，因为他不信天主为祂圣子所作的见证。这见证就是：天主赐给了我们永生，而这永生是在祂圣子内。拥有圣子的，就有生命；没有天主圣子的，就没有生命。}\BibleRef{若一 5:9}看来，婴孩若没有圣子，就不仅没有天国，也没有生命；而圣子只有通过圣洗才能得到。所以他又说：\BibleQuote{天主子之所以显现，是为了毁灭魔鬼的作为。}\BibleRef{若一 3:8}因此，如果天主子不在婴孩内毁灭魔鬼的作为，婴孩就无份于天主子的显现。\par

% structure: p0087 source=h2 target=subsection reason=relative-heading-rank; base=h2

\subsection{第四十三章——取自罗马书}

现在请容我请大家注意宗徒保禄在这件事上的见证。从他那里引用的话自然可以更多，因为他写了更多的书信，并且由他特别热切地推荐天主的恩宠，以对抗那些以自己行为夸耀、不认识天主的正义、企图建立自己的正义，而不服从天主正义的人。\BibleRef{罗 10:3} 在他的罗马书中他写道：\BibleQuote{天主的正义加于所有相信的人，因为没有区别；因为众人都犯了罪，缺乏天主的荣耀，却因他的恩宠白白地成义，藉着在基督耶稣内的救赎；天主立他为因信他的血而做赎罪祭，为显示他的正义，因为他宽容了先前所犯的罪；为在今时显示他的正义，为使他为正义的，也使相信耶稣的人成义。}\BibleRef{罗 3:22} 然后在另一段他说：\BibleQuote{对于工作的人，报酬不算恩宠，而是债务。但对于不工作、只相信那成就不虔敬之人成义者的人，他的信德被算为正义。正如达味也描述那人不靠自己行为、天主归义于他的福气，他说：\InnerQuote{罪恶得赦免、过犯得遮盖的人，是有福的。主不归罪于他的人，是有福的。}}\BibleRef{罗 4:4} 然后不久他注意到：\BibleQuote{但那话写给他，不是单为他记下的，也是为我们这些将要得算为正义的人，就是我们相信那使我们的主耶稣基督从死者中复活的人；他被交出是为我们的过犯，复活是为我们的成义。}\BibleRef{罗 4:23} 又过了一点他写道：\BibleQuote{因为我们还软弱的时候，基督就在适当的时候为不虔敬的人死了。}\BibleRef{罗 5:6} 在另一段他说：\BibleQuote{我们知道律法是属神的；但我属血肉，被卖给了罪。因为我所做的，我不明白；我所愿意的，我并不做；我所恨的，我倒去做。我若做我所不愿意的，我就同意律法是善的。既然如此，现在不是我做的，而是在我里面的罪做的。我也知道，在我里头（就是在我肉体中）没有善住着；因为愿意行善的能力在我，但怎样行善，我找不着。我所愿意的善，我不去做；我所不愿意的恶，我倒去做。我若做我所不愿意的，就不是我做的，而是在我里面的罪做的。于是我发现一条律：当我愿意行善的时候，恶就在我里。因为按着内心的人，我喜爱天主的律；但我看见在我肢体中另有一条律，与我心中的律交战，把我掳去，叫我去顺从在我肢体中罪的律。我真是苦啊！谁能救我脱离这死亡的身体呢？靠我们的主耶稣基督，天主的恩宠！}\BibleRef{罗 7:14} 让那些能说的人说吧，众人不是生在这死亡的身体中，好叫他们能断言不需要天主藉着耶稣基督的恩宠来脱离这死亡的身体。因此他后来加上几节：\BibleQuote{律法因肉体软弱所不能行的，天主却做到了：他派遣了自己的儿子，带着罪恶肉身的形状，为着罪，在肉身上定了罪。}\BibleRef{罗 8:3} 让那些胆敢说的人说吧，如果我们不是生在有罪的肉身中，基督就必须生在有罪肉身的形状中。\par

% structure: p0089 source=h2 target=subsection reason=relative-heading-rank; base=h2

\subsection{第四十四章——摘自格林多前后书}

同样，他给格林多人说：\BibleQuote{我首先传给你们的，也是我所领受的，就是基督照圣经所说，为我们的罪死了。}\BibleRef{格前 15:3} 同样，在他的格林多后书中：\BibleQuote{基督的爱催迫我们，因为我们是这样判断：一人既替众人死了，那么众人就死了；基督替众人死，是为使活着的人不再为自己活，而为替他们死而复活的那位活。所以，从今以后，我们不再按肉身认识人了；虽然我们曾按肉身认识基督，但如今不再这样认识了。所以若有人在基督里，他就是新受造物；旧的已成过去，看，一切都成了新的。一切都是出于天主，他藉着基督使我们与他和好，并将和好的职分赐给了我们。这是什么意思？就是天主在基督里使世界与自己和好，不将他们的过犯归到他们身上，并且将这和好的道理托付了我们。所以我们是基督的使者，好像天主藉着我们劝你们；我们替基督请求你们：与天主和好吧！因为他使那无罪的，替我们成为罪，好叫我们在他内成为天主的正义。\BibleRef{格后 5:14}我们与天主同工，也劝你们不要白白接受天主的恩宠。（因为他说：在悦纳的时候，我应允了你；在拯救的日子，我帮助了你。看，现在正是悦纳的时候；看，现在正是拯救的日子。）}\BibleRef{格后 6:1} 如今，如果婴儿不被包括在这和好与救恩中，谁会为他们要求基督的圣洗呢？但如果他们被包括，那么他们就算在基督替他们死的人中；除非他赦免他们的罪、不归咎于他们，他们不可能藉着他和好并被拯救。\par

% structure: p0091 source=h2 target=subsection reason=relative-heading-rank; base=h2

\subsection{第四十五章——摘自迦拉达书}

同样，宗徒给迦拉达人写道：\Quote{愿恩宠与平安，由天主父和我们的主耶稣基督赐与你们，祂为我们罪人舍弃了自己，为救我们脱离这邪恶的世代。}\BibleRef{迦 1:3} 在另一处又对他们说：\Quote{法律是为过犯而添设的，直到那蒙应许的后裔来到；且由天使经中保之手而立。中保并非为一方的，但天主却是一位。那么法律与天主的恩许相背吗？绝对不可！因为如果颁赐的法律能赋予生命，那么义德就确实出自法律。但圣经把众人都圈在罪中，好使恩许藉着对耶稣基督的信德，赐给相信的人。}\BibleRef{迦 3:19}\par

% structure: p0093 source=h2 target=subsection reason=relative-heading-rank; base=h2

\subsection{第四十六章——摘自致厄弗所书}

他给厄弗所人写了同样含义的话：\Quote{你们从前因过犯和罪恶原是死的人；那时你们随从这世界的潮流，顺从空中权能的首领，即现今在悖逆之子身上运行的邪神。我们从前也都在他们中间，放纵肉身的私欲，随从肉体和心意所喜好的行事，且生来就是义怒之子，和别人一样。然而富于慈悲的天主，因祂爱我们的大爱，竟在我们因过犯而死的时候，使我们同基督一起生活——你们得救是因着恩宠。}\BibleRef{弗 2:1} 稍后他又说：\Quote{你们得救是因着恩宠，藉着信德；这并不是出于你们自己，而是天主的恩赐；不是出于行为，免得有人自夸。因为我们原是祂的化工，在基督耶稣内受造的，为行善事，就是天主所预备叫我们行的。}\BibleRef{弗 2:8} 又过了一会儿：\Quote{那时你们没有基督，与以色列社团隔绝，无分于恩许的盟约，在这世上没有望德，也没有天主。但如今在基督耶稣内，你们从前远离的人，藉着基督的血，成为亲近的了。因为祂是我们的和平，祂使双方合而为一，拆毁了中间隔断的墙壁，以自己的肉身废除了那规条中的诫命律法，为把两者在自己身上造成一个新人类，这样和平便成就了，且藉着十字架使双方在一个身体内与天主和好，因为祂在自己身上消灭了仇恨。祂来向你们远方的人传了和平，也向近处的人传了和平，因为藉着祂，我们双方在一个圣神内，得以进到父面前。}\BibleRef{弗 2:12} 然后在另一处他写道：\Quote{正如真理在耶稣内，你们应当脱去照从前生活方式的旧人，这旧人是因迷惑的情欲而败坏的；并在你们心思的灵里更新，穿上新人，这新人是按照天主的肖像，在真实的义德和圣洁中受造的。}\BibleRef{弗 4:21} 又说：\Quote{不要使天主的圣神忧郁，因为你们是在祂内受了印记，以待救赎之日的来临。}\BibleRef{弗 4:30}\par

% structure: p0095 source=h2 target=subsection reason=relative-heading-rank; base=h2

\subsection{第四十七章——摘自致哥罗森书}

他给哥罗森人写了这些话：\Quote{感谢父，祂使我们有资格在光明中分享圣徒的基业；祂拯救了我们脱离黑暗的权势，把我们迁移到祂爱子的国度里；在祂内我们获得了救赎，即罪过的赦免。}\BibleRef{哥 1:12} 又说：\Quote{你们在祂内得以完全，祂是一切掌权者和权势的元首；在祂内你们也受了不是人手所行的割损，而是基督的割损，脱去肉躯之身；你们在圣洗中与祂一同埋葬，也因着那使祂从死者中复活的天主的运作，藉着信德与祂一同复活。你们从前因过犯和未受割损的肉躯，原是死的人，但天主却使你们与基督一同生活，赦免了我们的一切过犯；涂抹了那反对我们、与我们为敌的规条文字，把它从中间撤去，钉在十字架上；祂解除了掌权者和权势的武装，公开羞辱它们，藉着十字架向它们夸胜。}\BibleRef{哥 2:10}\par

% structure: p0097 source=h2 target=subsection reason=relative-heading-rank; base=h2

\subsection{第四十八章——摘自致弟茂德书}

然后他对弟茂德说：\Quote{这话是确实的，值得完全接纳：基督耶稣来到世上，是为拯救罪人；在罪人中我是个罪魁。然而我蒙了怜悯，为的是使在我身上，基督耶稣显示祂全部的坚忍，给那些将来要相信祂而得永生的人作榜样。}\BibleRef{弟前 1:15} 他又说：\Quote{因为只有一位天主，也只有一个在天主与人之间的中保，就是那人基督耶稣；祂舍了自己作众人的赎价。}\BibleRef{弟前 2:5} 在给同一弟茂德的第二封信中，他说：\Quote{所以你不要以给我们的主作证为耻，也不要以我这为主被囚的人为耻；但要为福音同我一起受苦，靠着天主的能力；祂拯救了我们，并以圣召召叫了我们，不是按照我们的行为，而是按照祂自己的旨意和恩宠；这恩宠是在万世以前，在基督耶稣内赐给我们的，如今却藉着我们救主基督耶稣的显现，显示了出来；祂已消灭了死亡，并藉着福音将生命和不朽彰显了出来。}\BibleRef{弟后 1:8}\par

% structure: p0099 source=h2 target=subsection reason=relative-heading-rank; base=h2

\subsection{第四十九章——摘自致弟铎书}

然后他给弟铎写了以下的话：\Quote{期待那所祝福的望德，以及伟大的天主和我们的救主耶稣基督光荣的显现；祂为我们舍弃了自己，为救赎我们脱离一切罪恶，并为自己洁净一个子民，热心行善。}\BibleRef{铎 2:13} 在另一处又有同样含义的话：\Quote{但当天主我们的救主的良善和慈爱显现时，祂拯救了我们，不是因为我们所行的义德，而是按照祂的仁慈，藉着重生之洗和圣神的更新；这圣神是祂藉着我们的救主耶稣基督，丰丰富富地倾注在我们身上的，好使我们因祂的恩宠成义，成为永生的承继者。}\BibleRef{铎 3:3}\par

% structure: p0101 source=h2 target=subsection reason=relative-heading-rank; base=h2

\subsection{第五十章——摘自\BibleBook{}{希伯来书}}

虽然有些人对\BibleBook{}{希伯来书}的权威存疑，然而，我发现那些反对我们关于婴儿圣洗观点的人有时也认为该书支持他们的看法，因此我们注意到该书对我们立场的有力见证；我引用它时更加有信心，因为东方教会明确将其列入正典圣经。在该书的序言中，我们读到：\BibleQuote{天主在古时，曾多次并以多种方式，藉着先知向我们的祖先说过话；但在这末后的日子里，祂藉着自己的儿子向我们说话，天主立祂为万有的承继者，并藉着祂创造了宇宙。祂是天主光荣的反映，是天主本体的真像，以自己大能的话支撑万有。当祂藉着自己洗净了我们的罪之后，便坐在高天至尊者的右边。}\BibleRef{希 1:1}随后作者说：\BibleQuote{如果藉着天使所说的话成为确定的，而一切过犯和悖逆都得到了公正的报应，那么，我们若忽略了如此伟大的救恩，又怎能逃脱呢？}\BibleRef{希 2:2}又在另一处说：\BibleQuote{孩子既然都有同样的血肉，祂也照样亲自成了血肉，为能藉着死亡，毁灭那握有死亡权势的，就是魔鬼，并解救那些因怕死而一生处于奴役中的人。}\BibleRef{希 2:14}不久之后，祂又说：\BibleQuote{因此，祂应当在各方面相似弟兄们，好成为一位仁慈忠信的大司祭，在天主面前，为人民的罪作补赎。}\BibleRef{希 2:17}在另一处写道：\BibleQuote{我们要坚持所明认的信仰，因为我们所有的，不是一位不能同情我们弱点的大司祭，而是一位在各方面与我们相似，受过试探的，只是没有罪。}\BibleRef{希 4:14}又说：\BibleQuote{祂具有不可更改的司祭职。因此，凡藉着祂接近天主的人，祂都能拯救到底，因为祂常活着，为他们转求。这样的大司祭才适合我们，祂是圣善的、无辜的、无玷的、从罪人中分别出来的，且高于诸天；祂不必像那些大司祭一样，每日先为自己的罪，然后为人民的罪献祭，因为祂一次而永远地奉献了自己。}\BibleRef{希 7:24}又一次：\BibleQuote{因为基督并非进入了一个人手所造的圣所——那只是真圣所的模型——而是进入了上天本身，如今在天主面前为我们出现；祂也不是多次奉献自己，如同大司祭每年带着别人的血进入圣所一样；否则，从创世以来，祂就必须多次受苦了。可是如今，在世界的末期，祂显现了一次，以祂自己的祭献除掉了罪。就如规定人只死一次，之后便是审判；同样，基督也只一次奉献了自己，为承担大众的罪；将来祂要第二次显现，与罪无关，而是为那些期待祂的人带来救恩。}\BibleRef{希 9:24}\par

% structure: p0103 source=h2 target=subsection reason=relative-heading-rank; base=h2

\subsection{第五十一章——摘自\BibleBook{}{默示录}}

同样，若望默示录告诉我们，在一首新歌中，这些赞美献给基督：\BibleQuote{你配取那书卷，并开启它的封印，因为你曾被宰杀，并用你的血，从各支派、各语言、各民族、各邦国中，将我们救赎归于天主。}\BibleRef{默 5:9}\par

% structure: p0105 source=h2 target=subsection reason=relative-heading-rank; base=h2

\subsection{第五十二章——摘自\BibleBook{}{宗徒大事录}}

同样，在\BibleBook{}{宗徒大事录}中，宗徒伯多禄称主耶稣为\BibleQuote{生命的创始者}，斥责犹太人杀害祂，说：\BibleQuote{你们却否认了那圣洁且正义者，反而要求把凶手赐给你们；你们杀害了生命的创始者。}\BibleRef{宗 3:14}而在另一处他说：\BibleQuote{这石头是你们建筑工人所弃用的，却成了屋角的基石。除祂以外，别无救恩，因为在天下人间，没有赐下别的名字，使我们赖以得救。}\BibleRef{宗 4:11}又在别处说：\BibleQuote{我们祖先的天主复活了耶稣，你们却把祂悬在木架上杀害了。天主以右手高举祂，使祂作元首和救主，为赐给以色列悔改和赦罪。}\BibleRef{宗 5:30}又一次：\BibleQuote{众先知都为祂作证：凡信祂的人，因祂的名必获得罪的赦免。}\BibleRef{宗 10:43}而在同一宗徒大事录中，保禄说：\BibleQuote{所以，诸位弟兄，你们要知道：藉着这耶稣，罪的赦免已向你们宣讲了；凡信祂的人，都因祂而成义，从一切你们不能因梅瑟法律而成义的事上获得成义。}\BibleRef{宗 13:38}\par

% structure: p0107 source=h2 target=subsection reason=relative-heading-rank; base=h2

\subsection{第五十三章——旧约经书的益处}

在如此众多见证的重压之下，谁还敢胆敢开口反对天主的真理呢？还有许多其他见证可寻，只是因我急于结束这篇论述——我不可轻视这种急切——而作罢。我认为从旧约经书中再引用许多受默感的话语来证实我们的教义是多余的，因为那些隐藏在尘世许诺面纱之下的事，在新约的宣讲中已清晰地启示出来。我们的主亲自简要地说明并定义了旧约经书的用途，祂说，凡在律法、先知和圣咏中关于祂所写的，都必须应验，并且基督必须受苦，第三天从死者中复活，并且因祂的名，从耶路撒冷开始，向万民宣讲悔改和罪的赦免。这与我已经引用过的伯多禄的话一致，即众先知都为基督作证，凡信祂的人，都因祂而获得罪的赦免。\BibleRef{宗 10:43}\par

% structure: p0109 source=h2 target=subsection reason=relative-heading-rank; base=h2

\subsection{第五十四章——借旧约的祭献，人得以知罪并被引向救主。}

然而，或许从旧约中也提出一些证据更好，这些证据应当具有补充性，甚至说是累积性的价值。主亲自借着圣咏作者说：\Quote{至于我在世上的圣徒，祂使我的所有旨意在他们身上受到赞叹。}不是\Emph{他们的功绩}，而是\Quote{\Emph{我的旨意}}。因为除了随后提到的——\Quote{他们的软弱增加了}——超过他们已有的软弱之外，他们还有什么呢？此外，律法也介入，为使过犯增多。但为何圣咏作者立即加上：\Quote{他们急忙追赶？}当他们的忧愁和软弱增多时（即当他们过犯增多时），他们便更急切地寻求良医，以便罪在哪里增多，恩宠在哪里更加增多。然后祂说：\Quote{我不会收集他们的血祭集会（连同他们的祭献）；}因为借由他们多次的血祭，当他们起初在会幕、后来在圣殿聚集时，他们更被定罪为罪人，而非被洗净。祂说：我不再收集他们的血祭集会；因为已有一次为多人流出的血，使他们得以真正洗净。随后又说：\Quote{我也不会用我的嘴唇提他们的名字。}如同他们是重生之人的名字。因为起初他们的名字是：血肉之子、世俗之子、愤怒之子、魔鬼之子、不洁、罪人、不敬；但后来，他们成为天主的儿女——这是新人的新名字，是新事的歌者所唱的新歌，借着新约。人不可对天主的恩宠忘恩负义，对伟大的事物吝啬；当不断从微小上升到伟大。整个教会的呼声是：\Quote{我如同迷失的羊走失了。}从基督的所有肢体听到这声音：\Quote{我们都如羊走迷了路；祂亲自为我们的罪被交付。}\BibleRef{依 53:6}这段先知预言的全部就是依撒意亚中著名的段落，由斐理伯向甘达刻女王的太监讲解，而那人信了耶稣。\BibleRef{宗 8:30}请看祂多么频繁地称赞这个主题，并仿佛反复向骄傲好争辩的人灌输它：\Quote{祂是一个受苦的人，熟知承担软弱；因此祂转过脸，被轻慢，不被重视。祂承担我们的软弱，为我们陷入痛苦；我们以为祂在痛苦、在不幸、在受罚。但祂是为我们的过犯受伤，为我们的罪孽受创；我们平安的惩罚落在祂身上，因祂的创伤我们得医治。我们都如羊走迷了路，主将祂交付为我们的罪。祂虽受虐待，却不开口；如羊被牵到宰杀之地，又如羊羔在剪毛者前无声，祂不开口。在祂的屈辱中，祂的审判被夺去；谁能述说祂的世代？因为祂的生命从地被夺去，为我百姓的罪孽被引至死。所以我要将恶人赐给祂作坟墓，将富人赐给祂作死亡，因为祂没有行不义，口中没有诡诈。主喜悦洁净祂脱离不幸。你们若能为自己的罪献上灵魂，就必看见长寿的后裔。主喜悦从痛苦中拯救祂的灵魂，给祂显示光明，并以祂的理解塑造它；使那位善服务众人的义人成义；祂要亲自承担他们的罪孽。因此祂要承受许多人，并分享强者的掠物；祂被列在罪犯之中，祂亲自承担多人的罪，并为他们的罪孽被交付。}\BibleRef{依 53:3}也请考虑同一位先知的一段话，基督亲自宣布它应验在祂身上，当祂在会堂里履行读经职务时诵读它：\Quote{主的神临于我，因为祂傅油于我：派遣我向贫穷者传报喜讯，以致我能安慰所有伤心人——向俘虏宣讲释放，向盲者宣讲复明。}\BibleRef{依 61:1}那么，让我们都承认祂；在我们这样的人中，不应有任何人例外，我们都愿依附祂的身体，借着祂进入羊栈，并达到祂所应许给祂自己之人的生命和永生。——我重复，让我们都承认那位没有犯罪、在木头上以自己的肉身承担我们罪孽的那一位，为使我们可以与罪分离而活在正义中；凭祂的伤痕我们得医治，当我们如迷羊般软弱时。\par

% structure: p0111 source=h2 target=subsection reason=relative-heading-rank; base=h2

\subsection{第五十五章[第二十八.]——他总结说，所有人需要基督的死亡，才能得救。未受洗的婴儿将陷入魔鬼的定罪。所有人在亚当内如何趋于定罪；在基督内趋于成义。没有人能与天主和好，除非借着基督。}

在这样的情况下，那些借着洗礼来到基督面前的人，按照纯正的信德和真正的教义，从未有人被视为例外于罪过的赦免之恩；也没有人曾认为永生可能离开祂的国度。因为这永生已准备好要在末次时期启示出来，\BibleRef{伯前 1:5}即在死人复活之时，这些死人被保留，不是为那被称为\Quote{第二次死亡,}的永死，而是为那不能撒谎的天主应许给祂的圣人和忠信仆人的永生。凡分享这生命的，除非在基督内，无人得以复活，如同众人都在亚当内死去一样。\BibleRef{格前 15:22}因为凡属于按肉欲而生的人，没有一个不是在亚当内死去，因为众人都犯了罪；同样，凡由圣神意愿而重生的人，除非在基督内，无人被赋予生命，因为众人都成义。因为正如因一人众人都被判定，同样因一人众人都成义。\BibleRef{罗 5:18}也没有任何中立的地位给任何人，因此人若不在基督内，便只能与魔鬼在一起。因此，主自己（愿意从错信者心中除去那模糊不确定的中立状态——有人为未受洗的婴儿提供这种状态，仿佛因他们的天真，他们被收纳于永生，却因未受洗而不与基督同在祂的国内）说出了祂那终结性的判决，堵住了他们的口：\Quote{不与我相合的，就是敌我的。}\BibleRef{玛 12:30}那么，请随意举一个婴儿为例：如果他已在基督内，为何还受洗？然而，按真理的说法，他受洗正是为了与基督同在，那么可以肯定，未受洗的就不与基督同在；正因他不\Quote{与}基督同在，他就是\Quote{敌}基督的；因为主已宣告了祂自己的判决，这判决如此明确，我们不应、也确实不能削弱或改变它。他怎能\Quote{敌}基督，若不是因着罪？因为这不可能是来自他的灵魂或他的身体，两者都是天主的受造物。若是因着罪，在这样的年龄能有什么罪，除了那古老的原罪？当然，众人生于定罪中的罪恶肉身是一回事，那取了\Quote{罪恶肉身的相似体}的肉身是另一回事，众人也因此得以摆脱定罪。然而，这绝不意味着所有生于罪恶肉身的人实际上被那\Quote{相似}罪恶肉身的肉身所洁净；\Quote{因为众人没有信德；}\BibleRef{得后 3:2}而是所有生于肉欲结合的人都完全是生于罪恶肉身，而所有生于属神结合的人都只被那具有罪恶肉身相似体的肉身所洁净。换言之，前者在亚当内被定罪，后者在基督内成义。这好比我们说，在某城有一位接生婆，她为所有人接生；在同一地方有一位博学的教师，他教导所有人。\InnerQuote{所有人}，在前者中，只能理解为那些出生的人；在后者中，只能理解为那些受教的人：但并非所有出生的人也都接受教导。但每个人都清楚，在第一种情况中，正确地说\Quote{她为所有人接生,}因为没有她的帮助，无人出生；在第二种情况中，正确地说\Quote{他教导所有人,}因为没有他的指导，无人学习。\par

% structure: p0113 source=h2 target=subsection reason=relative-heading-rank; base=h2

\subsection{第五十六章——除基督外，无人与天主和好。}

考虑到我所引用的所有受启示的陈述——无论是逐段评估每一节的价值，还是将它们的联合见证累积起来，甚或包括我未引用的类似章节——所能发现的，无非是天主教会为虔诚守护一切亵渎的新奇事物所持有的：每个人都与天主分离，除非他们通过中保基督与天主和好；没有人能与天主分离，除非因着罪过，只有罪过造成分离；因此，除非通过最仁慈的救世主的唯一恩宠——通过最真实司祭的唯一祭献，罪过得赦免，否则没有和好；凡从信了蛇并因此被情欲败坏的女人所生的人，\BibleRef{创 3:6}除非通过相信天使并因此无欲而怀孕的童贞女的圣子，\BibleRef{路 1:38}否则不能脱离这死亡的身体。\par

% structure: p0115 source=h2 target=subsection reason=relative-heading-rank; base=h2

\subsection{第五十七章 [XXIX.]——婚姻之善；婚姻善用与恶用之四种不同情况。}

婚姻的善，不在于情欲的冲动，而在于以合法而可敬的方式使用这种冲动，以繁衍子女为目的，而非满足淫欲。因此，在这死亡的身体中，那悖逆地激动于肢体中，并试图将我们整个堕落的灵魂拖入自身（既不受心灵命令而兴起，也不受其令而平息）的，便是人人所生而具有的罪的恶。然而，当它被抑制于非法欲望之外，仅被允许用于有序繁衍和更新人类时，这便是婚姻的善，人由此在指定的结合中出生。然而，没有人在基督的身体中重生，除非他先前生在罪的身体中。正如善用恶事是善，恶用善事是恶一样，这两对概念——\Emph{善}与\Emph{恶}，以及\Emph{善用}与\Emph{恶用}——当适当结合时，便产生四种不同的情形：—— [1] 人善用善事，当他将自己的节德奉献给天主；[2] 人恶用善事，当他将自己的节德奉献给偶像；[3] 人恶用恶事，当他放纵私欲偏情而行奸淫；[4] 人善用恶事，当他以婚姻约束自己的私欲偏情。既然善用善事比善用恶事更好——因为两者都是善的——所以\Quote{那嫁女儿的，做得好；那不嫁女儿的，做得更好。}\BibleRef{格前 7:38} 事实上，关于这一问题，我曾按天主根据我卑微的能力所赐，在我两部著作中——一部论\WorkTitle{婚姻之善}，另一部论\WorkTitle{神圣童贞}——更详尽、更充分地予以讨论。因此，那些称赞有罪受造物的血肉，以贬低救赎主的血肉的人，不可藉婚姻之善为私欲偏情的恶辩护；那些主从他们的婴孩时代就教导我们谦卑功课的人，也不可因他人的错误而骄傲自满。\BibleRef{玛 18:4} 唯有那一位生而无罪，祂是由童贞女感孕，没有丈夫的拥抱——不是藉肉体的私欲偏情，而是藉她心灵贞洁的顺从。\BibleRef{路 1:34} 只有她能生育那治愈我们创伤的一位，她未受罪的创伤，便生出了纯洁后裔的胚芽。\par

% structure: p0117 source=h2 target=subsection reason=relative-heading-rank; base=h2

\subsection{第五十八章（XXX）——伯拉纠派如何认为为婴儿施洗是必要的}

现在，让我们更仔细地（按主赐予我们的能力）考察福音书中那非常的一章，其中祂说：\Quote{人若不重生——藉水和圣神——不能进入天主的国。}若不是因这经文的权威，他们根本不会认为婴儿应该受洗。这是他们对这段经文的注释：\Quote{因为祂没有说：\InnerQuote{人若不重生于水和圣神，便不能得到救恩或永生}，而仅仅说：\InnerQuote{他不能进入天主的国}，因此婴儿应当受洗，以便他们能在天主的国里与基督同在，除非受洗，他们便不能在那里。然而，即使婴儿未受洗而死，他们也将有救恩和永生，因为他们不受任何罪的锁链束缚。}这种说法首先引人注意的是，他们从未解释，将毫无罪过的\Quote{天主的肖像}与天主的国分离，正义何在。其次，我们应当看主耶稣——这唯一良善的导师——是否在这段福音经文中暗示并启示我们，受洗者唯有通过罪的赦免才能到达天主的国；尽管对于正确理解的人，经文中的话本身应当已经足够明确：\Quote{人若不重生，不能看见天主的国}；\BibleRef{若 3:3} 以及：\Quote{人若不藉水和圣神重生，不能进入天主的国。}\BibleRef{若 3:5} 因为为何他要重生，若不是为了更新？他从什么更新，若不是从某种旧的状态？从何种旧的状态，若不是从那\Quote{我们的旧人已与祂同钉十字架，使罪身被消灭}的状态？\BibleRef{罗 6:6} 或者，为何\Quote{天主的肖像}不能进入天主的国，若不是因罪的阻碍？然而，让我们（如前所述）尽力殷勤而勤奋地查看，这段福音经文——关于所讨论的问题——的全部上下文。\par

% structure: p0119 source=h2 target=subsection reason=relative-heading-rank; base=h2

\subsection{第五十九章——他们主要经文的上下文}

\BibleQuote{有一个法利塞人，}我们读到，\BibleQuote{名叫尼苛德摩，是个犹太人的首领。这人夜间来见耶稣，向祂说：\Quote{辣彼，我们知道祢是由天主而来的师傅，因为祢若不与天主同在，谁也不能行祢所行的这些神迹。}耶稣回答说：\Quote{我实实在在告诉你：人除非由上而生，不能见到天主的国。}尼苛德摩说：\Quote{人已年老，怎样能重生呢？难道他还能再入母胎而重生吗？}耶稣回答说：\Quote{我实实在在告诉你：人除非由水和圣神而生，不能进入天主的国。由肉生的属于肉，由神生的属于神。你不要惊奇，因我给你说了：你们应该由上而生。风随意向哪里吹，你听到风的响声，却不知道风从哪里来，往哪里去：凡由圣神生的，也是这样。}尼苛德摩问说：\Quote{这事怎能成就呢？}耶稣回答说：\Quote{你是以色列的师傅，连这事你都不知道吗？我实实在在告诉你：我们知道的，我们才讲论；我们见过的，我们才作证；而你们却不接受我们的见证。如果我给你说地上的事，你们尚且不信；如果给你说天上的事，你们怎么信呢？没有人上过天，除了那自天降下而仍在天上的人子。正如梅瑟曾在旷野里高举了蛇，\BibleRef{户 21:9}人子也应照样被举起来，使凡信祂的人，在祂内获得永生。}天主竟这样爱了世界，甚至赐下了自己的独生子，使凡信祂的人不至丧亡，反而获得永生。因为天主没有派遣子到世界上来审判世界，而是为叫世界借着祂而得救。那信从祂的，不受审判；那不信的，已受了审判，因为他没有信从天主独生子的名字。审判就在于此：光明来到了世界，世人却爱黑暗甚于光明，因为他们的行为是邪恶的。的确，凡作恶的人，都憎恶光明，也不来就光明，免得自己的行为暴露出来。然而履行真理的，却来就光明，为显示出他的行为是在天主内完成的。}\BibleRef{若 3:1}至此，主的言论完全涉及我们当前探究的主题；从这一点起，圣史转入了另一个话题。\par

% structure: p0121 source=h2 target=subsection reason=relative-heading-rank; base=h2

\subsection{第60章 [XXXI.] — \Person{基督}，头和身体；由于基督位格中两性的结合，祂既留在天上，又在地上行走；一个基督如何能升天；头和身体，一个基督。}

\Person{尼苛德摩}不明白所告诉他的一切时，他向主询问这些事如何可能。让我们看看主如何回答他的询问；因为，既然他屈尊回答\Quote{这些事怎能成就？}的问题，他实际上将告诉我们，从肉性生殖而来的人，如何能获得属灵的再生。在简短地指出那自以为是老师而优越于他人者的无知，并责备所有这样的人因不接受他对真理的见证而不信之后，他继续探询并感到惊奇：既然他告诉他们关于地上的事，他们不信，他们是否会相信天上的事。尽管如此，他仍继续这个主题，并对所问的\Quote{这些事如何可能？}给出一个他人应当相信——尽管这些人拒绝——的答案。他说：\BibleQuote{没有人升过天，除了那从天上降下来的人子，他是在天上的。}\BibleRef{若 3:13} 因此，他说，属灵的出生将这样来临——人将从属地的变为属天的；他们只能通过成为我的肢体来获得这一点；这样，那降下的可以上升，因为没有人上升除非他先降下。因此，所有必须被改变并提升的，都必须在与基督的结合中相遇，以便那降下的基督可以上升，将他的身体（即他的教会）视同他自己，因为关于基督和教会，这一点最真实地被理解：\BibleQuote{二人要成为一体。}\BibleRef{创 2:24} 关于同一件事，他亲自明确地说：\BibleQuote{如此，他们不再是两个，而是一体。}\BibleRef{谷 10:8} 因此，他们完全不能上升，因为\BibleQuote{没有人升过天，除了那从天上降下来的人子，他是在天上的。}\BibleRef{若 3:13} 因为尽管他在地上成为人子，但他并不认为那神性不配——他虽然留在天上，却降下到地上——用人子的名称来指称它，正如他用天主的名称来尊崇他的肉体：好让他们不被视为两个基督——一个是天主，另一个是人——而是同一个天主和人——天主，因为\BibleQuote{在起初已有圣言，圣言与天主同在，圣言就是天主；}\BibleRef{若 1:1} 而人，因为\Quote{圣言成了血肉，寄居在我们中间。} 借此——借他的神性与谦卑之间的差异——他作为天主子留在天上，而作为人子在地上行走；同时，借着他位格的合一，使他的两个性体成为同一个基督，他既作为天主子在地上行走，又同时作为同一个人子留在天上。因此，对更可信之事的信德来自对更不可信之事的相信。因为，如果他的神性，虽是一个更为遥远的对象，且在其不可比拟的差异中更为崇高，却有如此能力，为我们而取了人性，成为一个位格，并且，在地上以软弱的肉身显现为人子的同时，能够始终留在天上，处于分享了肉体的神性中，那么，我们的信德更容易设想：其他人，即他忠信的圣人们，与那人基督成为同一个基督，以致当所有人都借他的恩宠和共融上升时，那从天上降下的同一个基督自己升到天上。正是在这个意义上，宗徒说：\BibleQuote{就如身体是一个，却有许多肢体，而且肢体虽多，仍是一个身体，基督也是这样。}\BibleRef{格前 12:12} 他没有说：\Quote{基督的也是这样}——意即基督的身体或基督的肢体——而是说：\Quote{\Emph{基督也是这样}}，从而称头与身体为同一个基督。\par

% structure: p0123 source=h2 target=subsection reason=relative-heading-rank; base=h2

\subsection{在旷野中被举起的蛇预象被悬在十字架上的基督；连婴孩自己也被蛇咬所毒害。}

既然这伟大而奇妙的尊严只能通过罪的赦免获得，他继续说：\BibleQuote{正如梅瑟在旷野里举起了蛇，人子也必须照样被举起来，使凡信他的人不至丧亡，反而获得永生。}\BibleRef{若 3:14} 我们知道那时在旷野里发生了什么事。许多人因蛇咬而丧命；人民于是承认他们的罪，并通过梅瑟恳求主除去他们身上的毒物；因此，梅瑟遵主的命令，在旷野里举起了一条铜蛇，并告诫百姓，凡被蛇咬的，都应仰望那被举起的像。他们一这样做，就立刻痊愈了。\BibleRef{户 21:6} 那被举起的蛇除了表示基督的死，又是什么呢？这是一种表达记号的方式，即所成就的事由那成就它的事来表征。死亡来自蛇，它引诱人犯罪，使人因此该死。然而，主将他自己的肉身中转移的不是罪，如同蛇的毒，而是转移了死亡，好使无罪的刑罚在罪恶肉身的形状中出现，从而在罪恶肉身中，罪过和刑罚两者都被除去。因此，正如当时凡仰望那被举起的蛇的人，既从毒中痊愈，又脱离死亡，同样现在，凡借着信德和圣洗，与基督的死亡的形状相结合的人，既由成义而从罪中得释放，又由复活而从死亡中得释放。因为这就是他所说的：\BibleQuote{使凡信他的人不至丧亡，反而获得永生。}\BibleRef{若 3:15} 那么，如果一个婴孩完全没有被蛇的咬伤所毒害，又有什么必要借圣洗与基督的死亡相结合呢？\par

% structure: p0125 source=h2 target=subsection reason=relative-heading-rank; base=h2

\subsection{除基督外，无人能与天主和好。}

他接着这样说：\BibleQuote{天主如此爱了世界，竟赐下祂的独生子，使凡信祂的人不致丧亡，反而获得永生。}\BibleRef{若 3:16}因此，每一个婴儿，若不藉着圣洗圣事信天主的独生子，就注定要丧亡并失去永生；然而，祂来不是为审判世界，而是为叫世界藉着祂得救。这在下一句中尤其明显，祂说：\BibleQuote{信祂的人不被定罪；不信的人已被定罪，因为他没有信天主的独生子的名字。}\BibleRef{若 3:18}那么，我们将受过洗的婴儿归入哪一类呢？岂不是如天主教会权威到处宣称的，归入信徒之列？因此，他们属于那些相信的人；因为这是藉着圣事的功效和代父母的答复为他们获得的。由此推出，未受洗的婴儿被算作不信的人。既然受洗的人不被定罪，那么这些未受洗的人就被定罪。祂接着又说：\BibleQuote{但定罪是此：光明来到了世界，世人却爱黑暗甚于光明。}\BibleRef{若 3:19}祂说\BibleQuote{光明来到了世界}，难道不是指祂自己的来临吗？而没有祂来临的圣事，婴儿怎能说在光明中？我们为何不也将以下事实包含在\BibleQuote{世人爱黑暗}中，即他们自己不信，也不愿认为他们的婴儿应当受洗，尽管他们害怕婴儿遭受身体的死亡？然而，祂宣称，\BibleQuote{那来到光明中的，是天主内所做的工作}\BibleRef{若 3:21}，因为他深知自己的成义不是来自任何功绩，而是来自天主的恩宠。因为宗徒说：\BibleQuote{是天主在你们内运行，使你们愿意并实行祂的美意。}\BibleRef{斐 2:13}这就是所有从肉身生育来到基督的人，在灵性上重生所成就的方式。祂亲自解释了这一点，并指明了这一点，当祂被问及这些事怎能成就时；祂不容许任何人凭人的推理来解决这样的问题，免得婴儿被剥夺赦罪的恩宠。没有别的道路通往基督；没有人能藉基督以外的任何方式与天主和好或来到天主面前。\par

% structure: p0127 source=h2 target=subsection reason=relative-heading-rank; base=h2

\subsection{第63章 [XXXIV.]——圣洗的形式或仪式。驱魔礼。}

关于这圣事的实际形式，我该说什么呢？我只希望持相反观点的某个人带一个婴儿来给我施洗。我的驱魔在那个婴孩身上起什么作用，如果他不在魔鬼的权下？带婴儿来的人当然必须充当他的代父母，因为他自己无法回答。那么，他怎么可能声明他弃绝魔鬼，如果婴孩内没有魔鬼？他怎么可能说皈依天主，如果他从未背弃天主？他怎么可能相信（除了其他信条外）罪的赦免，如果他没有罪？就我而言，如果我认为他的意见与这信仰相悖，我不会允许他带婴儿来领受圣事。我也无法想象他在人前以何种面目，或在天主前以何种心思，来处理这件事。但我不想说得过于严厉。他们中有些人认为，给婴儿施行一种虚假或骗人的圣洗形式，其中可能有某种行动的声音和表象，但实际上并没有罪的赦免发生，这已是所能提及或想象的最可憎可恶之事。同样，关于婴儿领受圣洗的必要性，他们承认婴儿也需要救赎——这是他们中某人写的一篇短论中承认的——但这篇著作中并没有公开承认任何单个罪的赦免。然而，根据你信中给我的暗示，如你所说，他们现在承认，婴儿藉着圣洗也获得了罪的赦免。这并不奇怪；因为\Emph{救赎}不可能以其他方式被理解。他们自己的话是：\Quote{然而，他们并非从起初，而是在他们出生后的实际生活中，才开始有罪的。}\par

% structure: p0129 source=h2 target=subsection reason=relative-heading-rank; base=h2

\subsection{第64章——关于婴儿的双重错误。}

你看到了，我在本书中如此漫长而执着地反对的那些人之间有多么巨大的差异——其中一人写了那本书，其中包含了我尽己所能驳斥的观点。如我所说，你看到他们之间的重要差异：有些人认为婴儿是绝对纯洁的，没有任何罪，无论是原罪还是本罪；而另一些人则认为，婴儿一出生就染上了自己实际犯下的本罪，因此需要透过圣洗圣事来洁净。后一类人确实通过查考圣经、考虑整个教会的权威以及圣事本身的样式，清楚地看到婴儿藉着圣洗圣事获得了罪过的赦免；但他们要么不愿意，要么不能承认婴儿所有的罪是原罪。然而，前一类人清楚地看到（这很容易），在所有人都能认识到的人的本性中，我们所说的这个幼小年龄在自己的行为中根本不可能犯任何罪；但为了避免承认原罪，他们断言婴儿根本没有罪。现在，在他们各自坚持的真理中，恰好发生的是：他们首先在某些方面彼此一致，随后又在重要方面与我们不同。因为如果一方让步于另一方，说所有受洗的婴儿都获得罪过的赦免，而另一方让步于他们的对手，说婴儿（正如婴儿本性本身无声地大声宣告）在自己的生活中尚未犯下任何罪，那么双方都必须同意让步于我们：除了原罪之外，没有什么可被圣洗圣事赦免给婴儿的了。\par

% structure: p0131 source=h2 target=subsection reason=relative-heading-rank; base=h2

\subsection{第六十五章 [三十五]——婴儿没有自己犯的罪}

这也会被质疑，我们必须花时间讨论它，以证明和说明：婴儿凭自己的意志——没有意志，在自己的生活中就不可能有罪——永远不可能犯罪；所有人正是因此习惯称他们为\Emph{无辜的？}难道他们心智和身体的极大软弱、对事物的极大无知、完全不能遵守诫命、缺乏对自然或成文法律的任何感知和印象、完全缺乏推动他们向任一方向的理性——这一切不正以一种远比我们任何论证更有说服力的无声见证，宣告并证明着我们面前的观点吗？事实的显而易见性必定大大有助于说服我们相信它的真实性；因为在我所说的事实中，没有一处找不到其踪迹，我们所谈论的事实如此普遍——它比我们所能说的任何证明更清楚易见。\par

% structure: p0133 source=h2 target=subsection reason=relative-heading-rank; base=h2

\subsection{第六十六章——婴儿的过失源于纯粹的无知}

然而，我希望任何在这点上明智的人告诉我，他在新生儿身上看到或想到了什么罪，以致于他认为圣洗圣事已经是必要的救赎；这个新生儿在自己的生活中，凭自己的心智或身体犯下了什么样的恶行。如果它碰巧哭闹，使长辈厌烦，我不知道我的告密者会将此归咎于邪恶，还是不幸。再者，他会对以下事实说什么：它被安抚而不哭，不是凭它自己的理性呼吁，也不是凭任何人的禁止？这来自它深陷其中的无知，由于这种无知，当它稍大一点（很快就如此），它会在小脾气中打母亲，甚至常常打它饥饿时吸吮的乳房。现在，这些小脾气在非常年幼的孩子身上不仅被忍受，而且实际上被喜爱——这是出于什么样的情感，除了肉体的情感？凭此，我们因聪明人的玩笑和笑话而开心，这些玩笑和笑话充满趣味和胡说，但如果按字面理解，它们就不会被当作诙谐而笑，而是被当作傻瓜而笑。我们也看到，那些被普通人称为\OriginalTerm{傻子}{Moriones}的傻瓜是如何被用作健康人的娱乐；而且在奴隶市场上，他们的售价高于健康人。所以，纯粹的自然情感对即使是绝非傻瓜的人也有如此大的影响，以致于从别人的不幸中取乐。然而，虽然一个人可能因另一个人的愚蠢而发笑，但他仍然不愿意自己成为傻瓜；如果父亲很高兴地寻找甚至引逗自己咿呀学语的孩子做出这些事，但预先知道孩子长大后会成为傻瓜，他无疑会认为孩子比死了更可悲。然而，只要还有成长的希望，智力的光芒预期会随着年岁的增长而增长，那么幼童对父母的冒犯不仅不被视为错误，反而被视为愉快和适意的。毫无疑问，明智的人不可能赞成在孩子能够被禁止时，不仅不禁止他们这样的言行，甚至为了长辈的虚荣取乐而鼓励他们。因为一旦孩子到了认识父母的年龄，他们就不敢对父母说错话，除非得到父母一方或双方的允许或命令。但这样的事情只能属于那些刚刚努力咿呀学语、心智刚刚能够以某种方式运动舌头的幼童。然而，让我们考虑一下新生婴儿的无知之深，随着他们年龄的增长，他们从这种无知中走出来，进入这种暂时的口齿不清的愚蠢——可以说是走在通往知识和言语的路上。\par

% structure: p0135 source=h2 target=subsection reason=relative-heading-rank; base=h2

\subsection{第六十七章 [三十六]——论婴儿的无知及其根源}

是的，让我们思索那他们\OriginalTerm{理性理智}{rational intellect}的\OriginalTerm{黑暗}{darkness}，由于这黑暗，他们甚至完全对天主无知，而他们实际上在受洗时对抗天主的圣事。现在我问：这黑暗是在何时、从何处使他们沉浸其中的？难道他们是在\Emph{这里}、在此生中——在母胎中经历了智慧和虔诚的生活之后——因过于疏忽而忘记了天主吗？让那些敢于这样说的人说吧；让那些愿意听的人听吧；让那些能够相信的人信吧。然而，我确信，任何没有被顽固坚持先入之见所蒙蔽的人，都不可能持有这样的意见。难道无知中没有邪恶——没有需要被炼净的东西吗？那祈祷说：\Quote{求你不要记念我幼年时的罪愆和我的无知？}是什么意思？因为虽然明知故犯的罪更应受谴责，但如果没有无知的罪，我们就不应读到圣经中我所引用的：\Quote{求你不要记念我幼年时的罪愆和我的无知。}现在，既然刚从母胎中出来的婴儿的灵魂仍然是人的灵魂——甚至理性受造物的灵魂——不仅未受教导，而且甚至不能接受教导，我问：它是为何、在何时、从何处陷入那厚重无知的黑暗中的？如果这是人的本性如此开始，而那本性尚未败坏，那么为什么亚当不是这样被造的？为什么他能够领受诫命？并且能够给他的妻子和所有动物命名？因为他论及她说：\BibleQuote{她应称为女人}\BibleRef{创 2:23}；关于其余的，我们读到：\BibleQuote{凡亚当给每种生物起的名字，就是那生物的名字。}\BibleRef{创 2:19}而这人，虽然他不知道自己在哪里、是什么、由谁创造、由何父母所生，却已有了罪的过犯，尚不能领受诫命，并且完全被无知的厚云所笼罩和淹没，以至于他无法从沉睡中醒来，意识到这些事实；而必须耐心等待时机，直到他能像摆脱某种奇怪的醉酒一样（确实，不是像最沉重的醉酒通常能在一夜间摆脱，而是）一点一点地，经过许多个月甚至许多年；而在此完成之前，我们在小孩子身上必须忍受许多我们在成年人身上惩罚的事情，多到我们无法枚举。那么，关于这无知和软弱的巨大邪恶，如果婴儿在此生一出生就染上了它，那么他们在何时、何处、如何通过犯下某个大罪而突然陷入这样的黑暗？\par

% structure: p0137 source=h2 target=subsection reason=relative-heading-rank; base=h2

\subsection{第六十八章【第三十七章】——如果亚当并非被造成像我们出生时的那种本性，那么基督虽然无罪，为何以婴儿和软弱之身诞生？}

有人会问：\Quote{如果这本性并不纯正，而是从起源就败坏了，因为亚当并非这样被造，那么基督——远为卓越，且由童贞女所生而无罪——为何也以这种软弱出现，并以婴儿之身来到世界？}对此问题的回答如下：亚当并非处在这样的状态中，因为在他之前没有来自父母的罪，所以他并非在\OriginalTerm{罪肉}{sinful flesh}中被造。然而我们却处于这样的境况，因为由于他先前的罪，我们在罪肉中出生。而基督却是在这样的状态中出生，因为\BibleQuote{祂为了在肉身上定罪恶，就取了罪肉的相似体}\BibleRef{罗 8:3}。我们现在讨论的问题并非关于亚当身体大小的方面——为何他不是以婴儿之身被造，而是以四肢完美成熟的形态。确实可以说，野兽也是这样被造的——并非因为它们犯罪，它们的幼崽才生来细小。关于这一切为何发生，我们现在不追问。但我们关心的是人智力和理性方面的活力，凭借这些，亚当能够接受教导，能够领会天主的诫命和律法，并且如果愿意，就能轻易遵守；而人现在出生时却完全不能做到，由于可怕的无知和软弱，不是身体上的，而是精神上的——尽管我们必须承认，在每个婴儿中都存在一个\OriginalTerm{理性灵魂}{rational soul}，其本质与第一个人相同。然而，这巨大的肉体软弱，在我看来清楚表明某种具有惩罚性质的东西。这引起疑问：如果最初的人类没有犯罪，他们的孩子是否会既不能使用舌头、也不能使用手和脚？他们出生时是孩子或许是必要的，因为子宫容量有限。但同时，不能因为肋骨是人身体的一小部分，就认为天主为那人创造了一个婴儿妻子，然后把她建造成一个女人。同样，天主的全能也有能力使她的孩子一出生就立刻长大成人。\par

% structure: p0139 source=h2 target=subsection reason=relative-heading-rank; base=h2

\subsection{第六十九章【第三十八章】——婴儿的无知与软弱}

但不必详细讨论这一点，至少有许多动物确实可能如此：它们的幼崽出生时很小，随着身体长大，心智并不进步（因为它们没有\OriginalTerm{理性灵魂}{rational soul}），然而，即使最幼小时，它们也能奔跑、认出母亲，想要吮吸时无需外界帮助或照料，而是以非凡的轻松自己找到母亲的乳房，尽管这些乳房寻常视线难以发现。相反，人在出生时既无适合行走的脚，也无能抓挠的手；除非母亲将他们的嘴唇实际贴到乳房上，他们不知道在哪里找到它；即使靠近乳头，尽管渴望食物，他们也会更会哭而不是吮吸。身体的这种完全无助与其心智的软弱相吻合；基督的肉身也不会成为\Quote{在\OriginalTerm{罪恶肉身}{sinful flesh}的模样上}，除非那罪恶肉身在如我们所述的方式中压迫着理性灵魂——无论这灵魂是由父母遗传而来，还是为每个个体分别创造，或是从上天赋予——关于这点我现在不予探究。\par

% structure: p0141 source=h2 target=subsection reason=relative-heading-rank; base=h2

\subsection{第七十章 ［第三十九］——论婴儿如何藉圣洗除去罪，以及成年人如何，并由此产生什么益处。}

在婴儿，确然，藉天主的恩宠，通过那曾在\OriginalTerm{罪恶肉身}{sinful flesh}的模样上来到的\Emph{祂的}圣洗，得以成就罪恶肉身被除去。然而，这一结果如此实现，以至于遍布于活肉身并生来具有的\OriginalTerm{私欲偏情}{concupiscence}并非一下子被移除，不再存在于其中；而只是使那在出生时固有的东西不在人死时有害。因为若一个婴儿在领洗后继续活着，到达能够服从法律的年龄，他会在那里发现某种需要与之战斗的东西，并藉天主的助佑去克服它，只要他没有白白领受祂的恩宠，并且不愿成为可弃绝的人。因为即使在年岁较长的人，在圣洗中也不被赐予（除非也许通过全能创造者不可言喻的奇迹）使他们肢体中那\Emph{罪恶的法律}（与他们的心神的法律交战）完全熄灭并不再存在；而是使一个人曾作为隶属于其私欲偏情的仆人的时候所\Emph{做过、说过、想过的一切邪恶}都被消除，被视为从未发生过。然而，这私欲偏情本身（尽管罪恶的束缚——魔鬼曾藉它来束缚灵魂——被解除，以及将人与造物主隔开的障碍被摧毁）仍然留在我们克制身体、使其服从的争战中，无论是为了合法且必要的用途而放松，还是藉节制加以约束。\BibleRef{格前 9:27} 但既然认识人类所有过去、现在和未来的圣神远胜于我们，祂预见了并预言了人的生命将是这样：\Quote{没有一个活着的人能在天主面前成义}，于是因无知或软弱，我们没有用意志的全部力量去对抗它，反而在犯各种不法之事中屈服于它——我们变得越坏，取决于我们投降的程度和频率；我们变得越好，取决于它不重要和稀少。然而，对我们现在感兴趣之点的探究——是否可能（或事实上是否存在、曾经存在或将要存在）一个在今世没有罪的人，除了那说过\BibleQuote{这世界的首领来了，他在我内一无所获}\BibleRef{若 14:30}的那一位——需要更充分的讨论；而本论著的编排使我们把这个疑问推迟到另一卷书的开始。\par

\SourceRecord{Translated by Peter Holmes and Robert Ernest Wallis, and revised by Benjamin B. Warfield.；Nicene and Post-Nicene Fathers, First Series；Vol. 5.；Edited by Philip Schaff.；Buffalo, NY: Christian Literature Publishing Co.,；1887.；<http://www.newadvent.org/fathers/15011.htm>.}

% structure: p0001 source=h1 target=section reason=page-title

\section{\WorkTitle{论功过与罪的赦免，及婴儿的圣洗（卷二）}}

\Emph{奥斯定在此反驳那些声称在今生有、曾有、且将有完全无罪之人的观点。他就此提出四项主张：首先，人藉天主的恩宠和自己的自由意志，可能在今生不犯罪；其次，他说明事实上没有人在今生完全无罪；第三，他陈述其原因——因为没有人能将自己的意愿完全限于每个情况之合理要求，而该合理要求若非未能察觉，便是无意实行；第四，他证明除了唯一的中保基督以外，没有、也不曾有、且永远不会有任何人是完全无罪的。}\par

% structure: p0003 source=h2 target=subsection reason=relative-heading-rank; base=h2

\subsection{第一章 [I]——以上所论及本书将探讨的内容。}

最亲爱的马尔切利诺，我想在前一卷书中已充分讨论了婴儿的圣洗——圣洗不仅使他们得以进入天主的国，也为获得救恩与永生，而永生不能无天主的国，也不能无与救主基督的结合，祂以自己的血救赎了我们。我今在本书中，将尽力探讨和解释一个问题：在这世界上，是否有人活着、或曾有人活过、或将来会有人活过，而没有任何罪，除了\Quote{天主与人之间的唯一中保，降生成人的基督耶稣，祂舍己作万人的赎价}？\BibleRef{弟前 2:5}——并仰赖祂自己赐予我的恩宠与能力。若在此讨论中，或因论证之必然或因偶然，涉及关于婴儿的圣洗或罪的任何问题，我既不惊讶，也不应畏缩，而在这些关键时刻，尽我所能回答任何引起我注意的问题。\par

% structure: p0005 source=h2 target=subsection reason=relative-heading-rank; base=h2

\subsection{第二章 [II]——有些人过于看重人的自由意志；无知与软弱。}

由于我们每日的祈祷，极为必要解决关于一个未受任何罪之欺骗或侵扰的人生的这个问题。因为有些人如此信任人的自由决断，以致认为它无需犯罪，也无需天主的助佑——将自由决断归因于我们本性的定数。其必然结果便是，我们不应祈祷\Quote{不要让我们陷于诱惑}——即不要被诱惑所胜，无论它在我们的\Emph{无知}中欺骗和突袭我们，还是在我们的\Emph{软弱}中压迫和纠缠我们。然而，若不恳求主赐予这样的恩惠，反而认为主祷文中\Quote{不要让我们陷于诱惑}\BibleRef{玛 6:13}这一祈求是虚妄无用的，这有多有害、多邪恶、多违背我们在基督内的救恩，又多暴烈地抵触我们所受教导的宗教本身以及我们崇拜天主的虔敬，我实在无法用言语表达。\par

% structure: p0007 source=h2 target=subsection reason=relative-heading-rank; base=h2

\subsection{第三章 [III]——天主如何不命令不可能的事；慈悲的工作是消除罪过的方法。}

这些人自以为敏锐（仿佛我们中无人知晓），当他们说：\Quote{若我们没有意愿，就不犯罪；天主也不会命令人去做人的意志不可能的事。}但他们没有看到，为了克服某些事物——无论是邪恶欲念或不当恐惧的对象——人需要意志的奋力奋斗，有时甚至需要全部精力；而那预见这一切的主，愿意先知说出如此真实的话：\Quote{在祢面前，凡活着的人没有一个得以成义。}因此，主预见我们会有这样的性格，乐意提供并赋予某些有益健康的补救方法以有效之能，来对抗甚至在圣洗后所犯的罪过的罪债和束缚——例如，慈悲的工作——正如祂所说：\Quote{你们宽恕，也必蒙宽恕；你们给予，也必蒙给予。}\BibleRef{路 6:37}因为，若没有人临终时能带着永生盼望，面对\Quote{凡遵守全部法律，而只在一条上失足的，就成了全律法的罪人}\BibleRef{雅 2:10}这句判决，除非紧接着有：\Quote{你们既然要按自由的法律受审判，就该照这法律说话行事；因为不行慈悲的人，受审判时没有慈悲；慈悲却胜过审判。}\BibleRef{雅 2:12}\par

% structure: p0009 source=h2 target=subsection reason=relative-heading-rank; base=h2

\subsection{第四章 [IV]——\OriginalTerm{私欲偏情}{concupiscentia}在我们内如何存在；已受洗者不受私欲偏情所害，只有同意它才受害。}

因此，身为罪的律的\OriginalTerm{私欲偏情}{concupiscence}存留于这死亡之身体的肢体中，与婴儿一同出生。在已受洗的婴儿身上，它被剥夺了罪责，留作（生命的）战场，但不追索那些在战斗前已死之人的定罪。未受洗的婴儿，它则将其定为有罪，作为义怒之子，即便他们死于襁褓，也引向定罪。然而在已受洗、拥有理智的成年人身上，他们的心神对这个私欲偏情为犯罪所给予的任何同意，都是他们自身意志的行为。在所有的罪都被抹去，那按本性将人束缚于被征服状态中的罪责被取消之后，它仍然存留——但丝毫不伤害那些不向它屈服以行非法之事的人——直到死亡被\BibleRef{格前 15:54}的胜利所吞没，并在那完美的平安中，再无可征服之物。然而，那些向它屈服以行非法之事的人，它则视其为有罪；除非通过悔改的药方和慈悲的工作，并藉着天上大司祭为我们所作的转求而获得医治，否则它便将我们引向第二次死亡和彻底的定罪。正是为此，主在教我们祈祷时，除其他祈求外，还劝告我们说：\Quote{宽免我们的罪债，如同我们也宽免得罪我们的人；不要让我们陷入诱惑，但救我们免于凶恶。}\BibleRef{玛 6:12}因为邪恶存留在我们的肉体内，并非由于天主以智慧创造人的本性，而是由于那由人自己的意志而陷入的过犯，在此过犯中，由于能力丧失，人并非以受伤时同样的意志力便能得医治。关于这邪恶，宗徒说：\Quote{我知道，在我的肉体内没有善好的东西存留；}\BibleRef{罗 7:18}同样，针对这同一邪恶，他劝诫我们不要顺从，当他说：\Quote{所以，不要让罪在你们必死的身体内为王，以顺从它的情欲。}\BibleRef{罗 6:12}因此，当我们因意志的非法倾向而向这些肉体的情欲屈服同意时，我们便为治疗这一过犯而说：\Quote{宽免我们的罪债，}\BibleRef{玛 6:12}同时，我们应用慈悲工作的疗法，加上：\Quote{如同我们也宽免得罪我们的人。}然而，为避免我们如此屈服同意，让我们祈求助佑，并说：\Quote{不要让我们陷入诱惑；}——并非因为天主亲自以这样的诱惑试探任何人，\Quote{因为天主不是邪恶的试探者，祂也不试探任何人；}\BibleRef{雅 1:13}而是为了使我们每当感到从我们的私欲偏情升起的诱惑时，不被祂的帮助所遗弃，以致我们能够藉此得胜，而不被诱惑所带走。然后，我们加上祈求那在末时才能完成的事，那时必死的将被生命所吞没：\BibleRef{格后 5:4}\Quote{但救我们免于凶恶。}\BibleRef{玛 6:13}因为那时将不再有那被命令与之战斗且不得同意的私欲偏情。因此，这三个祈求的全部内容可以简要表达如此：\Quote{赦免我们因被私欲偏情牵引而犯的过犯；帮助我们不被私欲偏情牵引；将私欲偏情从我们身上除去。}\par

% structure: p0011 source=h2 target=subsection reason=relative-heading-rank; base=h2

\subsection{第五章——人的意志需要天主的帮助。}

因为犯罪的施行，我们从天主得不到任何帮助；但若没有天主的帮助，我们就不能行正义，并完全履行义的法律。因为如同肉体的眼睛，如果闭上或转开，并不从光得到帮助，但光帮助它观看，并且除非光帮助它，它根本看不见；同样，天主，作为内在人的光，帮助我们的心智之眼，使我们能行一些善事，不是按照我们自己的义，而是按照祂的义。但我们若转离祂，那是我们自己的行为；那时我们是按肉性而智慧，那时我们便向肉体的\OriginalTerm{私欲偏情}{concupiscence}为非法之事而屈服同意。因此，当我们转向祂时，天主帮助我们；当我们转离祂时，祂舍弃我们。但祂甚至帮助我们转向祂；而这确实是肉体之眼的光所不做的事。因此，当祂藉着这些话命令我们：\Quote{转向我，我必转向你们，}\BibleRef{匝 1:3}而我们却对祂说：\Quote{天主，我们的救主，求祢转向我们，}又说：\Quote{万军的天主，求祢转向我们；}此外我们还说什么呢？岂不是说：\Quote{把祢所命令的赐下吧？}当祂命令我们，说：\Quote{如今，你们愚笨的，要明白，}而我们却对祂说：\Quote{求祢赐我明悟，好让我学习祢的诫命；}此外我们还说什么呢？岂不是说：\Quote{把祢所命令的赐下吧？}当祂命令我们，说：\Quote{不要随从你的情欲，}\BibleRef{德 18:30}而我们却对祂说：\Quote{我们知道，没有人能\OriginalTerm{节制}{continent}，除非天主赐给他；}\BibleRef{智 8:21}此外我们还说什么呢？岂不是说：\Quote{把祢所命令的赐下吧？}当祂命令我们，说：\Quote{你们要施行正义，}\BibleRef{依 56:1}而我们却说：\Quote{上主，求祢教训我祢的判断；}此外我们还说什么呢？岂不是说：\Quote{把祢所命令的赐下吧？}同样，当祂说：\Quote{饥渴慕义的人是有福的，因为他们要得饱饫，}\BibleRef{玛 5:6}我们应从谁那里寻求义的食物和饮料呢？岂不正是从那应许饥渴慕义者以祂的饱足的那一位？\par

% structure: p0013 source=h2 target=subsection reason=relative-heading-rank; base=h2

\subsection{第六章——法利塞人感谢天主时在何处犯罪；天主的恩宠上须加上我们自身意志的努力。}

让我们于是从我们的耳中和心中驱赶那些说我们应当接受我们自己自由意志的决定，而不必祈求天主帮助我们不犯罪的人。法利塞人甚至不曾被这样的黑暗所蒙蔽；因为虽然他错误地认为他不需要为自己的义德添加什么，并以为自己充满了义德，他仍然感谢天主，因为他\Quote{不像其他人，不义、勒索、奸淫，或者连这个税吏也不如；他每周禁食两次，凡他所有的都献上十分之一。}\BibleRef{路 18:11}他确实不想要为自己的义德添加什么；但是，借着感谢天主，他承认自己所有的一切都是从他那里领受的。尽管如此，他并不被认可，既因为他没有祈求义德的进一步食粮，好像已经饱足了一样，又因为他傲慢地将自己置于那饥渴慕义的税吏之上。那么，对那些虽然承认自己没有义德，或没有义德的圆满，却仍以为义德只能从他们自己而来，不必向创造主祈求——在他那里才是义德的仓库和源泉——的人，又该说什么呢？然而，这并不仅仅是一个关于祈祷的问题，好像我们意志的活力也不应被努力加上似的。据说天主是\Quote{ \Emph{我们的助佑}; }；但没有人能在不自愿做出一些努力的情况下得到\Emph{帮助}。因为天主在我们内成就我们的救恩，并不是如同他在无感觉的石头内，或在自然既未赋予理性也未赋予意志的受造物内工作那样。然而，他为什么帮助这个人，而不帮助那个人；或者为什么帮助这个人这么多，那个人那么多；或者为什么帮助一个人用一种方式，帮助另一个人用另一种方式——他按照他自己最隐秘的公义的方式和他的权能的卓越，保留给自己。\par

% structure: p0015 source=h2 target=subsection reason=relative-heading-rank; base=h2

\subsection{第七章 [VI.]——论义德的圆满的四个问题：（一）人能否在今世无罪。}

那些断言一个人能在今世没有罪而存在的人，不应立即以不谨慎的轻率加以反对；因为如果我们否认这种可能性，我们就既损害了人的自由意志——人出于自己的意愿而渴望它——也损害了天主的权能或仁慈——他借着其帮助而实现它。但这是一个问题：他是否\Emph{能够}存在；另一个问题：他是否\Emph{实际存在}。此外，一个问题：如果他在能够存在时却不存在，他\Emph{为什么}不存在；另一个问题：\Emph{是否}这样一个从未犯过罪的人，不仅实际上存在，而且是否曾经能够存在，或者将来能够存在。现在，如果按照这四组疑问命题的顺序，有人问我：人是否有可能在今世无罪？我会承认这种可能性，借着天主的恩宠和人的自由意志；并不怀疑自由意志本身应归于天主的恩宠，换句话说，归于天主的恩赐——不仅在于它的存在，而且在于它的善，即它转向遵行天主的诫命。因此，天主的恩宠不仅指明应当做什么，而且帮助人能够做它所指明的事。\BibleQuote{我们有什么不是领受的呢？}\BibleRef{格前 4:7}耶肋米亚也说：\BibleQuote{主，我知道，人的道路不由自己；行路、引导脚步，也不由人。}\BibleRef{耶 10:23}相应地，当人在圣咏中对天主说：\Quote{你曾命令我殷勤遵守你的诫命，}他立刻加上一句不是关于自己的信心，而是渴望能够遵守这些诫命的话：\Quote{但愿我的道路，}他说，\Quote{被引导遵守你的律例！那样我注视你所有的诫命时，就不至于羞愧。}现在谁曾为那些他已完全掌握、无需进一步帮助就能得到的东西而渴望呢？然而，他将他渴望的对象——不是指向运气、命运或除天主之外的任何别一位——在以下话语中足够清楚地表明，他说：\Quote{求你用你的言语引导我的脚步；不要让任何不义管辖我。}从这种可恶的辖制中，那些主耶稣赐给他们权能成为天主子女的人得到了解放。\BibleRef{若 1:12}从如此可怕的统治中，他们是得释放的人，主对他们说：\BibleQuote{如果子使你们自由，你们就真正自由了。}\BibleRef{若 8:36}从这些和许多其他类似的证言，我毫不怀疑天主没有给人下过不可能的命令；并且，借着天主的帮助和援助，没有什么是不可能的，他所命令的得以实现。这样，一个人如果愿意，可以借着天主的助佑而无罪。\par

% structure: p0017 source=h2 target=subsection reason=relative-heading-rank; base=h2

\subsection{第八章 [VII.]——（二）今世是否有无罪之人。}

\Emph{第二。} 若然有人问我所提议的第二个问题——是否有无罪的人——我认为没有。因为我更相信圣经，经上说：\BibleQuote{求祢不要与仆人进入审判，因为在祢面前凡活着的人没有一个是义的。}因此需要天主的仁慈，这仁慈\BibleQuote{欢喜胜过审判}\BibleRef{雅 2:13}，而若人不行仁慈，他将得不到这仁慈。\BibleRef{雅 2:13} 先知说：\Quote{我说：我要向主承认我的过犯，祢就赦免了我心中的罪孽。}接着立即又说：\Quote{为此，凡圣人都应在适宜的时候向祢祈祷。}并非每个罪人，而是\Quote{每个圣人}；因为是圣人的声音说：\BibleQuote{如果我们说我们没有罪，便是欺骗自己，真理不在我们内。}\BibleRef{若一 1:8} 因此，在同一位宗徒的《默示录》中，我们读到\Quote{十四万四千}圣人，\BibleQuote{他们没有与女人沾染，因为他们是童身；在他们口中找不出诡诈，因为他们是没有瑕疵的。}\BibleRef{默 14:3} 他们确实\Quote{没有瑕疵}，无疑是因为他们真正地指责自己；并且正因如此，\Quote{在他们口中没有发现诡诈}——\BibleQuote{因为他们若说没有罪，便是自欺，真理不在他们内。}\BibleRef{若一 1:8} 当然，真理不在的地方，就有诡诈；当义人以指责自己开始陈述时，他确实不说谎。\par

% structure: p0019 source=h2 target=subsection reason=relative-heading-rank; base=h2

\subsection{第九章 更新的开始；复活称为再生；他们是天主的儿子，过着适合新生活的生活。}

因此，在经文\BibleQuote{凡由天主生的，就不犯罪，他也不能犯罪，因为天主的种子存留在他内}\BibleRef{若一 3:9}以及其他类似意义的经文中，他们由于对圣经考虑不足而大大自欺。因为他们没有注意到，人是在开始以新精神生活，并依照创造者的形象更新内在之人时，才成为天主的儿子。因为从人受洗的那一刻起，并非他所有的旧病都消除了，而是更新始于所有罪的赦免，并开始有智慧的人就属灵地有智慧。其余的一切则在希望中完成，期待它们也实际实现，直到身体本身在那更好的不朽与不腐状态中得到更新，这状态是我们在死人复活时所穿上的。因为主也称此为\Quote{再生}——当然不像通过圣洗所发生的那样，但仍然是一种再生，其中在精神中现已开始的工作也将在身体中得以完成。祂说：\BibleQuote{在再生的时候，}祂说，\BibleQuote{当人子坐在祂光荣的宝座上时，你们也要坐在十二个宝座上，审判以色列十二支派。}\BibleRef{玛 19:28} 因为无论圣洗中罪的赦免多么完整和全面，如果籍此立刻使人在永恒的新意中得到完整和全面的改变——我不是说身体上的改变，这身体现在明显地不断趋向旧朽和死亡，之后要更新为完全和真正的新意——而是排除身体，如果在灵魂本身，即内在之人中，在圣洗中已经完成了完美的更新，宗徒就不会说：\BibleQuote{即使我们外在的人日渐朽坏，但内在的人却日日更新。}\BibleRef{格后 4:16} 无疑，那仍日日更新的人尚未完全更新；而就他尚未完全更新而言，他仍处于旧状态。因此，即使人在受洗后，仍在某种程度上处于旧状态，所以他们仍是世界之子；但就他们也进入了新状态而言，即藉由完全而完美的罪赦，并就他们是属灵的且以此行事而言，他们是天主的儿子。我们在内在脱去旧人，穿上新人；因为那时我们弃绝谎言，说真话，并做宗徒所说脱去旧人、穿上新人的其他事，这新人按天主的形象，在真实的正义和圣善中被创造。\BibleRef{弗 4:24} 他劝勉已经受洗和信主的人这样做——如果这完全而彻底的改变已在他们的圣洗中完成，这样的劝勉就不适合他们了。然而，这改变确实发生了，因为我们当时实际上已被\Emph{拯救}；因为\BibleQuote{祂藉着再生之洗拯救了我们。}\BibleRef{铎 3:5} 但在另一处，他告诉我们这事如何发生。他说：\BibleQuote{不但是他们，}他说，\BibleQuote{就是我们这些有圣神初果的人，也在自己内心叹息，等待着义子名分，即我们身体的救赎。因为我们得救是在于希望；但可见的希望不是希望；因为人看见的，为何还希望呢？但如果我们希望那看不见的，我们就耐心等待。}\BibleRef{罗 8:23}\par

% structure: p0021 source=h2 target=subsection reason=relative-heading-rank; base=h2

\subsection{第十章 [VIII.] 完美何时实现。}

因此，我们完全成为子女，要等到我们身体得赎之时。我们现在拥有圣神的初果，由此我们已经真正成为天主的子女；至于其余部分，因着望德我们得救和更新，同样我们也成为天主的子女。但因为我们尚未实际得救，我们也尚未完全更新，也尚未完全成为天主的子女，而是属于世界的子女。因此我们在更新和圣德的生活中前进——藉此我们成为天主的子女，也藉此我们不能犯罪——直到最后，那使我们至今仍属世界子女的一切都被转化为这——因为正是由于这一点，我们至今仍能犯罪。由此便发生：\Quote{凡由天主所生的，不犯罪；}\BibleRef{若一 3:9}同样，\Quote{如果我们说我们没有罪，便是欺骗自己，真理也不在我们内。}\BibleRef{若一 1:8}那时，那使我们成为血肉和世界之子的内在部分将被终结；而那使我们成为天主的子女并藉祂的圣神更新我们的部分将被成全。因此，同一若望说：\Quote{亲爱的，现在我们是天主的子女；但将来如何，还没有显明。}\BibleRef{若一 3:2}那么，这些表达中的变化——\Quote{\Emph{我们是}}和\Quote{\Emph{我们将是}}——是什么意思呢？无非是：\Emph{我们是}出于望德，\Emph{我们将是}出于实际。因为他接着说：\Quote{我们知道，当祂显现时，我们将相似祂，因为我们将看见祂实在怎样。}\BibleRef{若一 3:2}因此，我们现在已经开始相似祂，拥有圣神的初果；然而，由于旧性的残余，我们仍与祂不相似。因此，在相似祂的程度上，我们是藉着更新的圣神成为天主的子女；但在不相似祂的程度上，我们是血肉和世界之子。一方面，我们不能犯罪；但另一方面，如果我们说我们没有罪，我们便是自欺——直到我们完全进入子女的名分，罪人不再存在，你寻找他的地方，却找不着。\par

% structure: p0023 source=h2 target=subsection reason=relative-heading-rank; base=h2

\subsection{第十一章 [IX.]——伯拉纠派的一个异议：为何义人不能生义人？}

所以，他们中有些人这样辩论是徒劳的：\Quote{如果罪人生罪人，以致他初生的儿子必须藉领受圣洗才能除灭原罪的罪债，那么义人也同样应该生义子。}好像人是因自己的义而生子于肉身，而不是因那存在于他肢体中的贪欲而动，且罪恶之律被他心神之律用于生育的目的。因此，他生儿育女显示出他仍保留着这世界之子的旧性；这并非出于他被提升至天主子民中新生命的事实。因为\Quote{今世之子也娶也嫁。}\BibleRef{路 20:34}因此，凡由他们所生的也相似他们；因为\Quote{由肉身生的就是肉身。}\BibleRef{若 3:6}然而，只有天主的子女是义人；但就他们是天主的子女而言，他们并不按肉身生育，因为他们是由圣神而非由肉身所生。但他们中为人父母的，是由于他们尚未脱去旧性的全部残余以换得那等待着他们的完全更新而生育。因此，凡是在他父亲本性这种旧而软弱的状态中所生的儿子，必然也分沾同样的旧而软弱的状态。所以，为了使他被重生，他自己也必须藉着圣神通过罪的赦免而更新；如果这种改变不发生在他身上，他的义父对他毫无用处。因为他是因圣神而成为义人，但他并非因圣神生了他的儿子。另一方面，如果这种改变发生在他身上，不义的父也不会损害他：因为他因圣神的恩宠已进入了永恒新生命的希望中；而他的父亲则因属肉体的心思完全停留在旧性之中。\par

% structure: p0025 source=h2 target=subsection reason=relative-heading-rank; base=h2

\subsection{第十二章 [X.]——他调和圣经中的一些段落。}

因此，\Quote{凡由天主所生的，不犯罪}\BibleRef{若一 3:9}这句话与那处由天主子民所宣告的\Quote{如果我们说我们没有罪，便是欺骗自己，真理不在我们内}\BibleRef{若一 1:8}并不矛盾。因为无论一个人的现世望德多么完全，无论他因圣神的再生在他人性的那部分中更新多么真实，他仍然带着一个腐败的身体，这身体重压着灵魂；只要这样，即使在同一个人身上，也必须区分每个行为的来源和关系。我想，在圣经中很难找到像关于祂的三位仆人诺厄、达内尔和约伯所写的那样重的成圣见证，厄则克耳先知描述他们是唯独能从天主即将来临的愤怒中被解救的人。\BibleRef{则 14:14}在这三个人身上，他无疑预象了将要被解救的三类人：在诺厄身上，我想，他代表公义的民族领袖，因为他管理作为教会预象的方舟；在达内尔身上，是在节制中成义的人；在约伯身上，是在婚姻中成义的人——且不提这段经文的其他解释，现在无需考虑。无论如何，从先知的这个见证和其他默感陈述可以清楚看出，这些圣人在义德上多么卓越。然而，没有人应该因他们的历史而说，例如，醉酒不是罪，虽然如此好人曾陷入其中；因为我们读到诺厄曾一度醉酒，\BibleRef{创 9:21}但愿不会有人认为他是惯常醉酒者。\par

% structure: p0027 source=h2 target=subsection reason=relative-heading-rank; base=h2

\subsection{第十三章——伯拉纠派的一个遁词。}

达尼尔在向天主倾心祈祷之后，确实论及自己说：\Quote{当我正在祈祷，承认我的罪和我人民的罪，在上主我的天主面前。}\BibleRef{达 9:20} 如果我没有弄错的话，这就是为什么在以上提到的厄则克耳先知书中，有一个极其骄傲的人被问说：\Quote{那么你比达尼尔更聪明吗？}\BibleRef{则 28:3} 在这一点上，绝不能说某些人反对主祷文所主张的话。他们说：\Quote{因为虽然那祈祷是由宗徒们献上的，他们在成为圣善和成全、没有任何罪之后，但那祈祷不是为了他们自己，而是为了那些不完全且仍有罪的人，他们说：\InnerQuote{宽免我们的罪债，如同我们也宽免得罪我们的人。}他们用了\Emph{我们的}这个词，}他们说，\Quote{是为了表明在一个身体内，既包含那些仍有罪的人，也包含他们自己，他们已完全脱离罪。}现在这当然不能用在达尼尔身上，他（我想）作为先知预见了这种傲慢的意见，所以在祈祷中他屡次说：\Quote{我们犯了罪；}并向我们解释他为何这样说，并非要我们听到他说：当我正在祈祷，向我的天主承认我人民的罪；也不是混淆区别，以致于不确定他是否因同一身体的共融而说：当我向上主我的天主承认\Emph{我们的}罪；而是以如此清晰明确的言辞表达自己，仿佛他自己充满这种区别，并特别想引起我们的注意：\Quote{\Emph{我的罪}，}他说，\Quote{\Emph{和我人民的罪}。}谁能反驳这样的证据呢？只有那些更乐意辩护自己所想、而不愿发现应该所想的人。\par

% structure: p0029 source=h2 target=subsection reason=relative-heading-rank; base=h2

\subsection{约伯并非无罪。}

但让我们看看约伯在天主对他义德给予伟大见证之后，他自己怎么说。他说：\Quote{我确实知道是这样：因为凡人怎能在天主面前成为正义的呢？如果天主与他进入审判，他必不能听从祂。}\BibleRef{约 9:2} 不久之后他问道：\Quote{谁能抗拒祂的审判？即使我似乎义，我的口也会说出亵渎的话。}\BibleRef{约 9:19} 再后来他又说：\Quote{我知道祂必不放过我。既然我是不敬虔的，为什么我不死呢？我若用雪洗濯，用洁净的手洁净，你却将我染得污秽不堪。}\BibleRef{约 9:30} 在另一篇讲论中他说：\Quote{因为祢记述了恶事攻击我，用我青年的罪环绕我；祢把我的脚放在木枷中。祢察看我一切的行径，并窥探我的脚掌，它们旧如皮囊，或如虫蛀的衣服。因为妇人所生的人，时日短少，充满愤怒；如花盛开，随即凋谢；如影逝去，不复存留。祢岂不也察看他，使他与祢进入审判？因为谁能从污秽中洁净？连一个也没有；即使他的生命只存留一天。}然后稍后他说：\Quote{祢数算了我一切的需要；我哪一样罪没有逃过祢的眼目。祢将我的过犯封在囊中，并标记了我一切不愿做的事。}\BibleRef{约 14:16} 看，约伯也承认自己的罪，并说他确信在天主面前没有义人。所以他也确信，如果我们说自己无罪，真理就不在我们内。因此，当天主按照人类行为的标准赐予他崇高的义德见证时，约伯自己却从那义德的标准——他尽可能在天主内看见的标准——来衡量，确实知道事实如此；他随即说：\Quote{凡人怎能在天主面前成为正义的呢？因为如果天主与他进入审判，他必不能听从祂。}换句话说，如果被挑战进入审判，他希望表明在他内找不到任何可被定罪的事，\Quote{他必不能听从祂，}因为他甚至缺少那种能使他听从那位教导应当承认罪的主的服从。因此，〔主〕责备某些人说：\Quote{你们为什么与我争辩审判？}\BibleRef{耶 2:29} 〔圣咏作者〕避免这事说：\Quote{求祢不要与祢的仆人进入审判，因为在祢面前凡有血气的没有一个能成为义。}与此相符，约伯也问：\Quote{因为谁能抗拒祂的审判？即使我似乎义，我的口也会说出亵渎的话；}意思是：如果违背祂的审判，我称自己为义，而祂完美的义德准则证明我是不义的，那么我的口的确会说出亵渎的话，因为它会反对天主的真理。\par

% structure: p0031 source=h2 target=subsection reason=relative-heading-rank; base=h2

\subsection{血肉的世代因原罪而被定罪。}

他指出，血肉生育的这种绝对软弱，更确切地说，这种定罪，来自原罪的过犯。当讨论自己的罪过时，他似乎阐明了它们的原因，说：\BibleQuote{人由女子所生，时日短少，且充满忿怒。}这是什么样的忿怒呢？岂不是宗徒所说的\BibleQuote{本性的}，即本源的\BibleQuote{义怒之子}？\BibleRef{弗 2:3} 因为他们属于肉体和世俗的私欲偏情？他进一步表明，人的死亡也与这同样的忿怒有关。因为说了\BibleQuote{他时日短少，且充满忿怒}之后，他又加上：\BibleQuote{如绽放的花朵，他即凋谢；如阴影逝去，不再存留。}接着他附加说：\BibleQuote{祢岂不要他进入与祢的审判？谁能由不洁中洁净？连一个也没有；即使他的生命只存留一日。}这些话实际上是在说：祢将进入与祢审判的责任放在了短寿的人身上。因为无论生命多么短暂——即使只有一天——他都不可能洁净无污；因此，他完全有理由必须接受祢的审判。然后，他又说：\BibleQuote{祢数算了我所有的需要，没有一项罪逃过祢；祢将我的过犯封在囊中，并记下了我所不情愿做的一切；}这难道不够清楚吗？即使那些不是因快乐诱惑而犯、而是为了避免麻烦、痛苦或死亡而犯的罪，也正当地被归咎。而这些罪也被说成是在某种必要下犯的，尽管它们都应以对正义的爱和喜悦来克服。再者，他在\BibleQuote{祢记下了我所不情愿做的一切}这句话中所说的，显然可以与这句话相联系：\BibleQuote{我所愿意的，我不做；我所憎恨的，我反而去做。}\BibleRef{罗 7:15}\par

% structure: p0033 source=h2 target=subsection reason=relative-heading-rank; base=h2

\subsection{第十六章——约伯预知基督要来受苦；完美者之谦卑之路}

值得注意的是，主亲自在圣经中，即由天主圣神，向约伯作证说：\BibleQuote{在一切所遭遇的事上，他没有以嘴唇犯罪得罪主。}\BibleRef{约 1:22}但后来约伯自己告诉我们，主又用责斥对他说话：\BibleQuote{我为何还要辩护，而被训斥，且听到主的责斥？}\BibleRef{约 39:34}没有人被公正地责斥，除非他里面有应受责斥的事。[XI.] 这是怎样的责斥呢？——而且，这被认为是由我主基督的位格而来。祂向他述说祂全能的一切神圣作为，以此责斥他，似乎对他说：\BibleQuote{你能像我一样成就所有这些大能的工作吗？}但这全部目的岂不是为了让约伯明白（因为这种教导是他从神感而来的，使他预知基督要来受苦）——让他明白他应当如何忍耐地承受他所经历的一切，因为基督，虽然祂为我们成为人时完全无罪，而且作为天主拥有如此大的能力，却绝不拒绝服从以至受死？当约伯以更纯洁的心灵热忱领悟这一点时，他给自己的回答加上这些话：\BibleQuote{我以前只凭耳朵听闻到祢，但如今我的眼亲眼看见了祢；因此我厌恶自己，消灭自己，将自己看作尘土和灰烬。}\BibleRef{约 42:5}他为何如此深深地对自己不满？天主的工作，即他是人，本不应使他不满，因为甚至对天主自己说：\BibleQuote{不要轻视祢双手所造的。}正是鉴于他在其中发现自己的不义的公义，他才厌恶自己、消灭自己，并自视为尘土和灰烬——因为他心中看见了基督的公义，在祂身上不可能有任何罪，不仅就祂的天主性而言，而且就祂的灵魂和肉身而言。也是鉴于这来自天主的公义，宗徒保禄，虽然\BibleQuote{就法律上的义而言，他是无可指摘的}，却\BibleQuote{将万事看作}不仅是损失，而且是粪土。\BibleRef{斐 3:6}\par

% structure: p0035 source=h2 target=subsection reason=relative-heading-rank; base=h2

\subsection{第十七章 [XII.]——没有人在一切事上为义人}

因此，那句关于约伯被赞扬的卓著的天主见证，与那句所说\Quote{在祢面前，凡有血肉的都不能称为义}并不矛盾；因为它并不导致我们以为在他身上完全没有可被他自己真实地或被主天主正当地责备之处，尽管同时他也可能被真实地称为义人、真诚的天主敬拜者、远离一切恶行的人。因为关于他，天主这样说：\BibleQuote{你曾用心观察我的仆人约伯吗？因为地上没有像他这样的人，纯洁、正直、真诚敬拜天主、远离一切恶事}\BibleRef{约 1:8}。首先，他在这里因与地上所有人相比而受到称赞。因此他胜过当时地上所有能成为义的人；然而，因为他在这义上的优越性，他并非完全没有罪。其次，他被说成是\Quote{\Emph{纯洁的}}——没有人能公正地控告他生活中的任何过错；\Quote{\Emph{正直的}}——他在道德正直上进步如此之大，以至于没有人能与他相提并论；\Quote{\Emph{真诚敬拜天主的}}——因为他是一个真诚而谦卑地承认自己罪过的人；\Quote{\Emph{远离一切恶事的}}——如果这也延伸到一切恶言恶念，那将是了不起的。约伯到底是一个多么伟大的人，我们并没有被告知；但我们知道他是一个义人；我们也知道，在忍受可怕的痛苦和考验中，他是伟大的；我们还知道，他之所以要承受如此多的苦难，并非因为他的罪，而是为了彰显他的义。但主赞扬约伯的话，也可应用于那\BibleQuote{按内心是喜悦天主的法律，却看到肢体中另有一个法律，与理智的法律交战}\BibleRef{罗 7:22}的人；尤其是当他说：\BibleQuote{我愿意行的善，我反而不行；我不愿意作的恶，我倒去作。但如果我做我不愿意作的事，就不是我做的，而是住在我内的罪做的}\BibleRef{罗 7:19}。注意，他按内心也是远离一切恶行的，因为这样的恶行不是他自己造成的，而是住在他肉体中的邪恶造成的；然而，因为他连那喜悦天主法律的能力也只是来自天主的恩宠，所以他仍需要救援，喊道：\BibleQuote{天主的恩宠，藉着我们的主耶稣基督！}\BibleRef{罗 7:24}。\par

% structure: p0037 source=h2 target=subsection reason=relative-heading-rank; base=h2

\subsection{第十八章 [XIII.]——完美的人义是不完美的}

因此，世上有义人，有伟人（勇敢、明智、贞洁、忍耐、虔诚、慈悲），他们为了义德而以平静的心忍受各种暂时的恶。然而，如果这些话语是真实的——不，正因为是真实的——\BibleQuote{若我们说我们没有罪，便是欺骗自己}\BibleRef{若一 1:8}，以及\Quote{在祢面前，凡有血肉的都不能称为义}，那么他们并非没有罪；也没有一个人如此骄傲愚蠢，以至于不因自己众多的罪而认为主祷文是必需的。\par

% structure: p0039 source=h2 target=subsection reason=relative-heading-rank; base=h2

\subsection{第十九章——匝加利亚和依撒伯尔，罪人}

那么，对于经常在讨论这个问题时被人用来反对我们的匝加利亚和依撒伯尔，我们该说什么呢？除了圣经中明确记载\BibleRef{路 1:6}匝加利亚是司祭中的一位卓越义人，他的职责是献上旧约的祭献。然而，我们在\WorkTitle{致希伯来人书}（我在前一本书中已引用过那段话）中也读到，基督是唯一的大司祭，祂不需要像那些被称为大司祭的人那样，每日先为自己的罪献祭，然后为百姓献祭。经上说：\BibleQuote{因为这样的大司祭才合适我们，他是圣善的、无辜的、无玷的、从罪人中分别出来的、高于诸天的；祂不需要像那些大司祭那样，每日先为自己的罪献祭}\BibleRef{希 7:26}。在这里提到的司祭中有匝加利亚，其中有丕乃哈斯，是的，还有亚郎本人，这司祭职位是从他开始的，以及所有在这司祭职位上生活可称赞和正义的人；然而所有这些人都必须先为自己的罪献祭——唯独基督，他们预表祂的未来来临，作为一位无玷的司祭，没有这样的需要。\par

% structure: p0041 source=h2 target=subsection reason=relative-heading-rank; base=h2

\subsection{第二十章——配为宗徒之长的保禄，仍是罪人}

然而，给予匝加利亚和依撒伯尔的称赞，有什么不是包含在宗徒关于自己未信基督之前所说的话中呢？他说：\BibleQuote{就法律上的义而言，我原是无可指摘的。}\BibleRef{斐 3:6} 他们也同样被记载说：\BibleQuote{他们二人在天主面前都是义人，遵行主的一切诫命和典章，无可指摘。}这是因为他们所有的义，在人前并非虚伪，所以经上说他们\Quote{在\Emph{主面前}行走。}但记载匝加利亚和他妻子的话是\Emph{在一切诫命和典章上}，宗徒简单地将它表达为\Emph{在法律上}。原来，在福音之前，他有一部法律，他们也有一部法律，并非他有一部、他们另有一部。他们同一部法律，即梅瑟传给他们祖先的，按此法律，匝加利亚也作司铎，按他的班次献祭。然而，那时拥有同样义德的宗徒接着说：\BibleQuote{但是，凡以前对我有益的，我现在因基督都视为损失。不但如此，我也将万事看作损失，因为我以认识我主基督耶稣为至宝；为了祂，我不但将万事视为损失，更看作粪土，为赢得基督，并得以在祂内，不是拥有我自己那出于法律的义，而是藉着对基督的信德所有的义，即来自天主、基于信德的义；为使我知道祂和祂复活的德能，并参与祂的苦难，相似祂的死，希望我能达到从死者中的复活。}\BibleRef{斐 3:7} 因此，我们绝不能从圣经描述他们的字句，就以为匝加利亚和依撒伯尔持有毫无犯罪的完美义德；相反，我们甚至必须认为，宗徒本人按照同一规则，也并不完美，不仅在他所拥有、与他们相同的法律义德上——他将这义德视为损失和粪土，以换取那因信基督的至义——而且就在福音本身，他应得大宗徒职位的优越性。我不敢说这话，除非我认为拒绝相信他本人是极其错误的。他延伸我们所引的段落，说：\BibleQuote{这不是说我已经达到或已经成为完美的；我只在追赶，或许我能夺得基督耶稣之所以夺取了我的。弟兄们，我还不认为自己已经夺得了；我只有一件事：忘记背后的，努力面前的，向着标竿直跑，为要得着天主在基督耶稣里从高天召我来得的奖赏。}\BibleRef{斐 3:12} 在这里，他承认自己尚未达到，尚未完美于那渴望在基督内获得的圆满义德；他仍在向着标竿直跑，忘记背后，努力面前。因此，我们知道他在别处所说的话甚至对他自己也真实：\BibleQuote{我们外在的人虽然日渐朽坏，但内在的人却日日更新。}\BibleRef{格后 4:16} 虽然他已是一个完美的行者，却尚未达到旅程的完美终点。他愿意所有这样的人同他一起作赛跑的同伴。他在我们前一引文之后的话中表达了这一点：\BibleQuote{所以，我们凡是完美的人，都该怀有这种心情；若你们在什么事上有别的想法，天主也会将这事启示给你们。但我们已经达到了什么，就该照这个准则前行。}\BibleRef{斐 3:15} 这个\Quote{前行}不是用身体的双腿，而是以灵魂的情操和生活的品行，以致拥有义德的人能够达到完美；他们每日沿着信德的直道在更新中前进，在这同一义德中此刻已作为行路人而成为完美的。\par

% structure: p0043 source=h2 target=subsection reason=relative-heading-rank; base=h2

\subsection{第二十一章 [XIV.]——所有义人都是罪人。}

同样，所有在圣经中被描述为在今生表现出善意和义行的人，以及此后所有像他们那样生活的人——即使缺乏同样的圣经见证；或所有现在或将来如此生活的人：他们都是伟大的，都是义人，都确实值得称赞——但他们绝无罪；因为，根据同一圣经——它使我们相信他们的美德——我们也相信\Quote{在祢面前凡活着的人，没有一个是义人}，因此所有人都求祂\Quote{不要进入审判与祢的仆人争辩}；而且，不仅是所有信徒普遍需要，而且他们每一个人都特别需要主传授给祂门徒的主祷文。\par

% structure: p0045 source=h2 target=subsection reason=relative-heading-rank; base=h2

\subsection{第二十二章 [XV.]——伯拉纠派的一个异议；完美是相对的；在义德上进步很大的人，方可正确地说他义德完美。}

\Quote{可是，}他们说，\Quote{主说：\InnerQuote{你们应当是成全的，如同你们的天父是成全的一样。}，\BibleRef{玛 5:48}——祂不会下这个命令，如果祂知道祂所命令的是不可能做到的话。}\newline{}现在的关键不是问：如果人为了成全而领受了恩宠，在今生能否无罪——因为可能性的问题我们已经讨论过了；——我们现在要考虑的是，是否有人实际上达到了成全。不过，我们已经认识到，没有人愿意像责任所要求的那样多，正如我们上面大量引用的圣经的见证所宣告的那样。的确，当成全归于某个人时，我们必须仔细看他在哪方面被说成是成全的。因为我刚刚引用了宗徒的一段话，他在其中承认自己在所渴望的义德上尚未成全；但他立刻又说：\Quote{我们凡是成全的，都应怀有这种心情。}如果他不是在某一件事上成全，在另一件事上不成全，他肯定不会说出这两句话。例如，一个人作为追求智慧的弟子可以是成全的：而这还不能对那些格林多人说，他对他们说过：\Quote{我给你们喝的是奶，不是饭，因为那时你们还不能吃，就是如今还是不能。}\BibleRef{格前 3:2}而那些可以说的人，他对他们说：\Quote{我们在成全的人中，也讲智慧}——当然是指\Quote{成全的弟子}。所以，如我所说，一个人可能已经作为弟子而成全，但尚未作为智慧的老师而成全；可能作为学生而成全，但尚未作为行义者而成全；可能作为爱仇者而成全，但尚未在忍受他们的侵害上成全。即使对于那样成全、以致爱所有的人（因为他甚至达到了爱仇的程度）的人，仍然要问他在那爱中是否成全——换句话说，他是否按照真理不变的爱所命令的，去爱他所爱的人。因此，每当我们在圣经中读到关于任何人的成全时，必须仔细看是在哪方面说的，因为一个人并非因在某方面被描述为成全，就被理解为完全无罪；尽管这个术语也可以用来表示，并非没有余地让人继续走向成全，而是他确实进步了很多，因此配得上这个称号。这样，一个人可以说在法律科学上成全，即使有些事情他还不知道；同样，宗徒称那些人为成全的，同时对他们说：\Quote{然而，我们若在什么事上存有不同心意，天主也要将这事启示给你们。但是，我们到了什么地步，就当照着什么地步行。}\BibleRef{斐 3:15}\par

% structure: p0047 source=h2 target=subsection reason=relative-heading-rank; base=h2

\subsection{第二十三章——天主为何命令祂知道无法遵守的事。}

我们不可否认，天主命令我们应当在行义上成全，以致完全无罪。任何事，无论是什么，除非天主禁止它，就不是罪。那么，他们说，祂为何命令祂知道没有活人能做的事呢？同样，也可以问：祂为何命令起初只有两个人的人类去做祂知道他们不会服从的事呢？因为不能说祂发出那个命令，是为了如果我们不服从，有些人会服从；因为天主只命令他们，没有别人，不可吃那棵特定树上的果子。因为，正如祂知道他们不会履行多少义德，祂也知道祂自己打算对他们采取什么正义的措施。同样，祂命令所有的人不可犯罪，尽管祂预先知道没有人会遵守这命令；为的是在那些不虔敬地、可耻地藐视祂诫命的人身上，祂可以公平地惩罚他们；而在那些虔诚地、顺从地努力遵守祂的命令，虽未能完全遵守祂所吩咐的一切，却像他们自己希望被宽恕那样宽恕别人的人身上，祂可以仁慈地净化他们。因为，如果没有罪，天主的慈悲怎能赐予宽恕者宽恕呢？或者，如果有罪，天主的公义怎能不发出禁令呢？\par

% structure: p0049 source=h2 target=subsection reason=relative-heading-rank; base=h2

\subsection{第二十四章——白拉奇派的一个异议。宗徒保禄在有生之年并非无罪。}

\Quote{但是看}，他们说，\Quote{宗徒如何说：\BibleQuote{这场好仗，我已打完；信德，我已保持；赛程，我已跑完；从此以后，正义的冠冕已为我预备好了。}\BibleRef{弟后 4:7} 他若还有任何罪，就不会说这话。}那么，他们要解释，当还有那伟大的冲突、他所预告的那惨重而极度的苦难等待着他时，他怎么还能说这话。\BibleRef{弟后 4:6} 为了跑完他的赛程，难道只差一件小事吗？事实上，那还剩下来要忍受的，正是更凶猛、更残忍的敌人。然而，如果他因那启示他苦难临近的那一位使他确实、安稳而说出如此喜乐的话，那么他说这话，不是出于实现的圆满，而是出于希望的坚定，并把他预见将要发生的事当作已经完成的事来叙述。因此，如果他在那些话上加上了\Quote{我再没有罪}这样的声明，我们必须理解他那时所说的完美是源于未来的前景，而非已成的事实。因为他说这话时被认为已经完成的\Quote{无罪}，属于赛程的结束；正如他在苦难的决定性争战中战胜仇敌也属于赛程的结束——虽然连他们自己也必须承认，他说这话时，这事仍有待成就。因此，\Emph{我们}断言，当宗徒以对天主应许的完全信赖说这一切已经实现时，这一切还等待实现。因为正是在赛程的结束中，他宽恕了那些得罪他的人的罪，并祈求自己的罪也得到同样的宽恕，正是由于对这主应许的最确定的信赖，他相信在那仍属于未来的最后终结中，他将没有罪——即使他在信赖中把它说成已经实现。现在，撇开其他一切考虑，我想知道，当他说出被认为暗示他没有罪的话时，那\Quote{肉身上的刺}是否已经从他身上拔除了？他为此三次求主，并得到这样的回答：\BibleQuote{我的恩宠够你用，因为我的能力在软弱中才得以完全。}\BibleRef{格后 12:8} 为使如此伟大的人达至成全，需要那\Quote{撒殚的使者}不被拔除，他因此受打击，\BibleQuote{免得他因那极大的启示而过于高傲}，\BibleRef{格后 12:7} 那么，是否有任何人大胆到认为或说，任何一个必须在今生重担下弯腰的人完全没有一切罪？\par

% structure: p0051 source=h2 target=subsection reason=relative-heading-rank; base=h2

\subsection{第二十五章——天主在义怒中并在仁慈中惩罚}

虽然有些人在义德上如此卓越，以致天主从云柱中向他们说话，例如\BibleQuote{在祂的司祭中，梅瑟和亚郎，并在呼求祂名的人中，撒慕尔}——后者在真理的圣经中，从他幼年时起，因他母亲为还愿将他放在天主圣殿中，并奉献他为主仆人，而大受赞扬——然而，关于这样的人也写着：\BibleQuote{天主，祢曾向他们施恩，却惩罚了他们的一切计谋。}天主的愤怒惩罚义怒之子；而祂以仁慈惩罚恩宠之子，因为\BibleQuote{主爱谁就纠正谁，并鞭打祂所接纳的每个儿子。}然而，没有一种惩罚、纠正或鞭打不是由于罪，除了那一位，祂为击打者预备自己的背，以便在一切事上像我们一样，却没有罪，以便成为圣人的圣司祭，甚至为圣人代祷——这些圣人毫无虚假的牺牲，每人甚至为自己说：\BibleQuote{求祢宽恕我们的过犯，如同我们也宽恕那些冒犯我们的人。}因此，即使我们在争论中的对手，他们的生活贞洁，品行可嘉，他们并不迟疑去做主吩咐那富人的事——那富人曾询问主关于获得永生的事，在回答主的第一个问题时，他声称已经完全遵守了法律中的每一条诫命——即\BibleQuote{若他愿意成全，就要变卖一切所有，分施给穷人，并将财宝转移到天上}；\BibleRef{玛 19:12} 然而，他们没有一个人胆敢说自己无罪。但我们相信，他们不说这话是出于欺骗的意图；而如果他们撒谎，就在这行为中，他们开始增加或犯下罪。\par

% structure: p0053 source=h2 target=subsection reason=relative-heading-rank; base=h2

\subsection{第二十六章〔十七〕——（3）为何今世无人无罪}

[\Emph{3d.}] 现在让我们考虑我提到的第三个问题。既然藉着天主的恩宠扶助人的意志，人可能在此生中没有罪，为什么他没有？对这个问题，我可以很容易并真实地回答：因为人不愿意。但若问我为什么不愿意，我们就会陷入冗长的陈述。然而，在不妨碍更仔细的考察下，我可以简单地说这一点：人不愿意做正确的事，要么是因为他们不知道什么是正确的事，要么是因为它令他们不悦。因为我们越是确知一件事的美好，越是在其中感到喜悦，就越热切地渴望它。因此，无知和软弱是阻碍意志行动或行善或止恶的缺陷。但是，隐藏的事得以显明，不悦的事变得可喜悦，这乃是天主的恩宠，祂扶助人的意志；而人没有得到这恩宠的帮助，其原因也在他们自身，不在天主，无论他们是因骄傲的不义而被预定受罚，或是如果他们属于怜悯之子，则要因相反的骄傲而受审判和管教。因此，\Person{耶肋米亚}说：\BibleQuote{上主，我知道，人的道路不由自己，行走和指导自己的脚步也不属于任何人，}\BibleRef{耶 10:23}紧接着又说：\BibleQuote{上主，求祢管教我，但要有节制，不要在祢的震怒中管教；}\BibleRef{耶 10:24}这等于说：我知道，因为我祢给我的帮助太少，我的脚步未能完全导向正路，这是祢对我纠正；但求祢不要用祢的怒气对待我，如同祢决定惩罚恶人时那样，而要用祢的审判，藉此祢教导祢的孩子不要骄傲。因此，在另一处说：\BibleQuote{祢的审判将帮助我。}\par

% structure: p0055 source=h2 target=subsection reason=relative-heading-rank; base=h2

\subsection{第二十七章——天主对骄傲的补救}

因此，你不能将任何人过错的起因归给天主。因为所有人类犯错的起因都是骄傲。为了说服和消除骄傲，从天上来了一个伟大的补救。天主在怜悯中自谦，从天降下，向被骄傲抬举的人显示纯全而明显的恩宠，就是祂因巨大的爱为那些分享这爱者所取的人性。因为，就连这一位——如此与天主的圣言结合，以致因那结合祂同时成为天主的独生子，又是同一位的独生子——也不是凭祂自己意志的先在功行而行动。无疑，祂必须是独一的；若有两位、三位或更多，如果这种事可能发生，那就不是出于天主纯粹而单纯的恩赐，而是出于人的自由意志和选择。因此，这一点特别强调给我们；这一点，就我敢想而言，是隐藏于基督内的智慧与知识宝藏中特别教导和学习的属天课程。因此，我们每个人，有时知道，有时不知道——有时喜乐，有时不喜乐——为了开始、继续和完成我们的善工，好使他能知道，他之所以知道或喜乐，不是由于他自己的意志，而是由于天主的恩赐；这样他就被治愈了使他膨胀的虚荣，并知道这话是何等真实，不是指我们的土地，而是属灵地说：\Quote{上主将赐予仁爱和甘饴的恩宠，我们的土地将结出果实。}此外，善工所带来的喜乐更大，是按着天主越来越被爱为最高的不变之善，以及一切种类善事的作者的程度。并且为使天主被爱，\BibleQuote{祂的爱已倾注在我们心中，}不是由于我们自己，而是\BibleQuote{藉着所赐给我们的圣神。}\BibleRef{罗 5:5}\par

% structure: p0057 source=h2 target=subsection reason=relative-heading-rank; base=h2

\subsection{第二十八章 [XVIII.]——善意志来自天主}

然而，人们正努力在我们的意志中寻找一些属于自己的善——不是天主给我们的；但我无法想象如何能找到。宗徒论及人的善工时说：\BibleQuote{你有什么不是领受的呢？既然是领受的，为什么你还夸耀，好像不是领受的呢？}\BibleRef{格前 4:7}但除此之外，即使是理性本身——像我们这样的人在这些事上可以估量的——也严厉地约束我们每个人在研究时，以免我们如此辩护恩宠以致似乎取消了自由意志，或另一方面，如此断言自由意志以致在我们的傲慢不敬中被视为对天主的恩宠忘恩负义。\par

% structure: p0059 source=h2 target=subsection reason=relative-heading-rank; base=h2

\subsection{第二十九章——白拉奇派的遁词}

现在，关于我刚才引用的宗徒经文，有些人会维护其意思说：\Quote{一个人无论有多少善意志，都必须归于天主，理由是——即，若他不是人，连这数量也不能在他内。既然他只有从天主获得成为任何东西、成为人的能力，那么为什么他内的一切善意志，若不依赖其存在者就无法存在，不也应归于天主吗？}但用同样的方式，也可以说一个恶意志也可以归于天主作为其作者；因为连它也不能在人内存在，除非他是一个人在其中存在；但天主是他作为人的存在的作者；因此也是他的恶意志的作者，因为若没有一个人在其中，恶意志就无法存在。但这样论证是亵渎。\par

% structure: p0061 source=h2 target=subsection reason=relative-heading-rank; base=h2

\subsection{第三十章——一切意志，若是善的，便爱正义；若是恶的，便不爱正义}

因此，除非我们不仅获得意志的决定，即自由地向这边或那边转向，并属于那恶人可能恶用的本性之善之一；而且也获得善志，它属于那些不可能被恶用的善之一——除非这不可能性是从天主赐予我们的，我不知道如何辩护那所说的：\Quote{你有什么不是领受的呢？} 因为如果我们从天主那里获得某种自由意志，它仍可以是善或恶；但善志来自我们自己；那么来自我们自己的比来自祂的更好。但是，因为说这话是极其荒谬的，他们应当承认我们甚至从天主那里获得善志。如果意志能这样处于某种中间状态，既不善也不恶，那确实是一件奇怪的事；因为我们或者爱慕正义，那就是善的，如果我们爱慕得更多，就更善——如果爱慕得更少，就更不善；如果我们根本不爱慕它，就不是善的。谁能犹豫不决地肯定，当意志根本不爱慕正义时，它不仅是恶的，甚至是完全堕落的意志呢？因此，既然意志或是善的或是恶的，而且当然我们不是从天主那里获得恶志，那么剩下的就是我们是从天主那里获得善志；否则，我不知，既然我们的成义来自它，我们应当因祂的什么其他恩赐而喜乐呢？因此，我想经上记载着：\BibleQuote{意志是由主预备的；} \BibleRef{箴 8:35}，并在圣咏中：\Quote{人的脚步将由主正确引导，他的道路将是他意志的选择；}以及宗徒所说的：\BibleQuote{因为是天主在你们内工作，使你们愿意，也使你们力行，为成就祂的善意。} \BibleRef{斐 2:13}\par

% structure: p0063 source=h2 target=subsection reason=relative-heading-rank; base=h2

\subsection{第三十一章——恩宠是以仁慈赐给一些人，以公义和真理而不赐给另一些人。}

因此，既然我们离弃天主是我们自己的行为，这是恶志；但我们归向天主是不可能的，除非祂唤醒并帮助我们，这是善志——那么我们有什么不是领受的呢？如果我们领受了，为什么还要夸耀好像没有领受呢？因此，正如\Quote{夸耀的，应当因主夸耀}，这来自祂的仁慈，而非他们的功德，使得天主愿意将恩宠赐给一些人，但出自祂的真理而愿不赐给另一些人。因为对罪人理应处以惩罚，因为\Quote{主天主喜爱仁慈和真理}，并且\Quote{仁慈和真理彼此相遇}，并且\Quote{上主的一切道路都是仁慈和真理}。谁能数算圣经中结合这两种属性的无数例子呢？有时，通过变更术语，\Emph{恩宠}代替\Emph{仁慈}，如经上说：\BibleQuote{我们见了祂的光荣，即父的独生子的光荣，满溢恩宠和真理。} \BibleRef{若 1:14} 有时\Emph{审判}也代替\Emph{真理}，如经上说：\Quote{上主，我要向祢歌唱仁慈与公义。}\par

% structure: p0065 source=h2 target=subsection reason=relative-heading-rank; base=h2

\subsection{第三十二章——天主在其恩宠中的主权。}

至于祂为何愿意转化一些人，而惩罚另一些人因离弃——虽然无人能公正地指责那仁慈者在施恩中的作为，也无人能公正地挑剔那真理者在惩罚中的裁决（正如在葡萄园工人的比喻中，无人能公正地责备祂，给一些人约定的工钱，给另一些人非约定的慷慨赏赐 \BibleRef{玛 20:1}），然而，祂更深隐的审判的目的仍在祂自己的权能中。[XIX.] 就我们所领受的，让我们拥有智慧，并让我们明白，善主天主有时甚至向祂的圣人隐瞒善工的确定知识或得胜的喜乐，正是为了使他们发现，使他们那照亮黑暗的光和使他们的土地结果实的甜蜜恩宠并非来自自己，而是来自祂。\par

% structure: p0067 source=h2 target=subsection reason=relative-heading-rank; base=h2

\subsection{第三十三章——藉着恩宠我们既得识善，也得其带来的喜悦。}

当我们祈求祂赐予我们帮助去行并成就正义时，我们祈求的，岂不就是祂会开启那隐藏之事，将那不悦之事变为甘甜吗？因为就连这祈求祂的本分，也是藉着祂的恩宠才学会的——先前它是隐藏的；又藉着祂的恩宠我们才爱上它——先前它令我们无喜乐——以致\BibleQuote{凡夸耀的，必须夸耀不是在自我内，而是在主内。}的确，被高举而骄傲，是因人自己的意愿，而非天主的作为；因为天主既不催促我们，也不帮助我们去做这样的事。于是在人的意志中首次出现了一种对自身能力的欲望，藉着骄傲而成为悖逆。若不是有此欲望，实则没有什么是困难的；人愿意时，便可毫不困难地拒绝。然而，从罪有应得的惩罚中，产生了这样的缺陷：从此顺服于正义变得困难；除非有助佑的恩宠克服这缺陷，无人会转向圣善；除非有有效的恩宠医治它，无人能享有正义的平安。然而，征服并医治的恩宠，不属于祂——那祈祷所向的祂——又属于谁呢？\BibleQuote{求祢转化我们，我们救恩的天主，求祢转离对祢的愤怒。}而且，祂若这样做，是出于仁慈而做，以致论及祂时说：\BibleQuote{祂没有按我们的罪对待我们，也没有按我们的过犯报复我们}；当祂对任何人隐忍不做此事时，祂是出于审判而隐忍。谁又能对祂说：\BibleQuote{祢做了什么？}当圣人们以虔诚之心歌颂祂的仁慈和审判时，这又如何说呢？因此，即使在祂的圣人和忠信仆人身上，祂也对他们某些过犯施以较晚的治愈，好使他们卷入这些过犯时，在他们可能成就的任何善事（无论是隐藏的还是明显的）上，所体验到的喜乐少于成全全部正义所需的满足；以致于在祂最完美的公平和真理法则面前，\BibleQuote{凡有血肉的，没有一个人在祂面前能成义。}祂其实并不愿我们陷于定罪，而是愿我们谦卑；祂将自己全部同样的恩宠显示给我们。然而，当我们各方面都变得顺遂之后，不要以为那原本属于祂的，是我们的；因为那是最违背宗教和虔敬的错误。我们也不要以为因着祂的恩宠，就能继续固守旧日的罪；反而要对抗那骄傲——正因这骄傲我们才在其中受屈辱——我们尤其要警醒奋斗并热切地向祂祈祷，同时知道正是由于祂的恩赐，我们才有力量如此奋斗和祈祷；这样，在任何情况下，我们不看自己，而是举心向上，感谢主、我们的天主，并且每当夸耀时，只夸耀在祂内。\par

% structure: p0069 source=h2 target=subsection reason=relative-heading-rank; base=h2

\subsection{第34章 [XX.]— (4) 没有人，除基督外，曾活着或能活着而无罪。}

[\Emph{4th.}] 现在剩下我们的第四点，解释完这一点后，在天主助佑下，我们这篇冗长的论述终于可以结束。那就是：那从未有过罪或将来有罪的人，不仅现在活着作为人子之一，甚至是否曾在任何时候存在过，或将来在时间中存在？现在完全确定的是，这样一个人现在不存在，过去不曾存在，将来也不会存在，除非那唯一的天主与人之间的中保、人基督耶稣。我们在论及婴儿圣洗时已对此说了不少；因为若这些婴儿没有罪，那么不仅现在，而且过去将来都有无数无罪的。现在，若我们在第二点中所处理的论点确实成立，即事实上没有人没有罪，那么当然连婴儿也不是没有罪的。由此得出结论：即使一个人今生可能在美德上进展如此之深，以至达到了完全无罪的神圣生活的圆满境界，他先前也必曾无疑是一个罪人，并从罪的状态皈依至此后新的生命。当我们讨论第二点时，所面对的问题与第四点下面对的问题不同。因为那时我们须考虑的是：是否任何人在今生能藉着天主的恩宠和其意志的热切愿望，达到如此完美而绝对无罪？而现今在第四点中提出的问题是：在人子中间，是否有人，或可能曾有，或将来能有一个，并非是从罪中出来而成全正义，而是从未在任何时候被罪所束缚？因此，若我们关于婴儿所说的话是真实的，那么在人子中间，没有、不曾有、也不会有这样一个人，除了那一位中保，在祂内为我们赚得了赎罪和成义，藉此我们藉着消除由我们罪所产生的敌意而与天主和好。因此，重述一些思考并非不合适——只要当前主题似乎需要——从人类起源开始，以便它们能教导并坚固读者的心智，以回应对他可能困扰的一些异议。\par

% structure: p0071 source=h2 target=subsection reason=relative-heading-rank; base=h2

\subsection{第35章 [XXI.]— \Person{亚当}和\Person{厄娃}；天主最严厉地命令人服从。}

当最初的人——一人亚当，及其妻厄娃，厄娃出于亚当——不愿服从他们从天主所受的命令时，公正且应得的惩罚就临到他们。主曾警告说，他们吃禁果之日必定死。\BibleRef{创 2:17}现在，既然他们获准食用乐园中生长的所有树木（其中包括天主所种的生命树），但被禁止吃唯一的一棵树，那树被称为知善恶树，此名意指他们将发现：若遵守禁令会体验何种善，若违命会体验何种恶。他们无疑被认为在魔鬼恶意引诱之前禁食了禁果，而使用了所有允许的果子，因此，在一切果子中，尤其先吃了生命树。因为还有什么比假设他们吃了其他树的果子却未吃那同样被赐予的、且具有特殊效力的果子更荒谬呢？那树凭借其自身的能力，阻止他们属血肉的身体因岁月的衰败而改变，并因衰老而走向死亡，将此益处从自己的身体传递给人的身体，并以一种奥秘显示智慧（所象征的）赋予理性灵魂什么，即，被其果实赋予生命的灵魂不应改变为不义的衰败和死亡。因为论及智慧时说得对：\BibleQuote{凡紧握她的，她便成为生命树。}\BibleRef{箴 3:18}正如一棵树是为肉体乐园，另一棵是为精神乐园；一棵给予外在感官的活力，另一棵给予内在感官的活力，使其能历经时间而无任何恶化。因此他们事奉天主，因为那忠信的服从托付给他们，唯独借着服从才能敬拜天主。不可能有更合适的方式来暗示服从本身的重要性，或显示服从足以确保受造理性者安然处于造物主之下，如同禁止一棵本身并非邪恶的树那样。因为天主禁止，那造美好之物的造物主，祂创造了万物，\BibleQuote{看，一切都很好。}\BibleRef{创 1:31}祂不会在物质乐园的丰饶中种植任何邪恶的东西。然而，为向人显示服从这样一位主人何等有益——祂只要求仆人服从，这服从益处不在于主人的主权，而在于仆人的得益——他们被禁止用一棵树，若没有禁令，他们用这树不会遭受任何恶果；由此可以清楚明白，他们因不顾禁令使用那树而自招的任何灾祸，都不是由于那树果实的任何有害或恶毒的性质，而完全是由于他们违背了服从。\par

% structure: p0073 source=h2 target=subsection reason=relative-heading-rank; base=h2

\subsection{第三十六章 ［第二十二章］——人堕落前的状况}

在他们这样违命之前，他们悦乐天主，天主也悦乐他们；虽然他们带着属血肉的身体，但在其中并未感到任何不服的私欲对抗自己。这是正义的安排：既然他们的灵魂从主接受了身体作为仆人，正如灵魂服从主一样，身体也应服从灵魂，并毫无抗拒地为所赐的生命提供应有的服务。因此\BibleQuote{两人都赤身露体，并不羞耻。}\BibleRef{创 2:25}理性灵魂现在确实因一种自然的羞耻感而受影响，因为在那具肉体上，它曾获得管辖它的权力，但由于某种难以言说的软弱，它无法在自己不愿时阻止肢体的运动，也无法在自己愿意时激起它们。因此，在每个贞洁的人身上，这些肢体被正确地称为\Quote{\OriginalTerm{可耻之处}{pudenda}}，因为它们随意自身发动，违反作为主人的心意，仿佛自己是主人；而唯有德行的缰绳所拥有的权威是阻止它们接近不洁和非法的玷污。这种肉体的不服，在于这种自身兴奋，即使不允许它发生，在原初男女赤身不羞时并不存在。因为统治肉体的理性灵魂尚未对其主发展出如此不服，以致通过同态报应将自身仆役肉体的反叛连同那种困惑和烦恼的感觉引到自己身上——这种感觉它当然在不服天主时未施加于天主；因为当我们不服从天主时，祂并不感到羞耻或烦恼，我们也无法丝毫减弱祂对我们极大的权能；但我们羞耻于身体不服从我们的管辖——这是因犯罪而得的软弱的后果，被称为\Quote{住在我们肢体中的罪。}但这罪的特性是它乃是罪的惩罚。当那违命被实行，不服的灵魂远离主的律法时，它的仆役身体就开始对它抱有不服从的律法；于是男女二人都因觉察到肉体从前未感到的反叛运动而对自己的赤身感到羞耻，这种觉察被称为\BibleQuote{他们的眼开了；}\BibleRef{创 3:7}当然，他们并非闭着眼走在树木间。同样的话也论及哈加尔说：\BibleQuote{她的眼开了，看见了一口水井。}\BibleRef{创 21:19}于是那男女遮盖了他们的羞处，那本是天主为他们造的肢体，但他们却使之成为羞处。\par

% structure: p0075 source=h2 target=subsection reason=relative-heading-rank; base=h2

\subsection{第三十七章 ［第二十三章］——本性因罪而败坏，因基督而更新}

从罪恶的法律生出罪恶的肉身，这肉身需要藉着那一位以罪恶肉身的形状而来的圣事得以洁净，使罪恶的身体被摧毁，这身体也称为\Quote{死亡的身体}，只有天主的恩宠能藉着我们的主耶稣基督拯救可怜的人脱离它。\BibleRef{罗 7:24} 因为这法律是死亡的根源，从第一对夫妇传给后代，正如所有人在地上劳苦工作、妇女在分娩的疼痛中受产难所见到的那样。这些痛苦是他们在被定罪时，因天主的判决所应得的；我们看到这些痛苦不仅在他们身上，也在他们的后代身上实现，有的多些，有的少些，但总之在所有人身上。然而，最初人类原始义德在于服从天主，不在他们的肢体中有违反心中法律的私欲偏情的法律；现在，由于他们的罪，在我们从他们而生的罪恶肉身中，那些服从天主的人获得了一个巨大的收获，即他们不顺从这邪恶私欲偏情的欲望，反而在自己身上钉死肉身及其私欲和情欲，好使他们属于耶稣基督，祂在祂的十字架上象征了这一点，并藉着祂的恩宠赐给他们能力成为天主的子女。因为不是所有人类，而是凡接受祂的人，祂赐给他们：在由肉身生于世界之后，由圣神重生归于天主。关于这些人确实写着：\Quote{凡接受祂的，祂赐给他们能力成为天主的子女，这些人不是从肉身、也不是从血、也不是从人的意愿、也不是从肉身的意愿，而是从天主生的。}\BibleRef{若 1:12}\par

% structure: p0077 source=h2 target=subsection reason=relative-heading-rank; base=h2

\subsection{第三十八章 [XXIV.]——圣言降生成人赐予了我们什么恩惠；基督在肉身的诞生，何处相似、何处不相似于我们的诞生。}

他接着加上：\Quote{圣言成了肉身，并居住在我们中间}\BibleRef{若 1:14}；这等于说：他们在中间确实成就了一件大事，即他们从肉身生于世界之后，又从天主生于天主，虽然后者是由天主自己创造的；但一件更奇妙的事成就了：虽然他们出于本性是从肉身生的，而藉着天主的仁慈是从天主生的——为将这如此大的恩惠赐予他们，那按自己本性是从天主生的，仁慈地屈尊也从肉身而生——\Quote{圣言成了肉身，并居住在我们中间}这句话所表达的正是这个意思。他实际上说，由此成就了：我们这些从肉身生而为肉身的，后来从圣神而生，可以成为精神并居住在天主内；因为天主，那从天主生的，后来从肉身而生，成了肉身，并居住在我们中间。因为那成了肉身的圣言，在起初就存在，并与天主同在。\BibleRef{若 1:1} 但同时，祂参与我们较低的状态，好使我们参与祂较高的状态，在祂的肉身诞生中保持了一种中道：我们确实是在罪恶的肉身中生的，但祂是在罪恶肉身的形状中生的；我们不仅从肉和血，也从人的意愿和肉身的意愿而生，但祂只从肉和血而生，不是从人的意愿，也不是从肉身的意愿，而是从天主生的；因此，我们因罪而死，祂无因罪而为我们而死。同样，正如祂降到我们这里时所取的低等情况下，并不在每一点上都完全等同于我们被祂在这里发现时我们的低等情况；同样，我们升到祂那里时的高等状态，也不会完全等同于祂在那里我们找到祂时的祂的高等状态。因为我们藉着祂的恩宠将要成为天主的子女，而祂从永远就是天主子；我们归依后要依附于天主，但并非与祂平等；祂从未离开天主，永远与天主平等；我们分享永生，祂是永生。因此，唯有祂成了人，但仍然是天主，从未有任何罪，也没有取一个罪恶的肉身，尽管是从一个罪恶的母亲的肉身出生的。因为祂那时从肉身所取的，祂要么洁净了才取，要么因取了而洁净。因此，祂的童贞母亲，她的受孕不是按罪恶肉身的法律（即不是靠肉身的私欲偏情的激动），而是因她的信德应得那圣洁的种子在她内成形，祂塑造了她为了拣选她，拣选了她为了从她成形。那么，当那未受罪玷污的肉身受洗是为了树立效法的榜样时，罪恶的肉身岂不更需要受洗以逃脱审判吗？\par

% structure: p0079 source=h2 target=subsection reason=relative-heading-rank; base=h2

\subsection{第三十九章 [XXV.]——佩拉纠派的一个异议。}

我们已对那说\Quote{如果罪人生了罪人，那么义人也应该生义人}的人所作的答复，现在用来回答那些主张由受洗者所生的人应被视为已受洗的人。因为他们问：\Quote{为什么他不能在父亲的腰中受洗呢？因为按\BibleBook{}{希伯来书}，肋未能在亚巴郎的腰中献了什一税？}提出这个论点的人应当注意：肋未后来并未因此就不再缴纳什一税，因为他已经在亚巴郎的腰中缴纳了什一税，而是因为他被任命为司铎的职务以收取什一税，不是缴纳什一税；否则他的兄弟们，他们都向他缴纳什一税，也不会被什一税——因为他们在亚巴郎的腰中时也已经向默基瑟德缴纳了什一税。\par

% structure: p0081 source=h2 target=subsection reason=relative-heading-rank; base=h2

\subsection{第四十章。——一个预料的论点。}

但愿无人争辩说，亚巴郎的后代已在他们祖先的腰中缴纳了什一之税，因而可以不必再缴纳，因为缴纳什一之税是一种义务，需要每人不断地重复履行，就如以色列人每年终身向他们的肋未人缴纳各种出产的什一之税一样；而圣洗圣事却是一种只能领受一次的圣事，如果一个人在父亲怀中已经领受了，那么他只能被认为是已经受洗的人，因为他由一位已受洗的人生出。那么，谁若如此争辩（我只简单地说，而不详细讨论），他应当看割损礼，那是只施行一次的，却仍须对每个人分别施行。因此，正如在那古老圣事的时代，受割损者的儿子也必须受割损，同样，如今已受洗者的儿子也必须领受圣洗。\par

% structure: p0083 source=h2 target=subsection reason=relative-heading-rank; base=h2

\subsection{第四十一章——宗徒称信友的子女为\Quote{圣洁的}。}

宗徒的确说：\Quote{\Emph{不然}，你们的子女便\Emph{不洁}，但\Emph{如今}他们是\Emph{圣洁的}；\BibleRef{格前 7:14} }因此他们推断\Quote{信友的子女不必受洗。}我甚为惊讶，那些否认原罪由亚当传下来的人竟会使用这样的话。因为如果他们把这句宗徒的话理解为信友的子女生来就是圣洁的，那么他们自己为何对受洗的必要性毫无疑虑？最后，他们为何拒绝承认从有罪的父母会传来原罪，如果从圣洁的父母会传来某种圣洁？实际上，即使从信者会生出\Quote{圣洁}的子女，这也不反驳我们的主张，因为我们坚持除非他们受洗，否则将进入永罚，而我们的对手自己就是关闭了天国之门给这些人，尽管他们坚持这些人没有罪，无论是本罪还是原罪。或者，如果他们以为\Quote{圣洁的人}受罚是不合适的，那么将\Quote{圣洁的人}排除在天国之外又如何合适呢？他们更应特别注意这一点：如果从圣洁的父母传来某种圣洁，从不洁的父母传来不洁，那么从有罪的父母怎能不传来有罪？因为当他说道\Quote{不然，你们的子女便不洁，但如今他们是圣洁的}时，就肯定了这种双重原则。他们还应向我们解释，为何信友的圣洁子女与不信者的不洁子女，尽管情况不同，若未受洗，同样被禁止进入天国？那圣洁对他们又有何用？如果他们主张不信者的不洁子女要受罚，但信友的圣洁子女虽未受洗不能进入天国，但并非因此受罚，因为他们是\Quote{圣洁的}，那倒是一种区分；但事实上，他们同样宣称：圣洁父母的圣洁子女和不洁父母的不洁子女，都因为没有罪而不受罚，并且都因为未受洗而被排除在天国之外。何其荒谬！谁能想见如此卓越的天才们竟看不到这一点？\par

% structure: p0085 source=h2 target=subsection reason=relative-heading-rank; base=h2

\subsection{第四十二章——圣化的多重方式；慕道者的圣事。}

我们在这点上的意见与宗徒本人完全一致，他曾说：\Quote{由一人而来，众人都被判定}，以及\Quote{由一人而来，众人都获得生命的成义}。如今，当他在另一处论及其他要点时所说的话：\Quote{不然，你们的子女便不洁，但如今他们是圣洁的}，与这些陈述如何协调，请思考片刻。【XXVI.】圣化并非只有一种尺度；因为我认为，连慕道者也在他们自己的尺度上因基督的标记和覆手祈祷而被圣化；他们所领受的是神圣的，尽管那不是基督的身体——却比任何构成我们日常饮食的食粮更为神圣，因为它是圣事。然而，那维持我们现世生命所需的饮食，按同一位宗徒所说，\Quote{藉着天主的话和祈祷而被圣化}\BibleRef{弟前 4:5}，即我们祈求身体获得滋养的祈祷。因此，正如日常食物的这种圣化并不阻止进入口中的东西下到肚腹，又排泄到茅厕\BibleRef{谷 7:19}，并参与一切尘世之物的腐朽，主因而劝勉我们为那永不朽坏的另一食粮而劳碌\BibleRef{若 6:27}：同样，慕道者的圣化，若他未受洗，对于他进入天国或获得罪赦也毫无益处。同理，那圣化，不论其尺度如何，按宗徒所言存在于信友子女中，也与圣洗圣事及罪的起源或赦免毫无关系。宗徒在这段我们所关注的经文中说，一个不信的配偶因信主的配偶而被圣化：\Quote{因为不信的丈夫因妻子而成了圣洁的，不信的妻子因丈夫而成了圣洁的；不然，你们的子女便不洁，但如今他们是圣洁的。}\BibleRef{格前 7:14}我敢说，没有一个人的心思因不信而扭曲到这种地步，以致无论他对这些话作何解释，竟会认为一位未信主的丈夫不必受洗，因为他的妻子是基督徒，并且他已经获得了罪赦，有进入天国的确定希望，因为他被描述为由妻子而圣化。\par

% structure: p0087 source=h2 target=subsection reason=relative-heading-rank; base=h2

\subsection{第四十三章【XXVII.】——为何已受洗者的子女应当受洗。}

若有人仍对为何已受洗者的子女需要受洗感到困惑，愿他简要思量：正如借着一个人——亚当，有罪的肉身所生的后代都被带入定罪，同样，借着一个人——耶稣基督，恩宠之神的后代将一切因预定而分享此重生者带入永生的成义。然而，圣洗圣事无疑就是重生的圣事。因此，正如从未活过的人不能死，从未死过的人不能复活，同样，从未出生的人不能重生。由此得出结论：没有人能在其父内重生，除非他已出生。然而，人在出生之后必须重生，因为\Quote{人除非重生，不能见到天主的国}\BibleRef{若 3:3}。因此，即使婴儿也必须领受重生的圣事，否则他将不幸地离开此生；而这圣洗圣事除为赦罪外，不另举行。基督在此段经文中也明确指示我们：当人问\Quote{这事怎能成就？}时，祂提醒问者梅瑟在举蛇时所行的。如此，既然婴儿借着圣洗圣事与基督的死亡相似，就必须承认他们也脱离了蛇的毒咬，除非我们故意偏离基督信仰的规则。然而这毒咬并非在他们自己的实际生命中领受，而是在那最初受伤者身上。\par

% structure: p0089 source=h2 target=subsection reason=relative-heading-rank; base=h2

\subsection{第四十四章　佩拉热异端者的一个异议}

他们也看到这一点：父母在皈依之后，自己的罪不再有害于他们；因此他们提出疑问：\Quote{既然如此，这些罪又怎能妨碍他的儿子呢？}但如此想的人没有注意到：正如父亲自己的罪之所以不伤害他，正是因为他由圣神重生，同样，在他儿子身上，除非他以同样的方式重生，他从父亲所继承的罪将对他有害。因为即使更新的父母生育子女，也不是出于他们更新状态的首果，而是出于旧本性的残余，以血肉之躯生育；这些子女作为父母残留旧本性的后代，生在罪身之中，借着属灵重生与更新的圣事逃脱了属于旧人的定罪。这是一个我们应当铭记在心的考量，因为围绕着这个问题已经发生并可能继续发生争论：在圣洗圣事中确实发生了罪过的完全而完美的赦免，但实际之人的品性并不会立即完全改变；只有在那些相称行走的人中，圣神的首果借着日复一日增长的更新过程，将旧的肉性改变为同样品性，直至整个旧本性如此更新，以致自然身体的软弱达到属灵身体的刚强与不朽坏。\par

% structure: p0091 source=h2 target=subsection reason=relative-heading-rank; base=h2

\subsection{第四十五章［XXVIII］——罪的律法被称为罪；在已受洗者中，私欲偏情的邪恶被移除后仍如何存留}

然而，这罪的律法——宗徒也称之为\Quote{罪}，他说：\Quote{所以不要让\Emph{罪}在你们必死的身体上为王，使你们顺从它的私欲}\BibleRef{罗 6:12}——在那些由水和圣神重生的人身上，并非以其未得赦免的方式存留在肢体中，因为我们的罪有着完全而完美的赦免，一切使我们与天主分离的仇敌已被杀灭；但它存留在我们旧的肉性中，仿佛已被克服和摧毁，除非它因同意于非法之事而某种程度地复苏，并恢复其统治与主权。然而，在这旧的肉性——其中罪的律法或罪已被废除——与圣神之生命——借着这种新生命，受洗者因天主的恩宠重生——之间，可看到如此明显的区别，以致宗徒认为只是说他们不在罪中还不够，除非他也说他们不在肉体本身之中，甚至在他们离开这必死生命之前。他说：\Quote{随从肉体的人不能得天主的欢心；但你们\Emph{不在肉体内}，而在于圣神，只要天主的圣神住在你们内。}\BibleRef{罗 8:8}事实上，正如那些将肢体的能力用于善工、不再处于那肉体中的人——因为他们不再按肉体原则塑造理解或生活——甚至将肉体本身（虽然它是可朽坏的）用于善工；同样，他们也善用死亡（即原罪的惩罚），为兄弟、为信仰、为一切真、圣、义之事勇敢而坚韧地面对死亡——同样，在信仰中所有\Quote{真正的同轭者}也善用那仍存留在已得赦免的旧肉性中的罪的律法本身，因为他们拥有在基督内的新生命，不容私欲统治他们。然而，正是这些人，因为他们仍带着亚当的旧本性，以必死的方式生育子女，使这些子女得以不朽地重生，并带有罪的传递；在罪中，那些重生之人不再被束缚，而出生者因重生而得释放。因此，只要这律借着私欲偏情居住在肢体中，虽然它仍存留，它的罪责已被赦免；但只有领受了重生圣事并已开始更新的人，这赦免才生效。凡从仍存留其私欲偏情的旧本性所生的，都需要重生才能得医治。既然相信的父母（他们既以肉身出生，又已属灵重生）自己以肉身的方式生育了子女，那么他们的子女在第一次出生之前，又怎能已经重生呢？\par

% structure: p0093 source=h2 target=subsection reason=relative-heading-rank; base=h2

\subsection{第四十六章——罪责可被除去，但私欲偏情仍存留}

你不可因我所说的话感到惊奇，即：虽然罪的律法连同其私欲偏情仍然存在，但其罪债却透过圣事的恩宠被消除了。因为正如邪恶的行为、言语和思想，就其心灵和身体的纯粹活动而言，已经过去并停止存在，然而它们在过去之后、不再存在时，其罪债仍然存留，除非藉着罪的赦免而消除；同样，相反地，在这尚未消除而仍存留的私欲偏情的律法中，其罪债被消除，不再存留，因为在圣洗圣事中发生了完全的赦罪。事实上，如果一个人在领洗之后立即离开此生，那么他完全没有留下任何东西可让他承担责任，因为所有使他承担责任的东西都已被解除。因此，一方面，过去的思想、言语和行为的罪债在赦免之前仍然存在，这毫不奇怪；另一方面，存留的私欲偏情的罪债在赦免之后消失，这也无需令人惊讶。\par

% structure: p0095 source=h2 target=subsection reason=relative-heading-rank; base=h2

\subsection{[XXIX.]——所有预定者皆藉唯一中保基督并藉同一信德得救。}

既然如此，自从罪藉着一人进入了这世界，死亡藉着罪进入了世界，这样死亡就传遍了众人，直到这血肉生育和败坏世界的终结——这世界的儿女们生养并受生——从未有过，也永远不会有一个处于我们这种生活中的人，能说他是毫无罪的，除了那唯一的中保，祂藉着罪的赦免使我们与我们的创造主和好。现在，这位我们的主，在人类历史的任何时期，从未拒绝，也直到最后审判，永远不会拒绝祂的医治，给予那些祂在极其确定的预知和未来的慈爱中预定与祂一同为王、得到永生的人。因为，在祂降生成人、在苦难中软弱、在祂自己的复活中显大能之前，祂教导所有当时活着的人，信靠那些当时\Emph{未来}福分的信德，使他们能继承永生；而当所有这些事在基督身上成就、并目睹预言应验时活着的人，祂教导他们信靠这些当时\Emph{现在}的福分；至于此后活着的人、我们现在活着的人、以及所有将来要活着的人，祂并未停止教导他们信靠这些现已\Emph{过去}的福分。因此，是\Quote{同一信德}拯救所有在肉身的出生之后由圣神重生的人，而它归结于祂，祂来为我们受审而死——审判生者死者之主。但这\Quote{同一信德}的圣事，为了适切的象征，在不同时期有所不同。\par

% structure: p0097 source=h2 target=subsection reason=relative-heading-rank; base=h2

\subsection{第四十八章。——基督甚至是婴儿的救主；基督在婴儿时，免于无知和精神软弱。}

祂因此同时是婴儿和成年人的救主，关于祂天使曾说：\Quote{今天为你们诞生了一位救主；}\BibleRef{路 2:11} 又关于祂向童贞玛利亚宣告说：\Quote{你要给祂起名叫耶稣，因为祂要把自己的民族从他们的罪恶中拯救出来，} 这里清楚表明祂被称为耶稣，是因为祂赐予我们的救恩——耶稣相当于拉丁文的\OriginalTerm{救主}{Salvator}，\Quote{救主}。那么，谁敢如此大胆，主张主基督只是成年人的耶稣，而不是婴儿的耶稣呢？祂取了有罪肉身的模样，为要摧毁罪恶的身体，却带着婴儿的肢体，适合于那身体极度软弱中的无用，而祂理性的灵魂被悲惨的无知所压迫！我不能相信在道成为肉身、居住在我们中间的婴儿中，存在这种全然无知；我也不能想象在婴儿基督中，存在我们在一般婴儿身上所见的那种心智官能的软弱。因为由于这种软弱和无知，婴儿被非理性的情感所困扰，且不受任何理性的命令或管束所约束，而是藉着痛苦和惩罚，或对这些的恐惧；所以你完全可以看见，他们是那悖逆的儿女，这悖逆在我们身体的肢体中激动，违背心灵的律法——即使理智愿意，它也不肯安静；反而常常只能藉着实际施加身体上的痛苦（例如用鞭打）来压制；或只能藉着恐惧或某种心灵的情感来抑制，而非藉着意志的任何劝诫。然而，因为祂身上有有罪肉身的模样，祂愿意经历各个年龄阶段的变化，甚至从婴儿期开始，这样看起来，如果祂不是在年轻时被杀，祂的血肉甚至可能因衰老的逐渐逼近而到达死亡。不过，死亡是因悖逆而加于有罪肉身的刑罚，但在有罪肉身的形像中，祂却出于自愿的服从而承受了死亡。因为当祂走向死亡、即将受苦时，祂说：\Quote{看，这世界的首领来了，他在我身上一无所有。但为使众人知道我是遵行我父的旨意，起来，我们从这里走吧。}\BibleRef{若 14:30} 祂说了这些话后，便立刻前往，承受了祂不应得的死亡，成为服从至死。\par

% structure: p0099 source=h2 target=subsection reason=relative-heading-rank; base=h2

\subsection{[XXX.]——白拉奇派的一个异议。}

因此，那些说：\Quote{如果因第一个人的罪，我们必须死；那么因基督的来临，相信祂的人就应当不死；}并且他们加上一个自认为的理由，说：\Quote{因为第一个犯罪者的罪，损害我们，不可能大于救主的降生成人或救赎对我们的益处。}然而，为何他们不更加留心倾听，并毫不迟疑地相信宗徒如此明确陈述的话：\BibleQuote{因为死亡既由一人而来，死者的复活也由一人而来；就如在亚当内，众人都死了，照样，在基督内，众人都要复活；}？\BibleRef{格前 15:21}宗徒所说的，无非就是身体的复活。他说因一人，众人身体的死亡临到，然后加上应许：因一人，即基督，众人身体复活至永生将要发生。怎能说\Quote{一人因犯罪损害我们多于另一人因救赎裨益我们}呢？因为因前者的罪，我们经历暂时的死亡；但因后者的救赎，我们复活不是到暂时的，而是到永远的生命。因此，我们的身体因罪而死亡，但基督的身体只死了而无罪，为要无过犯地倾流祂的血，\Quote{涂抹}那包含一切过犯记录的\Quote{债券}，这些债券使如今相信祂的人从前成为魔鬼的负债者。因此祂说：\BibleQuote{这是我的血，为许多人流出，以赦免罪过。}\BibleRef{玛 26:28}\par

% structure: p0101 source=h2 target=subsection reason=relative-heading-rank; base=h2

\subsection{第五十章【XXXI.】——为何在圣洗中与罪一同被取消的，死亡本身却不被取消。}

然而，祂本也可以将此恩赐给信徒，使他们甚至连身体的死亡也不经历。但如果祂这样做了，无疑会给肉体增添某种福乐，但信德的刚毅会减少；因为人们如此畏惧死亡，以致他们会仅仅因免于死亡而称基督徒为有福。没有人会为了死后极其幸福的生命，借着藐视死亡本身的力量，而奔赴基督的恩宠；反而为了免除死亡的痛苦，他们会倾向于一种更舒适的相信基督的方式。因此，祂赐予那些相信祂的人比这更大的恩宠，无疑给了他们一份更伟大的礼物！一个人看到别人成为信徒后不死，自己也相信他将不死，这有什么了不起呢？多么更伟大的事，多么更勇敢、更值得称颂的事：虽然确定要死，仍能希望永远活着！最终，在末日，有些人将蒙受此恩赐，在突然的变化中不再感受死亡本身，而是与被复活的人一同被提到云中，在空中迎接基督，如此他们永远与主同在。这些领受此恩宠的人理应如此，因为在他们之后将没有后代需要被引导去相信——不是借着未见之事的希望，而是借着所见之事的爱。这样的信德是软弱无力的，根本不能称为信德，因为信德如此定义：\BibleQuote{信德是所希望之事的担保，是未见之事的确证。}\BibleRef{希 11:1}因此，在《希伯来书》同一书信中，这段话之后，在接下来的句子中列举了一些凭信德悦纳天主的古人，他说：\BibleQuote{这些人都怀着信德死了，没有获得所应许的，只从远处望见，且表示欢迎，明认自己在世上只是客旅和寄居的。}\BibleRef{希 11:13}然后后来他用这些话结束对信德的赞辞：\BibleQuote{这些人虽然因信德获得了褒扬，却没有领受到所恩许的，因为天主为我们预备了更好的事，以致他们没有我们，也不会成为完美的。}\BibleRef{希 11:39}如果人们相信时追求看得见的报偿——也就是说，如果信徒在今世就得到不死的赏报，那么这就不是对信德的赞美，也不是，正如我所说的，信德了。\par

% structure: p0103 source=h2 target=subsection reason=relative-heading-rank; base=h2

\subsection{第五十一章——为何说魔鬼掌有死亡的权利和统治。}

因此，主自己愿意死亡，\Quote{为能，}如经上论祂所记：\BibleQuote{借着死亡，毁灭那握有死亡权势的，即魔鬼；并解救那些因死亡的恐惧，一生处于奴役之中的人。}\BibleRef{希 2:14}这段话充分表明，甚至身体的死亡也是因魔鬼的煽动和作为——简言之，因祂引诱人所犯的罪而来；没有其他理由可以说他严格地握有死亡的权势。因此，那位没有任何原罪或本罪而死的，在我已引用过的那段话中说：\BibleQuote{看，这世界的首领}，即握有死亡权势的魔鬼，\BibleQuote{来了，在我身上一无所有}——意思是，他找不到我有什么罪，因这些罪他使人死亡。仿佛有人问他：那你为何要死？祂说：\BibleQuote{为叫世人知道我爱父，并且父怎样命令我，我就怎样行；起来，我们从这里走罢；}\BibleRef{若 14:30}\par

% structure: p0105 source=h2 target=subsection reason=relative-heading-rank; base=h2

\subsection{第五十二章【XXXII.】——为何基督复活后，从世界撤回祂的临在。}

因此，虽然主行了许多可见的奇迹，为使信德起初萌发并依靠婴儿般的滋养而成长，并逐渐从这种软弱中达到全备的力量（因为信德越坚强，寻求这种帮助就越少）；祂却愿意我们安静地等待，在没有可见诱导的情况下，盼望所应许的恩许，以便\Quote{义人因信德而生活；}\BibleRef{哈 2:4}而且祂的这种愿望如此强烈，以至于祂虽然第三天从死者中复活，却不愿留在人中间，而是在留下祂复活的证据——向那些祂屈尊作为此事件的见证人显现肉身之后，祂升了天，这样从他们的视线中隐退，并未给他们任何人的肉身赐予像祂在自己的肉身中所显示的那样，好叫他们也\Quote{因信德而生活，}并在现今世界中耐心等待，没有可见的诱导，等待那使人因信德而生活的义德的赏报——这赏报将在日后可见地、公开地赐予。我认为，祂论及圣神所说的这段话必须指向这个意义：\Quote{我若不离去，祂就不会来。}\BibleRef{若 16:7}因为这实际上等于说：你们将无法因信德而正义地生活——这信德将作为我的恩赐，即来自圣神——除非我从你们的眼前离去，你们如今所注视的那位，好使你们的心通过将信德集中于不可见的事物上而获得灵性的成长。祂不断地向门徒们推荐这种信德的义德。论及圣神，祂说：\Quote{祂要指证世界关于罪、关于义德、关于审判：关于罪，因为他们没有信从我；关于义德，因为我往父那里去，而你们将不再看见我。}\BibleRef{若 16:8}那种义德是什么——人们因之不再看见祂——岂不正是\Quote{义人因信德而生活，}并且我们不注视那可见的，而注视那不可见的，在圣神内等待那因信德而来的义德的希望吗？\par

% structure: p0107 source=h2 target=subsection reason=relative-heading-rank; base=h2

\subsection{第五十三章【三十三】——佩拉纠派的异议。}

但那些说\Quote{如果身体的死亡是因罪而来的，那么在我们罪过的赦免——这是救赎主所赐予我们的——之后，我们当然不该死亡，}的人，并不明白为何有些事物，其罪债已被天主勾销，以免在今生之后阻碍我们，祂却允许它们存留，作为信德的争斗，以便成为那些在圣德争战中前进者的教训和操练。也许有人因不理解这一点而提出疑问：如果天主因人的罪而对人说：\Quote{你必汗流满面，才有饭吃；地要给你生出荆棘和蒺藜，}\BibleRef{创 3:18}那么为什么在罪过赦免之后，这种劳动和辛劳仍然继续，而且信徒的土地仍然给他们这种粗糙而可怕的收成？再者，既然因妇女的罪而对她说：\Quote{你应在痛苦中生子，}\BibleRef{创 3:16}为什么信主的妇女在罪过赦免之后，在分娩过程中仍然遭受同样的痛苦？然而，这是无可争议的事实：由于他们所犯的罪，原初的男人和女人从天主那里听到了这些判决，并且应得这些判决；也无人反对我所引用的关于男人劳动和女人分娩之痛苦的圣经话语，除非有人完全敌对公教信仰，并敌视默感写作的圣经。\par

% structure: p0109 source=h2 target=subsection reason=relative-heading-rank; base=h2

\subsection{第五十四章【三十四】——为什么罪过已被赦免，惩罚却仍被施加。}

但是，既然不乏这样的人，正如我们答复那些提出此问题的人时所说，那些在赦免之前是罪的惩罚的事物，在赦免之后成为义人的争斗和操练；同样，对于因身体的死亡而同样困惑的人，我们的答复也应如此提出，以表明我们既承认死亡是因罪而来的，也并未因惩罚被留给我们作为修练的操练而气馁，好让我们在圣德上进时能战胜对死亡的巨大恐惧。因为如果只有小小的德行归于\Quote{以爱德行事的信德}去征服对死亡的恐惧，那么殉道者便不会有大荣耀；主也不能说：\Quote{人若为自己的朋友舍命，再没有比这更大的爱了；}\BibleRef{若 15:13}若望在他的书信中这样表达：\Quote{祂为我们舍命，我们也应当为弟兄舍命。}\BibleRef{若一 3:16}因此，如果死亡本身中没有真实而非常严酷的考验，那么对因义德而忍受或藐视死亡的最卓越受苦的赞美便是徒然的。一个人藉信德战胜了对死亡的恐惧，便为信德本身赢得了极大的荣耀和公义的赏报。因此，要么身体的死亡在人犯罪之前不可能发生——因为它是作为罪的惩罚——要么在罪过赦免之后它临到信徒，为的是在他们战胜对死亡的恐惧时，能操练义德的刚毅，这一点都不该令人惊讶。\par

% structure: p0111 source=h2 target=subsection reason=relative-heading-rank; base=h2

\subsection{第五十五章——要恢复因罪而失去的义德，人必须用大量的劳苦和忧伤来奋斗。}

起初受造的肉身并非那有罪的肉身，人在乐园的安乐中未能持守其义德，因此天主定意，使有罪的肉身在犯罪后自行繁衍，并在诸多劳苦与磨难中为恢复圣德而挣扎。因此，亚当被逐出乐园后，不得不居住在伊甸对面——即乐园的对面——以表明，有罪的肉身在最初的欢愉中未能持守圣德、沦落为有罪的肉身之后，必须通过劳苦与忧患——这些正是欢乐的反面——来磨炼。同样，我们的原祖父母虽后来悔改度义德生活，并借主的血得以免除判决中最重的罚，但他们在世时仍不配被召回乐园。照样，我们的有罪肉身，即使人在罪赦后在此肉身中度义德生活，也不配立刻免除从罪之繁衍而来的死亡。\par

% structure: p0113 source=h2 target=subsection reason=relative-heading-rank; base=h2

\subsection{第五十六章——以达味为例证}

关于达味圣王，我们在\WorkTitle{列王纪}中读到类似的想法。先知受派到他那里，以天主因他犯的罪而发怒的灾祸恐吓他；他认罪后获得了宽恕，先知回答说他的羞耻和罪行已被赦免。然而，天主恐吓他的灾祸仍然按部就班地来临，他因儿子而受屈。此时，人们为何不立即反对说：\Quote{如果天主因他的罪而恐吓他，为何在罪赦后仍然实现他的恐吓？}除非——若有人挑剔——最正确的回答是：罪赦是为了使人不被阻碍获得永生，但恐吓的灾祸仍被执行，以使人在降卑的处境中锻炼并证实其虔敬。同样，天主因人犯罪而使人遭受身体的死亡，但在罪赦后并不赦免死亡，好让人在义德中受锻炼。\par

% structure: p0115 source=h2 target=subsection reason=relative-heading-rank; base=h2

\subsection{第五十七章[三十五]——不偏左右}

因此，让我们坚定不移、毫不犹疑地持守这信仰的宣认。唯有那一位，生于有罪肉身的肖像中，无原罪，在别人的罪中活出无罪，并因我们的罪而死而无罪。\BibleQuote{不可偏右偏左。}\BibleRef{箴 4:27}因为偏右是自欺，声称自己无罪；偏左则是以某种我不知道的、多么悖逆而败坏的放纵，放任自己于罪中。\Quote{天主确实知道右手的道路，}祂是唯一无罪、且能抹去我们罪过的那一位；\Quote{但左手的道路是乖僻的，}与罪为友。那种坚定不移的品格，正是那些二十岁的青年所预示的，他们预表了天主的子民；他们进入了恩许之地，据说\BibleQuote{不偏右也不偏左。}\BibleRef{苏 23:6}这个二十的年龄不可与孩童的无辜相比，但若我未弄错，这数字是奥秘的影子和回响。因为旧约在\WorkTitle{梅瑟五书}中卓越非常，新约则在\WorkTitle{四福音}的权威中最灿烂。这些数字相乘得到二十：四乘以五，或五乘以四，等于二十。这样的子民——如我所说——受两个约（旧约与新约）中天国的教导，既不偏右地骄傲自恃义德，也不偏左地放荡以罪为乐，将进入恩许之地，在那里我们不再祈祷赦免罪过，也不害怕它们在我们身上受罚，因为我们都已藉着那救赎者从一切罪中得释放。祂并非\BibleQuote{被卖于罪恶之下}\BibleRef{罗 7:14}，却\BibleQuote{救赎以色列脱离他的一切过犯。}\par

% structure: p0117 source=h2 target=subsection reason=relative-heading-rank; base=h2

\subsection{第五十八章[三十六]——\BibleQuote{有罪肉身的肖像}暗示实在}

那些人不愿公开承认婴孩需要罪的赦免，却承认他们需要救赎，这对圣经的权威和真实是不小的让步。他们所说的话与源自基督徒训导的另一种说法并无不同。而那些忠信诵读、忠信聆听、忠信持守圣经的人，毫不怀疑：从最初因选择犯罪而成为有罪的肉身，并随后代代相传的肉身，传衍了有罪的肉身——唯有那\BibleQuote{有罪肉身的肖像}\BibleRef{罗 8:3}是例外。然而，这\Emph{肖像}之所以存在，是因为也有有罪肉身的\Emph{实在}。\par

% structure: p0119 source=h2 target=subsection reason=relative-heading-rank; base=h2

\subsection{第五十九章——灵魂是否传衍；关于圣经未启示之奥秘，我们必须谨防仓促判断与意见；圣经有关救恩之必要课题已足够清晰}

关于\Emph{灵魂}，确实有这样的问题：它是否也以同样的方式[像肉体那样]传下来，并被同样的罪债所束缚（这罪债被赦免）——因为我们不能说，只有婴儿的肉体需要救主和救赎者帮助，而他的灵魂不需要，或者说，灵魂不得包含在那首圣咏的感恩中（我们读到并重复）：\BibleQuote{我的灵魂，请赞美上主，不要忘记他的一切恩惠；是他赦免了你的一切罪恶，是他治愈了你的一切疾病，是他救赎了你的性命脱离死亡。} 如果灵魂不是同样传下来的，我们可以问：是否仅仅因为灵魂与有罪的肉体混合并被其压制，仍然需要赦免它自己的罪，需要它自己的救赎，天主作为审判者，以祂预知的高度，知道什么婴儿在出生之前，或在任何情况下曾行过善或恶，不应被赦免那罪债？问题还在于：天主（即便祂不通过自然生殖创造灵魂）又如何能不是那罪债的创造者——由于这罪债，婴儿的灵魂需要圣事的救赎？这是一个广泛而重要的问题，需要另一篇论文。然而，据我判断，讨论应当以温和与适度进行，以便配得上谨慎探究的称赞，而不是固执主张的谴责。因为每当一个问题出现在异常晦涩的论题上，而圣经的清晰而确定的证据不能提供帮助时，人的臆断应当自我克制，也不应偏向任何一面试图得出确定结论。但如果我确实对这些点的某些方面无知——关于它们如何能被解释和证明——我仍应相信这一点：恰恰从这种情况中，圣经将拥有最清晰的权威，每当出现一个问题，人若不了解它就会危及已应许给他的救恩时。现在，[我亲爱的\Person{马尔塞利诺}，]你面前有这篇论文，我尽最大努力完成。我只希望它的价值与它的长度相称；因为它的长度我或许能够辩解，只是我担心，再加上辩解，我会使冗长超过你的忍受限度。\par

\SourceRecord{Translated by Peter Holmes and Robert Ernest Wallis, and revised by Benjamin B. Warfield.；Nicene and Post-Nicene Fathers, First Series；Vol. 5.；Edited by Philip Schaff.；Buffalo, NY: Christian Literature Publishing Co.,；1887.；<http://www.newadvent.org/fathers/15012.htm>.}

% structure: p0001 source=h1 target=section reason=page-title

\section{论功绩与罪的赦免及婴儿的圣洗（第三卷）}

\Emph{以书信形式致同一马赛利努斯。}\par

\Emph{奥古斯丁在此信中驳斥了佩拉吉乌斯关于罪的功绩和婴儿圣洗问题上的若干错误——佩拉吉乌斯曾将他的各种论点穿插在他对圣保禄的注释中，以反对原罪。}\par

致他心爱的儿子马赛利努斯，奥古斯丁，主教、基督的仆人和基督众仆人的仆人，在主内问安。\par

% structure: p0005 source=h2 target=subsection reason=relative-heading-rank; base=h2

\subsection{第一章 [I.]——佩拉吉乌斯被尊为圣人；他对圣保禄的注释}

你曾提议要我写信答复的那些问题——反对那些说即使亚当没有犯罪也会死、并且亚当的罪没有通过自然遗传传给后代的人；特别是关于婴儿的圣洗（普世教会以最虔诚和慈母般的关怀坚持不断施行）；以及今生是否有、曾有过、或将会有没有任何罪的人之子——我已经在两卷长篇论著中讨论过了。我冒昧地认为，如果在这两卷书中我没有解答所有在这些问题上困惑人心的要点（我认为——不，我确信——这既超出我，也超出任何其他人的能力），至少我已经准备了一些坚固基础，使那些捍卫我们祖先所传信仰、反对其对手新观点的人，随时可以站稳立场，不至于在论战中手无寸铁。然而，在最近几天里，我读到了佩拉吉乌斯的一些著作——据我所知他是一位圣洁的人，在基督徒生活中取得了不小的进步——其中包含一些关于宗徒保禄书信的非常简短的注释；当我读到宗徒所说的那段话：\Quote{因一人罪进入世界，罪又带来死亡；这样死亡就传给了众人，} \BibleRef{罗 5:12} 时，我发现了一个被那些说婴儿不带有原罪的人所采用的论点。我承认，我在我的长篇论著中没有反驳这个论点，因为它确实从未出现在我的脑海中，以至于我认为有人会如此思考。然而，由于我不愿意在那部已经完成的作品上再增加内容，我认为应当在这封信中插入这个论点本身（以我读到的原话）以及我认为适当的答复。\par

% structure: p0007 source=h2 target=subsection reason=relative-heading-rank; base=h2

\subsection{第二章 [II.]——佩拉吉乌斯的异议；婴儿被列在信者和信徒之中。}

于是，论点是这样陈述的：——\Quote{那些否认罪的遗传的人试图这样反驳：如果（他们说）亚当的罪伤害了那些不犯罪的人，那么基督的义也同样有益于那些不信的人；因为\InnerQuote{同样，不仅如此，更是如此}，他说，\InnerQuote{因一人而得救的人，比先前因一人而灭亡的人更多。}}对于这个论点，我再说一遍，我在先前写给您的两卷书中没有进行答复；事实上，我并未给自己设定这样的任务。但现在我请您首先注意，当他们说：\Quote{如果亚当的罪伤害了那些不犯罪的人，那么基督的义也同样有益于那些不信的人}，他们如何认为这是荒谬和虚假的，即基督的义应该有益于那些不信的人；然后他们试图由此构造这样的论点：即第一人的罪不可能伤害不犯罪的婴儿，正如基督的义不能有益于任何不信的人一样。因此，请他们告诉我们，基督的义对受洗的婴儿有什么益处；请他们务必告诉我们他们是什么意思。当然，既然他们不忘记自己也是基督徒，他们毫不怀疑有某种益处。但无论这益处是什么，它（正如他们自己所声称的）不能有益于那些不信的人。因此，他们被迫将受洗的婴儿列在信者之中，并赞同神圣普世教会的权威，该教会不认为那些因他们而得基督之义（按照他们的说法，只有信者才能得此义）的人不配称为信者。因此，正如通过那些使他们重生的人的答复，义德之神将信德传递给他们，而他们自己的意志还不能拥有这信德；同样，通过那些使他们出生的人，有罪的肉体将伤害传递给他们，而他们在自己的生命中尚未犯下这伤害。正如生命之神在基督内重生他们为信者，同样死亡之体在亚当内生他们为罪人。一次出生是肉体的，另一次是属神的；一次产生肉体之子，另一次产生精神之子；一次是死亡之子，另一次是复活之子；一次是世界之子，另一次是天主之子；一次是义怒之子，另一次是仁慈之子；因此，一次将他们束缚在原罪之下，另一次将他们从一切罪的束缚中释放出来。\par

% structure: p0009 source=h2 target=subsection reason=relative-heading-rank; base=h2

\subsection{第三章。——佩拉吉乌斯使天主不义。}

我们最终被迫基于神圣权威而同意那些即使以最清晰的理智也无法探究的事。他们自己提醒我们，基督的义德除了信徒外不能使任何人获益，而他们却承认它在一定程度上能使婴儿获益；根据这一点（正如我们已说过的），他们必须毫无回避地在信徒之列中为受洗婴儿找到位置。因此，如果他们不受洗，他们将被列于不信者之中；因此，他们甚至不会得到生命，反而\BibleQuote{天主的愤怒常在他们身上}，因为\BibleQuote{不信从子的，必不得见生命，天主的愤怒常在他身上}\BibleRef{若 3:36} 他们处于审判之下，因为\BibleQuote{不信的，已被定罪}\BibleRef{若 3:18} 他们将被定罪，因为\BibleQuote{信而受洗的，必得救；不信的，必被定罪}\BibleRef{谷 16:16} 现在，让他们考虑，如果一个人没有罪——即，既然他们不可能有本罪，那么他们内也没有原罪——他们凭什么正义可以坚持或努力主张人无份于永生，反而处于天主的愤怒之下，招致天主的审判和定罪？\par

% structure: p0011 source=h2 target=subsection reason=relative-heading-rank; base=h2

\subsection{第四章}

至于\Person{佩拉吉乌斯}让那些反对原罪的人所提出的其他论点，我想，我已经在我长篇论著的前两卷中作了充分而明确的答复。现在，如果我的答复在任何人看来过于简短或晦涩，我请求他们的原谅，并请他们与那些也许指责我的论著不是过于简短而是过于冗长的人达成谅解；而对于那些仍然不理解那些我认为已经尽可能根据主题性质解释清楚的观点的人，我肯定不会因为他们冷漠或不理解而责备或斥责他们。我更希望他们祈求天主赐予他们智慧。\par

% structure: p0013 source=h2 target=subsection reason=relative-heading-rank; base=h2

\subsection{第五章 [III.]——某些人称赞佩拉吉乌斯；佩拉吉乌斯在他的注释中提出的反对原罪的论点。}

但我们确实不可忽略观察的是，这位良善可称赞的人（如认识他的人所描述的那样）并非以自己的名义提出这个反对罪自然传生的论点，而是复述了那些不赞同这一教义的人所主张的，而且不仅限于我刚才引述并反驳的主张，还包括其他我现已进一步着手答复的观点。现在，在说完\Quote{如果（他们说）\Person{亚当}的罪伤害了那些不犯罪的人，那么基督的义德也造福那些不信的人}——这句话，您从我对它的答复中可以看出，不仅与我们所持的立场不相抵触，反而提醒我们应当持守什么——之后，他立刻接着补充说：\Quote{然后他们争辩说，如果圣洗洗去了那旧罪，那么由两个受过洗的父母所生的孩子必然没有这罪，因为他们无法把自己没有的东西传给子女。}他说，\Quote{此外，如果灵魂不是由传生而来，而只是肉体，那么只有后者有罪的传生，只有它应受惩罚；因为他们断言，让那刚刚出生、并非来自\Person{亚当}之团块的灵魂，承担如此古老的、外来的罪，是不公义的。}同样，\Person{佩拉吉乌斯}说，\Quote{绝不能承认那赦免人自己罪过的天主，会将别人的罪归算给他。}\par

% structure: p0015 source=h2 target=subsection reason=relative-heading-rank; base=h2

\subsection{第六章——为什么佩拉吉乌斯不以自己的名义发言。}

请问，难道你没有看到\Person{佩拉吉乌斯}是如何将整个这段文字插入他的著作中的吗？他并不是以自己的名义，而是以别人的名义，他如此清楚这前所未闻的新奇教义（它现在开始发声反对教会古老深入人心的观点），以至于他羞于或害怕自己承认它？也许他自己并不认为一个人生而无罪——因为他承认圣洗对于赦罪是必要的；也不认为一个没有罪的人会被定罪——因为他必须将未受洗的人列在不信者之列，因为福音当然不会欺骗我们，它最清楚地断言：\BibleQuote{不信的，必被定罪}\BibleRef{谷 16:16}；最后，也不认为天主的肖像，当它没有罪时，不能进入天主的国，因为\BibleQuote{人除非由水和圣神而生，不能进天主的国}\BibleRef{若 3:5}——因此它要么没有罪而被抛入永死，要么（更荒谬的是）必须在天国之外拥有永生；因为主在预言最后要对祂的子民说什么时——\BibleQuote{我父所祝福的，你们来承受自创世以来为你们预备的国}\BibleRef{玛 25:34}——也清楚指出了祂所说的国是什么，以这样的话语作结：\BibleQuote{这些人要进入永罚，而义人却要进入永生}\BibleRef{玛 25:46}。那么，这些看法以及其他源于这核心错误的看法，我相信如此可敬的人、如此良好的基督徒绝不会有丝毫接受，因为它们太过悖谬，与基督徒真理相抵触。但是很可能，他仍然被那些否认罪传生的人的论据本身所困扰，以至于渴望听到或知道能回答他们什么；因此，他既不愿沉默地保留那些否认罪传生的人所提出的主张，以便在适当的时候讨论这个问题，同时又不以自己的名义报告这些看法，以免被人认为他自己持有这些看法。\par

% structure: p0017 source=h2 target=subsection reason=relative-heading-rank; base=h2

\subsection{第七章 [IV.]——婴儿原罪的证明。}

现在，虽然我自己或许无力反驳这些人的论点，但我看到紧守圣经最清晰的陈述是多么必要，以便藉助它们解释晦涩的段落，或者，如果心智尚不足以在解释时领悟它们，或当它们隐晦时探究它们，至少也应毫不怀疑地相信它们。但有什么比许多神圣宣告的沉重见证更为清楚呢？这些见证给了我们最清晰的证据：没有与基督的结合，没有人能获得永生和救恩；且没有人能被天主的审判不义地定罪——即与那生命和救恩分离。从这些真理得出的必然结论是：正如婴孩领受圣洗时，除了被纳入教会——换句话说，就是联合于基督的身体和肢体——之外，别无其他效果，除非这一恩惠赐予了他们，否则他们显然处于被定罪的危险中。然而，如果他们确实没有罪，就不能被定罪。既然他们幼小的年龄不可能在自己的一生中犯下罪，那么，即使我们尚不能理解，至少也当相信婴孩继承原罪。\par

% structure: p0019 source=h2 target=subsection reason=relative-heading-rank; base=h2

\subsection{第八章——耶稣甚至也是婴孩的救主。}

因此，如果宗徒说\BibleQuote{藉一人罪恶入了世界，藉罪恶死亡，且如此传给了众人}，\BibleRef{罗 5:12}并且可能被曲解应用于其他含义，那么\BibleQuote{人除非由水和圣神重生，不能进入天主的国}这句话还有歧义吗？\BibleRef{若 3:5}再问，\BibleQuote{你要给祂起名叫耶稣，因为祂要把自己的民族从他们的罪恶中拯救出来}有歧义吗？\BibleRef{玛 1:21}\BibleQuote{健康的人不需要医生，而是有病的人}这句话有任何疑问吗？\BibleRef{玛 9:12}——即耶稣不是为无罪的人所需，而是为那些要从罪中被拯救的人所需。还有，\BibleQuote{人除非吃人子的肉}——即成为祂身体的分享者——\BibleQuote{就不会有生命}这句话有任何歧义吗？通过这些以及我现在略过的类似陈述——在天主的光照下绝对清晰，藉祂的权威绝对确定——真理不是毫不含糊地宣告：未领受圣洗的婴孩不仅不能进入天主的国，而且除非在基督的身体内，就不能获得永生，而正是为了被结合于基督的身体，他们才在圣洗圣事中被洗涤？真理不是毫无怀疑地见证：他们被虔诚的手带到耶稣（即基督，救主和医生）面前，没有其他原因，只是为了藉祂圣事的药物治愈他们罪的瘟疫？那么，我们为什么迟延这样理解宗徒的话语呢？这些话语我们或许过去有些疑惑，如今却要与这些我们毫不怀疑的陈述相一致？\par

% structure: p0021 source=h2 target=subsection reason=relative-heading-rank; base=h2

\subsection{第九章——\BibleQuote{亚当是那要来者的预象}这句话的歧义。}

然而，对于宗徒论及因一人的罪而有多人被定罪，因一人的义而有多人被成义的整段经文，我并无疑惑，除了这句话：\BibleQuote{亚当是那要来者的预象}。因为这句话实际上不仅适用于理解亚当的后裔将同他一样带着罪出生的含义，而且这句话还能被引申出几种不同的意义。因为我们自己也许在不同时期曾从这句话中争辩过各种含义，很可能我们还会提出另一种观点，但这与这里提到的含义并不冲突；甚至伯拉纠也不总是以一种方式解释这段经文。然而，这段含有可疑词语的经文的其他部分，若仔细考察并处理，如同我在本论著第一卷尽力而为的那样，尽管因主题本身必然产生的文风晦涩，但不会不显示出任何其他含义都与普世教会从最初以来所持守的立场不相容——即相信的婴孩藉基督的圣洗获得了原罪的赦免。\par

% structure: p0023 source=h2 target=subsection reason=relative-heading-rank; base=h2

\subsection{第十章[V.]——他表明西彼连从未怀疑婴孩的原罪。}

因此，真福\Person{西彼廉}仔细地指出教会从一开始就将此视为一条清楚明了的信德条款，这并非没有理由。当他主张刚出生的婴孩适合领受基督的圣洗时，有一次有人咨询他是否应在第八天之前施行此圣事，他竭尽全力证明婴孩是完美的，以免有人因日期的数目（因为从前婴孩是在第八天受割损）而认为他们尚缺乏完美。然而，在给予他们充分的论证支持之后，他仍然承认他们并非没有原罪；因为他若否认这一点，就会取消他所主张的婴孩适合领受的圣洗的全部理由。你若愿意，可以亲自阅读这位光荣殉道者的著作\WorkTitle{论幼童圣洗}；因为在迦太基不难找到。但我认为将其中与眼前问题相关的少许陈述转录到我此信中是合适的；并请你们仔细注意。\Quote{现在关于婴孩的情况，}他说，\Quote{你们宣称若在出生后第二、三天带来受洗是不合适的，因为应当尊重旧时的割损礼律法，以致你们认为婴孩不应在出生后第八天之前受洗并成圣——在我们的会议中对这问题形成了截然不同的看法。那里没有一个人赞同你们认为应当做的事；但我们全体反而决定，不应拒绝天主的仁慈和恩宠给予任何由人所生的人。因为主在福音中说：\BibleQuote{人子来不是要毁灭人的生命，而是要救他们，}\BibleRef{路 9:56}因此，就我们所能做的，不容任何灵魂失落，如果可以的话。}你们看到，在这些话中，他认为一个人若没有那救恩的圣事而离开此生，不仅是肉身，连灵魂也陷入毁灭和死亡。因此，即使他没有再说别的，我们也有理由从他的话语中得出结论：没有罪，灵魂就不能灭亡。然而，请看他在稍后维护婴孩的无辜时，同时以最明白的措辞承认了什么：\Quote{但是，如果有什么能阻碍人获得恩宠，那么更重的罪反而更应阻碍那些已达到成年、高龄和老年的人。然而，既然罪的赦免甚至赐给那些信后最严重的罪人，无论他们先前如何得罪天主，而且没有人被禁止领受圣洗和恩宠，那么一个新生婴孩岂不更不应被禁止？他除了因按肉体出生自亚当，从出生时就已感染了原始死亡的传染之外，并没有犯过任何罪。此外，这一事实本身岂不也更有助于他们接受罪的赦免，因为他们所得的赦免不是自己的罪，而是别人的罪！}\par

% structure: p0025 source=h2 target=subsection reason=relative-heading-rank; base=h2

\subsection{第十一章——古人假定原罪。}

你们看，这位伟人依据古代无疑的信德规范，以何等的信心表达自己。在提出这些非常明确的陈述时，他的目的是借助这些牢固的结论来证明由他信友所提交的、并告知他会议已通过的谕令中那个不确定之点，即：如果一个婴孩甚至在出生后第八天之前被带来，没有人应犹豫为他付洗。当时，会议并未像对待新事物那样确定或确认婴孩受原罪束缚，也没有受到任何人的反对；而是在进行另一场辩论时，讨论了关于肉身割损律法的问题：是否应在第八天之前为他们付洗。没有人赞同那个否认的人；因为这不是一个可以讨论的开放问题，而是一个固定不可动摇的真理：如果灵魂在未获得圣洗圣事的情况下结束此生，它将丧失永远救恩；但同时，刚出母腹的婴孩被认为只受原罪罪债的影响。因此，虽然在他们的情况下罪的赦免更容易，因为这些罪来自别人，但赦免仍然是必不可少的。正是在这些确凿的基础上，解决了第八天的悬疑问题，会议决定：人出生后，不应失去一天的时间来提供帮助，以防止他永远丧亡。当时，对肉身割损律法作为未来之事预像的理由的解释，其目的不是要我们理解圣洗应在出生后第八天施行，而是我们要在基督的复活中受精神割损，基督在受难后第三天复活，但在按时间计算的一周的日子中，是在第八天，即安息日后的第一天。\par

% structure: p0027 source=h2 target=subsection reason=relative-heading-rank; base=h2

\subsection{第十二章 [VI.]——有关原罪的普遍共识。}

如今，那些人又以新争论中一种奇怪的胆量，企图使我们对一个我们的先辈在处理某些人认为不确定的问题时，总是作为最确定的前提提出来的要点感到不确定。我现在无法确切地说出这场争论最初是什么时候开始的；但我知道一件事，就连我们当代因在教会文献方面的巨大勤奋和学识而闻名的\Person{圣热罗尼莫}，在他著作中处理各种问题时，也使用同样最确定的假设，而不加论证。例如，在他关于\Person{约纳}先知的注释中，当谈到那个提到婴儿因禁食而受处罚的段落时，他说：\Quote{年龄最大的先来，然后所有其余的人一直向下传递到最小的。\BibleRef{纳 3:3} 因为没有人没有罪，无论他的年龄只有一天，还是他活了多少岁。\BibleRef{约 25:4} 因为如果连星星在天主眼前也是不洁的，那么虫子和腐朽之物，那些受违抗天主命令的\Person{亚当}所犯之罪奴役的，又该如何呢？} 如果我们能轻易询问这位最有学问的人，他会向我们提起多少用两种语言讨论过神圣圣经并写过基督教争论的作者呢？他们自从\Person{基督}的教会建立以来，从未持有过任何别的观点——他们既没有从他们的前辈那里接受任何别的观点，也没有传给他们的后代别的观点？就我个人的阅读来说，虽然非常有限，但我并不记得曾从接受两约的基督徒那里听到过任何有关这个问题的别的观点，无论他们是在天主教会内，还是在任何异端或分裂团体中。我不记得，我说，我曾在任何对这种文献有所贡献的作者中，无论他们是遵循了正典圣经，还是自以为遵循了，或者希望别人这样以为，发现任何别的观点。我不知道这个问题是从哪里突然冒到我们面前的。不久前，在\Place{迦太基}与某些人一次偶然的交谈中，我的耳朵突然被这样的主张冒犯了：\Quote{婴儿领洗不是为了获得赦罪，而是为了在\Person{基督}内被圣化。} 尽管我对此新奇的见解感到非常不安，但由于没有机会反驳，而且提出者也不是让我担心的有影响力的人，我很容易就让它被忽略和遗忘了。看啊！如今它却带着炽热的热情被用于反对教会；看啊！它被写作下来永久地提交给我们注意；不，事情甚至达到了如此令人分心的地步，以致我们的弟兄也来咨询我们；我们实际上不得不在辩论和写作中反对它的进展。\par

% structure: p0029 source=h2 target=subsection reason=relative-heading-rank; base=h2

\subsection{第十三章 [VII.] — \Person{约维尼安}的错误并没有延伸到那么远。}

几年前，在\Place{罗马}住着一个叫\Person{约维尼安}的人，据说他说服了甚至年老的修女结婚——当然，不是通过勾引，好像他想让其中任何一位成为他的妻子，而是通过争辩说：奉献自己于苦修生活的贞女在天主面前并不比有信德的妻子更有功德。然而，他从未想到要在这个奇思怪想之外，还敢断言人的孩子生来没有原罪。如果他确实添加了这样的意见，那些妇女可能会更容易同意结婚，以生下如此纯洁的后代。当这个人的著作（因为他竟敢写作）由弟兄们寄给\Person{热罗尼莫}以反驳时，他不仅没有在其中发现这样的错误，而且，在寻找他的奇思怪想加以反驳时，他在其他段落中发现了这个非常明确的关于人原罪的教义的见证，\Person{热罗尼莫}确实由此确信那个人相信那项教义。以下是他在讨论这事时所说的话：\Quote{那说自己住在\Person{基督}内的，就应当照祂所行的去行。\BibleRef{若 2:6} 我们给我们的对手选择他喜欢的选项。他住在\Person{基督}内吗？还是没有？如果他住在内，那么，让他像\Person{基督}一样行走。然而，如果承担效法\Person{基督}的美德是一件轻率的事，那么他就不住在\Person{基督}内，因为他并不像\Person{基督}所行的去行。祂没有犯罪，口中也没有诡诈；\BibleRef{依 53:9} 祂被骂不还口；如同羊在剪毛者前缄默，祂也不开口；\BibleRef{依 53:7} 这世界的首领来到祂这里，在祂内一无所有；\BibleRef{若 14:30} 祂虽然没有犯罪，天主却使祂成为罪。\BibleRef{格后 5:21} 然而我们，按照雅各伯书信，都犯许多过失；\BibleRef{雅 3:2} 我们中没有一个人是纯洁无玷的，即使他的生命只有一天。\BibleRef{约 14:5} 因为谁能夸口说自己有洁净的心呢？谁能自信自己纯洁无罪呢？我们都按\Person{亚当}过犯的样式被定罪。为此，\Person{达味}也说：\InnerQuote{看，我是在罪过中生成的，我的母亲在罪恶中怀了我。}}\par

% structure: p0031 source=h2 target=subsection reason=relative-heading-rank; base=h2

\subsection{第十四章——然而，所有争论者的意见并不是正统权威；原罪为什么是别人的罪；我们在\Person{亚当}内曾是一个人。}

我引用这些话，并不是说我们可以倚仗每个辩论者的意见如同倚仗正典权威一样；我这样做，是为了让人看到，从起初直到这个产生了这种新奇意见的时代，原罪教义如何作为教会信仰的一部分被极其坚定地守护着，以至于它通常被当作最确凿的根据来驳斥其他错误意见，而非本身被任何人当作错误而遭受驳斥。此外，在圣经正典的神圣书籍中，这一教义的权威以最清楚、最充分的方式被有力地宣告。宗徒大声疾呼：\Quote{藉着一个人，罪恶进入了世界，藉着罪恶，死亡；这样死亡就传到了众人，因为众人都犯了罪。}\BibleRef{罗 5:12} 如今从这些话中，绝不能断言亚当的罪恶甚至伤害了那些\Emph{不犯罪的人}，因为圣经说：\Quote{\Emph{众人都犯了罪。}}诚然，婴儿的那些罪也并非被称为\Emph{别人的}，仿佛它们完全不属于婴儿一样，因为众人在亚当内都犯了罪，那时在他们的本性中，凭借那使他能产生他们的天生力量，他们全体仍是一个亚当；但它们被称为\Emph{别人的}，是因为他们那时还没有活自己的生命，而那个人的生命包含了日后后裔中的一切。\par

% structure: p0033 source=h2 target=subsection reason=relative-heading-rank; base=h2

\subsection{第十五章 [VIII.]——我们都犯了亚当的罪。}

\Quote{确实，}他们说，\Quote{天主绝不会赦免一个人的罪，却将别人的罪归咎于他。}祂确实赦免，但只赦免那些由圣神重生的人，而非由血肉所生的人；但祂不再将别人的罪，而只将人自己的罪归咎于他。那些罪无疑是别人的，当那些承受它们的人被生育时还没有存在；但现在通过血肉的生育，这些罪属于他们，而他们还没有通过属灵的重生而被赦免。\par

% structure: p0035 source=h2 target=subsection reason=relative-heading-rank; base=h2

\subsection{第十六章。——错误的起源；从割礼的包皮和小麦的糠秕中寻求的比喻。}

\Quote{但是，}他们说，\Quote{如果圣洗圣事洁净了原罪，那么由两位已领洗的父母所生的人应当没有这个罪；因为他们不能将自己没有的东西传给他们的儿女。}现在请看错误通常如何滋生：就是当人能够提出他们不能理解的话题时。因为在什么听众面前，用什么言语，我能解释为何有罪的、有死的开端对那些已经开启了另一种不朽开端的人不构成障碍，同时却对那些由这些本身不构成障碍的人从同一个有罪开端所生的人构成障碍呢？一个心智受累于其成见和自己顽固的锁链的人，怎能理解这些事呢？如果我确实在反对那些完全禁止婴儿领洗、或者主张给他们领洗是多余的人，认为他们因生于信徒父母而必然享有父母功德的人面前诉讼，那么我就有责任唤醒自己或许以更大的劳苦和努力来驳斥他们的意见。那样，如果我在迟钝好辩的人面前，由于主题本质的晦涩而在驳斥谬误、灌输真理中遇到困难，我或许会诉诸那些明显易得的比喻；我反过来会问他们一些问题：例如，如果他们困惑于知道罪在通过圣洗圣事被洁净后如何仍留在已领洗的父母所生的人身上，他们如何解释割礼除去包皮后包皮如何仍留在受割礼者的儿子身上？或者，又如，被人工仔细簸扬的糠秕如何仍存在于从被簸净的麦粒所生的谷粒中？\par

% structure: p0037 source=h2 target=subsection reason=relative-heading-rank; base=h2

\subsection{第十七章 [IX.]——基督徒并不总是生养基督徒，纯洁者并不总是生养纯洁的子女。}

用这些和类似明显的论据，我尽力说服那些认为洁净的圣事对洁净之人的子女是多余的人，使他们明白，为已受洗父母的婴儿施洗是多么正确的判断；以及，在一个人内有两重种子——肉身的死亡种子和灵魂的不朽种子——的情况下，他因圣神而重生，而这对他因其肉身而生的儿子却可能成为障碍，这如何可能发生；并且，在一个人中藉着赦免得以洁净的，在另一个人中却同样需要藉着赦免来洁净，正如所假设的割礼和簸扬与打谷的例子。但现在，当我们与那些同意应给已受洗者的子女施洗的人辩论时，我们可以更方便地进行讨论，并且可以说：你们这些声称从罪恶玷污中得洁净之人的子女按理应当生而无罪的人，为什么不明白，按照同样的规则，你们同样可以说基督徒父母的子女按理应当生而为基督徒？那么，你们为什么反而主张他们应当成为基督徒？难道在他们父母身上——就是那些听到\BibleQuote{你们不知道你们的身体是基督的肢体吗？}\BibleRef{格前 6:15}的人——没有一个基督徒的身体吗？也许你们认为，一个基督徒的身体可以从基督徒父母生出，而没有领受基督徒的灵魂？那么，这就更加奇妙了。因为你们会随心所欲地认为灵魂有两种可能之一——当然，你们与宗徒一致，认为出生前灵魂没有行善或作恶：\BibleRef{罗 9:11} ——或者灵魂是藉着传衍而来的，正如基督徒的身体是基督徒，他们的灵魂也应当是基督徒；或者灵魂是由基督创造的，要么在基督徒身体内，要么为了基督徒身体的缘故，因此它应当在基督徒的状态中被创造或赋予。除非你们或许会声称，虽然基督徒父母有能力生出基督徒的身体，但基督自己却未能产生一个基督徒的灵魂。那么，请相信真理，并且看到，正如（你们自己也承认）一个非基督徒可以从基督徒父母生出，一个不属基督的肢体可以从基督的肢体生出，并且（为回答那些尽管虚假却仍以某种方式拥有宗教感的人）一个未祝圣的人可以从祝圣的父母生出；同样，一个未洁净的人也完全可以从洁净的父母生出。现在，你们将给我们什么解释，为什么从基督徒父母生出非基督徒，除非不是因为出生，而是因为重生使人成为基督徒？因此，用同样的理由解答你们自己的问题，即：罪的洁净不是由于出生而来，而是藉着重生而来。因此，任何从已洁净（因为已重生）的父母生下的孩子，他自己也必须重生，以便他也能被洁净。因为父母完全可能传给子女他们自己所没有的东西——这不仅像麦子产生糠秕、受割礼者产生包皮，而且也像你们自己所引用的例子，即信徒将不信传给后代；然而，这不归因于因圣神而重生的信者，而是由于他们从肉身而生的可朽坏种子的缺陷。因为关于你们认为有必要通过信者的圣事使他们成为信者的婴儿，你们不否认他们是生于不信之中，尽管父母是信者。\par

% structure: p0039 source=h2 target=subsection reason=relative-heading-rank; base=h2

\subsection{第十八章 [X.]——灵魂是由自然传衍而来的吗？}

嗯，但是\Quote{如果灵魂不是被传衍的，而只是肉身被传衍，那么只有肉身有罪的传衍，也只有肉身应受惩罚：}这是他们所想的，说\Quote{那新近产生、且不是出于亚当实质的灵魂，却要承担另一个许久以前所犯的罪，这是不义的。}现在，我请你们注意，谨慎的\Person{白拉奇}是如何感到关于灵魂的问题非常困难，并且相应地行事——因为我刚才引用的这些话是从他的书中抄来的。他并非绝对地说\Quote{因为灵魂不是被传衍的，}而是假设地说\Emph{如果灵魂不是被传衍的}，在如此模糊的主题上（对此我们在圣经中找不到确定而明显的证据，即使能找到也非常困难）正确地决定以犹豫而非自信来说话。因此，我也同样，不轻率地断言来回答这个命题：如果灵魂不是被传衍的，那么正义何在，那新近被创造、完全不受罪的感染的东西，却要在婴儿身上被迫忍受肉体的情欲和其他折磨，而且更可怕的是，甚至要受恶灵的攻击？因为肉体从未如此受苦，而活着的、有感觉的灵魂却不受惩罚。如果这确实被证明是正义的，那么同样可以证明，原罪以何等正义存在于我们有罪的肉身中，之后藉着圣洗圣事和天主的仁慈恩宠得以洗净。如果前一点无法证明，我想后一点同样无法证明。因此，我们要么沉默地忍受这两个立场，记住我们是人；要么在另一时间，如有必要，准备另一部关于灵魂的著作，以谨慎和清醒的态度讨论整个问题。\par

% structure: p0041 source=h2 target=subsection reason=relative-heading-rank; base=h2

\subsection{第十九章 [XI.]——在亚当内的罪与死亡，在基督内的义与生命。}

宗徒说：\BibleQuote{藉着一个人，罪进入了世界，藉着罪，死亡也进入了世界；这样死亡就传遍了众人，因为众人都犯了罪。}\BibleRef{罗 5:12}然而，目前我们必须如此接受，以免轻率而愚昧地反对圣经中许多伟大的篇章，它们教导我们：没有一个人能获得永生，除非与基督有那在祂内、并与祂一起实现的合一，即当我们领受了祂的圣事，并被纳入祂身体的肢体时。如今，宗徒致罗马人的这句话：\BibleQuote{藉着一个人，罪进入了世界，藉着罪，死亡也进入了世界；这样死亡就传遍了众人，因为众人都犯了罪，}与他对格林多人所说的意思吻合：\BibleQuote{因为死亡是因一个人而来，死者的复活也是因一个人而来。就如在亚当内众人都死了，照样在基督内众人都要复活。}\BibleRef{格前 15:21}因为无人怀疑这里所指的是身体的死亡，因为宗徒极其热切地谈论身体的复活；他在这里似乎对罪保持沉默，原因是：所讨论的并不是正义的问题。这两点都在\WorkTitle{罗马书}中提到，并且宗徒以极长的篇幅坚持这两点——在亚当内的罪，在基督内的正义；在亚当内的死亡，在基督内的生命。然而，如我早已指出的，我已在本论集的第一卷中，尽我所能、并按其所需，详尽地查考并开释了宗徒论证的全部这些言词。\par

% structure: p0043 source=h2 target=subsection reason=relative-heading-rank; base=h2

\subsection{第二十章——死亡的刺是什么？}

但是，即使在致格林多人书那段（他充分谈论了复活）中，宗徒也以下列方式结束他的论述，不容我们怀疑身体的死亡是罪的结果。因为他先说：\BibleQuote{这可朽坏的必须穿上不可朽坏的，这可死的必须穿上不可死的。当这可朽坏的穿上了不可朽坏的，这可死的穿上了不可死的，那时}接着又说：\BibleQuote{就要应验那经上所记载的话：死亡被吞灭在胜利中。死亡，你的胜利在哪里？死亡，你的刺在哪里？}最后他附上了这些话：\BibleQuote{死亡的刺就是罪，而罪的权势就是法律。}\BibleRef{格前 15:56}现在，由于（正如宗徒的话最清楚地宣告）死亡要在那可朽坏和可死的穿上不可朽坏和不可死时被吞灭在胜利中——也就是说，当天主借着那住在我们内的圣神，连我们必死的身体也要复活起来时——显然，这死亡之身的刺，即身体复活的相反，就是罪。然而，刺是那造成死亡的，而不是死亡造成的，因为我们因罪而死，而不是因死亡而犯罪。因此，它被称为\Quote{死亡的刺}，正如形成\Quote{生命树}这一表达的原则一样——不是人的生命产生它，而是藉着它人的生命被造成。同样，\Quote{知善恶树}是那造成人的知识的东西，而不是人藉着知识造成的。同样，\Quote{死亡的刺}是那产生死亡的，而不是死亡产生的。我们也类似地使用\Quote{死亡的杯}这一表达，因为有人藉着它已经死亡或可能死亡——当然不是指一个垂死或死人制成的杯。死亡的刺因此就是罪，因为藉着罪的刺，人类已被杀害。何必再问：是那一种死——灵魂的死还是身体的死？是第一个我们如今都在经历的死亡，还是第二个恶人将来要受的死亡？不必如此好奇地追问这个问题，没有狡辩的余地。宗徒表达这一情况的言词回答了这些问题：\BibleQuote{当这可死的，他说，穿上了不可死的，那时就要应验那经上所记载的话：死亡被吞灭在胜利中。死亡，你的胜利在哪里？死亡，你的刺在哪里？死亡的刺就是罪，而罪的权势就是法律。}他是在谈论身体的复活，那时死亡将被吞灭在胜利中，当这可死的穿上了不可死。那时，在身体的复活中，死亡被吞灭在胜利中，要向死亡发出胜利的欢呼；那时要对它说：\BibleQuote{死亡，你的胜利在哪里？死亡，你的刺在哪里？}因此，这话是对身体的死亡说的。因为胜利的不朽将吞灭它，当这可死的穿上了不朽。我重复，这是对身体的死亡说的：\BibleQuote{你的胜利在哪里？}——你藉以征服一切的胜利，连天主子都与你交战，祂没有退缩而是与你角力，并胜过了你。在那些死去的人中，你征服了；但在这些复活的人中，你自身被征服了。你的胜利只是暂时的，你在其中吞没了他们的身体。我们的胜利将是永恒的，你在其中被吞没在那些复活的身体中。\BibleQuote{你的刺在哪里？}——即那使我们被刺伤而中毒的罪，你藉此将自身固定在我们的身体中，并长久地占有它们。\BibleQuote{死亡的刺就是罪，而罪的权势就是法律。}我们都在一人内犯了罪，所以我们都因一人而死；我们领受了法律，不是要按它的诫命改正而终止罪，而是因违犯而增加罪。因为\BibleQuote{法律闯进来，为的是使罪增多；}\BibleRef{罗 5:20}并且\BibleQuote{圣经把众人都禁锢在罪中；}\BibleRef{迦 3:22}但是\BibleQuote{感谢天主赐给了我们因我们的主耶稣基督而获得的胜利；}\BibleRef{格前 15:57}为的是\BibleQuote{在那里罪增多的地方，恩宠更加增多；}\BibleRef{罗 5:20}并且\BibleQuote{使那因信耶稣基督的恩许，能赐予信者们；}\BibleRef{迦 3:22}好使我们能藉着不朽的复活战胜死亡，并藉着白白成义战胜死亡之刺，即罪。\par

% structure: p0045 source=h2 target=subsection reason=relative-heading-rank; base=h2

\subsection{第二十一章——关于触碰经期妇女的诫命不可按比喻理解；圣事的必要性。}

因此，在这件事上，谁也不可受欺骗或欺骗人。我们所考察的圣经的明显含义，消除了所有的模糊之处。正如死亡在我们这必死的身体中从起初而来，同样，从起初罪就被带入了我们这有罪的血肉之中。为了医治这罪（它既通过传衍而来，又因故意的过犯而加剧），也为了使我们的血肉本身复生，我们的医生以罪孽血肉的形像而来，祂不需要健康的人，只需要患病的人——祂来不是召义人，而是召罪人\BibleRef{谷 2:17}。因此，当宗徒劝告信徒不要与不信的伴侣分离时所说的话：\BibleQuote{因为不信的丈夫因妻子而成了圣洁的，不信的妻子因丈夫而成了圣洁的；不然，你们的儿女就是不洁的，但如今他们是圣洁的}\BibleRef{格前 7:14}，必须要么像我们自己在别处以及白拉奇在他对同一封致格林多人书信的注释中所解释的那样，按照已提到的段落的含义来理解——即有时妻子赢得了丈夫归向基督，有时丈夫使妻子皈依，而即使是父母一方基督徒的意愿也足以使他们的儿女成为基督徒；要么（正如宗徒的话似乎更倾向于指示，并在一定程度上迫使我们这样理解）这里要理解某种特别的圣化，借此不信的丈夫或妻子因信主的伴侣而成为圣洁，并且信主父母的儿女也借此成为圣洁——不论是因为丈夫或妻子在妇女经期时避免同居，因为他们在法律中学会了这一义务（因为厄则克耳将此列为不可按比喻理解的诫命之一\BibleRef{则 18:6}），还是由于某种那里没有明确规定的自愿圣化——一种因婚姻生活和子女的紧密联系而产生的圣洁沾溉。然而，无论所指的是哪一种圣化，必须坚定持守这一点：使人成为基督徒并赦免罪过的有效方法，只有通过人按照基督和教会的制度领受圣事而成为信徒。因为不信的丈夫和妻子，尽管与圣洁公义的配偶有亲密的结合，也不能洗刷那使人隔离天主的国并被判入永罚的罪；同样，无论父母如何公义圣洁，所生的子女也不能免除原罪的罪债，除非他们受洗归于基督；我们为这些婴儿应更恳切地祈求，因为他们自己越不能提出请求。\par

% structure: p0047 source=h2 target=subsection reason=relative-heading-rank; base=h2

\subsection{第二十二章——我们应当热切确保婴儿领受圣洗圣事。}

这正是争议所指向的目标——我们必须借助古老真理与之斗争的新奇主张：即婴儿受洗完全是多余的。这不是因为这种意见被明确宣告，免得教会如此牢固建立的习俗无法承受其攻击者。但是，如果我们被教导要帮助孤儿，那么对于这些儿童，尽管他们在父母的保护下，但如果基督的恩宠被拒绝赐予他们（他们自己完全无法要求这恩宠），他们将比孤儿更无依无靠、更为可怜，我们岂不更应为他们的缘故而努力？\par

% structure: p0049 source=h2 target=subsection reason=relative-heading-rank; base=h2

\subsection{第二十三章。——结语}

至于他们说，有些人使用理性，在这世上已经生活并且正在生活而没有罪，我们希望这是真的，我们应当努力使之成真，我们应当祈求它成真；但同时，我们应当承认它尚未成真。因为对于那些渴望、努力并相称地为此结果祈祷的人，他们身上所存的任何罪过都每天被赦免，因为我们真诚地祈祷：\BibleQuote{宽免我们的罪债，如同我们宽免得罪我们的人}\BibleRef{玛 6:12}。凡否认这祈祷在此生对于每个认识并承行天主旨意的义人是必要的（除了一位圣中之圣），就大大地错了，并且全然无法使祂喜悦——他赞美祂，却与自己相悖。此外，如果他自认为是那样的人，\BibleQuote{他欺骗自己，真理就不在他内}\BibleRef{若一 1:8}——没有别的原因，只是因为他思想虚假的事。因此，那位不需要健康者、而需要患病者的医生，知道如何医治我们，并通过医治使我们达致永生；祂在现世并没有从那些祂赦免其罪过的人身上除去因罪而施加的死亡，为使他们进入战斗，以最真诚的信德克服对死亡的恐惧。在一些情况下，祂甚至拒绝帮助祂的义仆（只要他们还能达到更高的境界以获致完美的义德），为使我们（因为在祂眼中没有一个活人是义人）永远感到有义务感谢祂对我们的慈悲包容，从而藉圣洁的谦卑，医治我们一切过犯的最初原因——即骄傲的膨胀。这封信，按我最初的意图，原是一封短信；它已增长为一本长书。但愿它已变得完美，正如它终于完成了一样！\par

\SourceRecord{Translated by Peter Holmes and Robert Ernest Wallis, and revised by Benjamin B. Warfield.；Nicene and Post-Nicene Fathers, First Series；Vol. 5.；Edited by Philip Schaff.；Buffalo, NY: Christian Literature Publishing Co.,；1887.；<http://www.newadvent.org/fathers/15013.htm>.}
